\documentclass[11pt,a4paper]{article}
\usepackage[T1]{fontenc}
\usepackage[utf8]{inputenc}
\usepackage{lmodern}
\usepackage[margin=24mm,headheight=15pt]{geometry}
\usepackage{amsmath,amssymb,amsthm,mathtools,mathrsfs}
\usepackage{booktabs,longtable,array,enumitem}
\usepackage{graphicx,microtype,xurl,xcolor}
\usepackage[colorlinks=true,linkcolor=blue!55!black,citecolor=blue!55!black,urlcolor=blue!55!black]{hyperref}
\usepackage{fancyhdr}
\pagestyle{fancy}
\fancyhf{}
\lhead{\small SIGN\_NEXT}
\rhead{\small Hodge index, integral Fock models, and heights}
\cfoot{\thepage}
\setlength{\emergencystretch}{2em}
\setlength{\parskip}{3pt}
\allowdisplaybreaks
\newtheorem{theorem}{Theorem}[section]
\newtheorem{proposition}[theorem]{Proposition}
\newtheorem{corollary}[theorem]{Corollary}
\newtheorem{lemma}[theorem]{Lemma}
\theoremstyle{definition}
\newtheorem{definition}[theorem]{Definition}
\theoremstyle{remark}
\newtheorem{remark}[theorem]{Remark}
\DeclareMathOperator{\CH}{CH}
\DeclareMathOperator{\Spec}{Spec}
\DeclareMathOperator{\Pic}{Pic}
\DeclareMathOperator{\Tr}{Tr}
\title{\textbf{SIGN\_NEXT}\\[4pt]
Grothendieck's Hodge-index route after sign-only purity}
\author{Research and proof audit: GPT-6 Astra Pro\\
\small Run SIGN\_NEXT, with three parallel research and review agents}
\date{21 November 2026}
\begin{document}
\maketitle
\begin{abstract}
We separate the Lefschetz correspondence condition from its critical
cycle-surjectivity criterion and from the index inequality. A literal
lower-critical ordinary-Gysin reading of cote 16 scan 050 fails already
on the square of a projective line. An explicitly weighted Gysin
condition gives an exact corrected criterion, and the missing
inverse on the actual hyperplane is identified. The supplied CM Hodge
theorem, Milne's theorem, and the classical abelian Lefschetz theorem
complete the $C+I$ route for abelian varieties. A rigid tensor
countermodel shows why pure eigenvalues and coefficient traces alone
do not supply the positive trace form used for semisimplicity.

On the integral side we construct Fock/exterior and Weyl-lattice
intertwiners, prove saturation of the divided-power order in every
rank, give bounded index matrices in every rank, and compute the
ordinary orders through rank six. The finite
theta coefficient kernels are treated with their actual lattices,
duals, moments and determinant factors, including an inhabited
two-step algebraic theta-jet family and its exact size limitations.
Finally we give the vertical
Gram correction to cote 44 Theorem 5 and compute all three nodal
contributions in the $\mathbf P^2\times\mathbf P^2$ example. Every
assertion is assigned a stable identifier; proof checks and the
boundaries of the comparisons are recorded separately.
\end{abstract}
\tableofcontents
\clearpage
\section*{Inputs, notation, and check conventions}
\addcontentsline{toc}{section}{Inputs, notation, and check conventions}

\paragraph{Settled input, not reproved.}
HODGE\_SIGN SG-08--SG-11, with GCT\_NSRH F03, is the stipulated
sign-only purity theorem. We use the predecessor and its supplied
independent check, \path{CHECK_a3ccbc11_20261121.md}; the latter
uses Katz's Appendix A.3--A.5 for its generic weak Lefschetz input.
Nothing in this paper reopens or reproves that theorem. It gives the
purity conclusions of Deligne, \emph{Weil I}, Theorem 1.6, and
\emph{Weil II}, Theorem 3.3.1 over a point, in the precise scope
specified by the supplied run.

The second stipulated input is OpenAI, supplied 032,
Theorem 1.1 (statement p.~1; concluding proof p.~47),
equivalently van Hoften--Sempliner, Hexagon 2610.00259v2,
Theorem I, p.~1: the rational Hodge conjecture for complex CM
abelian varieties. Its use here is a cited theorem, as the task
requires. The exact downstream implication is supplied 032,
Corollary 8.4, p.~50, through Milne 2002, Theorem 3.3(a),(b).

Alp\"oge's \emph{Remark}, with its supplied Lean archive, is also
treated as cited input, including the RH and Dirichlet-GRH statements
in \path{ALPOGE_Main.lean} and \path{ALPOGE_Statements.lean}.
Those global RH statements are not needed to prove the new geometric
or integral identities below. The Fock comparison uses the actual
definitions and closed declarations located in the paper and source
manifests; no new build of the entire Lean archive is claimed.
No Masullo statement is an input.

\paragraph{Historical notation.}
$C_{\mathrm G}(X)$ means the algebraicity of the normalized
Lefschetz lowering operator in cote 16, hence modern $B(X)$.
$A(X,L)$ means Lefschetz surjectivity on rational algebraic cycle
classes. $I(X,L)$ is the Hodge standard index inequality.
Neither $C_{\mathrm G}$ nor $A$ is silently identified with the
modern K\"unneth condition $C$. In the integral sections $D$
denotes ordinary contraction, while elsewhere a statement
$D(X)$, when explicitly used, denotes equality of numerical and
homological equivalence. Ordinary contraction is not normalized
inverse-string lowering.

\paragraph{Archive locators.}
An archive \emph{scan} is the supplied one-based PDF page, not the
manuscript folio. Cote 16 scans 044--047 form the letter dated
4.1.1967; scans 048--050 form the separate letter dated 6.1.1966
as written. Cote 44 scan 004 contains Theorem 5 on manuscript
page 2. The overtyped glyph and mixed degree conventions on
cote 16 scan 050 are treated as textual ambiguities, separately
from the mathematical counterexample to a specified reading.

\paragraph{Classical foundational results.}
We use hard Lefschetz for the chosen Weil realization; for
$\ell$-adic cohomology the exact general locator is Deligne,
\emph{Weil II}, Theorem 4.1.1. Weak Lefschetz and its Gysin form
are SGA 5, VII, Theorem 7.1 and Corollary 7.2. These known results,
the cited classical standard-conjecture implications, and the
height comparison inputs are not presented as new findings.

\paragraph{What a check means.}
A new proof is placed in Part I only after a complete reading,
with a recorded outcome. Independent reconstruction, targeted
whole-proof reading, source checking, and computation are separate
event types in \path{CHECK_EVENTS.tsv}. The integer computations
are reproducible exact-arithmetic checks, not Lean formal
certificates. All agents used the inherited GPT-6 Astra Pro model;
the separate effort setting is not exposed and is not invented.
Earlier drafts and corrections remain identifiable in the working
review notes. Part II contains no conjecture disguised as a proved
finding.

\part{Checked findings and their proofs}



% SIGN_NEXT main-critical contribution. Author: GPT-6 Astra Pro, child /root/critical_lefschetz.
% Reference keys to integrate: GrothCoates1966, GrothCoates1967, Kleiman1968, SGA5.
\section{The critical-degree proposition: a corrected formulation}
\label{sec:critical-proposition}

Throughout this section, $H$ is a characteristic-zero Weil cohomology
with hard Lefschetz and the usual weak Lefschetz theorem for the smooth
hyperplane sections in use.  Tate twists are suppressed in formulas; an
algebraic class of degree $2r$ belongs to $H^{2r}(T)(r)$.
Write
\[
 \operatorname{Alg}^r(T)=\operatorname{im}\bigl(
 \mathrm{CH}^r(T)_{\mathbf Q}\longrightarrow H^{2r}(T)(r)\bigr).
\]
Algebraicity always refers to rational cycle classes.  We use
$C_{\mathrm G}(X)$ for the condition called $C(X)$ in the Coates
letters: the normalized Lefschetz lowering correspondence is algebraic.
It is the condition called $B(X)$ by Kleiman in 1968 and in modern
standard-conjecture notation.  It is not the modern K\"unneth condition.
The cycle condition $A(X,L)$ says that the hard-Lefschetz isomorphisms
restrict to isomorphisms on the spaces $\operatorname{Alg}^r(X)$.

\paragraph{SN-C01 (the source and the left-hand label).}
The left side of the proposition relevant to the task is
$C_{\mathrm G}(X)$, not the weaker cycle condition $A(X,L)$.
This is the mathematical reconstruction of the ambiguous overtyped
letter in cote~16, scan~050 (folio~49).  It is corroborated by the
CSG transcription of the letter dated 6.1.1966, printed p.~3,
\cite{GrothCoates1966}, and by the preceding archived letter dated 4.1.1967,
printed p.~2, Proposition (i), (ii bis), \cite{GrothCoates1967}.
The exact published result to use is Kleiman's Theorem~2.13(i),
(ii$'$), pp.~372--373 \cite{Kleiman1968}: if $Y\subset X$ is a
smooth hyperplane section and $\dim X=n$, then
\begin{equation}
 C_{\mathrm G}(X)\quad\Longleftrightarrow\quad
 C_{\mathrm G}(Y)\ \text{and}\
 L_{X\times X}^{\,2}:
 \operatorname{Alg}^{n-1}(X\times X)
 \xrightarrow{\ \sim\ }
 \operatorname{Alg}^{n+1}(X\times X),
 \label{eq:kleiman-critical}
\end{equation}
where the product is polarized by the Segre class
$L_{X\times X}=L_X\otimes1+1\otimes L_X$.
This known theorem is cited, not reproved here.  In particular,
a statement about $A(X,L)$ alone must not be used in its place.
The same source's condition (iv) gives the criterion along a flag
of smooth hyperplane sections.

There is a second issue in scan~050, distinct from the left-hand letter.
Its inequalities use $2i$ while several cohomological superscripts
are printed $i$; the critical-degree prescription says $2i=d-2$
when $d$ is even, whereas the preceding ordinary-Gysin variant uses
the upper half of the range.  One cannot silently combine these
pieces into a valid ordinary-Gysin condition.  The following
counterexample requires $i=0$, so it is unaffected by the missing
factor of two in the superscripts.

\begin{proposition}[SN-C02: the literal lower critical Gysin condition fails]
\label{prop:literal-critical-fails}
Read the even critical ordinary-Gysin condition as
\begin{equation}
 j_*\operatorname{Alg}^{m-1}(W)=\operatorname{Alg}^{m}(T),
 \qquad \dim T=2m,
 \label{eq:bad-low-gysin}
\end{equation}
for a smooth hyperplane section $j:W\hookrightarrow T$.
It is not a consequence of $C_{\mathrm G}(X)$ when $T=X\times X$,
and therefore cannot occur on the right side of an equivalence with
that condition.  This failure occurs for $X=\mathbf P^1$ and also
for an elliptic curve.
\end{proposition}
\begin{proof}
Take $X=\mathbf P^1$, so that $T=\mathbf P^1\times\mathbf P^1$,
$m=1$, and let $a,b$ be the two ruling classes.  Take $W$ to be a
smooth divisor of bidegree $(1,1)$.  It is isomorphic to
$\mathbf P^1$, and
\[
 H^*(T,\mathbf Q)=\mathbf Q[a,b]/(a^2,b^2),\qquad
 j_*1=a+b.
\]
Every cohomology class in question is algebraic.  The left side of
\eqref{eq:bad-low-gysin} is the line $\mathbf Q(a+b)$, while its
right side is $\mathbf Qa\oplus\mathbf Qb$.  They are not equal.
The restriction in degree zero is an isomorphism and satisfies the
critical $A'$ condition.  Both $X$ and its zero-dimensional hyperplane
section satisfy $C_{\mathrm G}$: on $\mathbf P^1$, the fundamental
class of $\mathbf P^1\times\mathbf P^1$ is the correspondence taking
the degree-two generator to $1$ and vanishing on the other degrees.
Moreover $A(T)^0$ itself holds, since
$(a+b)^2\cdot1=2ab$.  Thus the failed weak condition cannot be
attributed to a failure of the genuine critical Lefschetz statement.

For an elliptic curve $E$, choose any sufficiently ample product
polarization on $E\times E$ and a smooth connected member $W$.
The two factor divisor classes are independent: they have squares
zero and intersection one.  Hence
$\dim_{\mathbf Q}\operatorname{Alg}^1(E\times E)\ge2$, whereas
$j_*H^0(W)$ has dimension one.  The same failure follows.  For a
curve, $C_{\mathrm G}$ holds directly: the only nonzero lowering
map is $H^2\to H^0$, induced, after division by the degree of the
polarization, by the fundamental product correspondence.
\end{proof}

We now give a replacement stated in terms of actual maps.  It uses
the final even-dimensional diagram in scan~050, but it makes its
extra $L_T$ part of the hypothesis.  This is a mathematical correction,
not a claim that the preceding definition of $A''$ in the manuscript
already contained that factor.

\begin{theorem}[SN-C03: the exact even critical factorization]
\label{thm:weighted-critical}
Let $T$ be smooth projective of dimension $2m\ge2$, and let
$j:W\hookrightarrow T$ be a smooth hyperplane section with class $L_T$.
Put $L_W=j^*L_T$.  Define the following conditions:
\begin{align*}
 R_0(T,W):\ &x\in H^{2m-2}(T),\quad
             j^*x\in\operatorname{Alg}^{m-1}(W)
             \ \Longrightarrow\ x\in\operatorname{Alg}^{m-1}(T);\\
 G_{L,0}(T,W):\ &L_Tj_*\operatorname{Alg}^{m-1}(W)
                    =\operatorname{Alg}^{m+1}(T);\\
 A_0(T):\ &L_T^2:\operatorname{Alg}^{m-1}(T)
                    \xrightarrow{\ \sim\ }
                    \operatorname{Alg}^{m+1}(T).
\end{align*}
Then
\begin{equation}
 A_0(T)\quad\Longleftrightarrow\quad
 R_0(T,W)\ \text{and}\ G_{L,0}(T,W).
 \label{eq:weighted-equivalence}
\end{equation}
No algebraicity of the inverse Lefschetz operator on $W$ is an
additional hypothesis in this equivalence.

For comparison, put
\[
 G_0(T,W):\quad
 j_*\operatorname{Alg}^{m}(W)=\operatorname{Alg}^{m+1}(T),
 \qquad
 A_0(W):\quad L_W:\operatorname{Alg}^{m-1}(W)
                       \xrightarrow{\ \sim\ }\operatorname{Alg}^{m}(W).
\]
There is an exact further equivalence
\begin{equation}
 G_{L,0}(T,W)\quad\Longleftrightarrow\quad
 G_0(T,W)\ \text{and}\ A_0(W).
 \label{eq:gysin-hidden-inverse}
\end{equation}
Thus replacing $G_{L,0}$ by ordinary upper critical Gysin
surjectivity drops precisely the indicated inverse-Lefschetz
assertion on the hyperplane $W$.
\end{theorem}
\begin{proof}
Weak Lefschetz and its Poincar\'e-dual form give isomorphisms
\[
 r=j^*:H^{2m-2}(T)\xrightarrow{\sim}H^{2m-2}(W),\qquad
 g=j_*:H^{2m}(W)\xrightarrow{\sim}H^{2m+2}(T).
\]
The projection and self-intersection formulas give
\begin{equation}
 L_T^2=gL_Wr=L_Tj_*r.
 \label{eq:critical-factorization}
\end{equation}
All the displayed maps preserve algebraic cycle classes whenever
source and target degrees match.

Suppose first that $A_0(T)$ holds.  If $r(x)$ is algebraic, then
$L_Tj_*r(x)=L_T^2x$ is algebraic.  Surjectivity of $A_0(T)$ supplies
an algebraic $x_0$ with $L_T^2x_0=L_T^2x$.  Hard Lefschetz makes
$L_T^2$ injective on all cohomology, so $x=x_0$.  This proves
$R_0(T,W)$.  If $\beta\in\operatorname{Alg}^{m+1}(T)$, write
$\beta=L_T^2x_0$ with $x_0$ algebraic.  Then $r(x_0)$ is algebraic
and \eqref{eq:critical-factorization} proves $G_{L,0}(T,W)$.

Conversely, assume $R_0$ and $G_{L,0}$.  Given an algebraic
$\beta$ in the upper critical degree, choose an algebraic
$y\in H^{2m-2}(W)$ with $L_Tj_*y=\beta$.  The isomorphism $r$
gives a unique $x$ with $r(x)=y$.  Condition $R_0$ makes $x$
algebraic, and \eqref{eq:critical-factorization} gives
$L_T^2x=\beta$.  Injectivity is again hard Lefschetz.  This proves
\eqref{eq:weighted-equivalence}.

Next assume $G_{L,0}$.  Since $L_Tj_*=j_*L_W$, it immediately
implies $G_0$.  For any algebraic $y\in H^{2m}(W)$, the class
$g(y)$ is algebraic.  Apply $G_{L,0}$ to choose algebraic
$z\in H^{2m-2}(W)$ with $g(y)=gL_Wz$.  Injectivity of $g$ gives
$y=L_Wz$.  Hard Lefschetz on $W$ gives the injectivity needed for
$A_0(W)$.  In the other direction, $G_0$ supplies, for an algebraic
$\beta$ on $T$, an algebraic $y$ with $g(y)=\beta$; $A_0(W)$
supplies an algebraic $z$ with $L_Wz=y$.  Therefore
$\beta=gL_Wz=L_Tj_*z$.  This proves
\eqref{eq:gysin-hidden-inverse}.
\end{proof}

\begin{corollary}[SN-C04: a usable corrected proposition for scan 050]
\label{cor:corrected-scan050}
Let $X$ be smooth projective of dimension $n\ge1$, let $Y\subset X$
be a smooth hyperplane section, and polarize $T=X\times X$ by
$L_T=L_X\otimes1+1\otimes L_X$.  Let $W\subset T$ be a smooth
hyperplane section for this polarization.  Assume the hard and weak
Lefschetz hypotheses of Kleiman's Theorem~2.13.  Then
\begin{equation}
 \boxed{\quad C_{\mathrm G}(X)\ \Longleftrightarrow\
 C_{\mathrm G}(Y)+R_0(T,W)+G_{L,0}(T,W).\quad}
 \label{eq:corrected-scan050}
\end{equation}
Along a flag of smooth hyperplane sections of $X$, the conjunction
of the two corrected weak conditions on each square is equivalent
to $C_{\mathrm G}(X)$.  The hyperplane $W\subset X\times X$
is not $Y\times Y$; their dimensions are $2n-1$ and $2n-2$,
respectively.
\end{corollary}
\begin{proof}
Apply Theorem~\ref{thm:weighted-critical} to $T$, with $m=n$.
It identifies the two displayed weak conditions with exactly the
critical condition in \eqref{eq:kleiman-critical}.  Kleiman's
Theorem~2.13(i), (ii$'$) now proves the assertion.  Iterating down
the flag and using the same theorem's (iv), or the fact that
$C_{\mathrm G}$ holds in dimension at most one, proves the final
assertion.  No implication from the uncorrected ordinary-Gysin
condition is used.
\end{proof}

\begin{proposition}[SN-C05: the upper-degree repair alone does not justify the diagram argument]
\label{prop:formal-critical-obstruction}
Restriction detection $R_0$ and ordinary upper critical Gysin
surjectivity $G_0$, together with the cohomological hard-Lefschetz
isomorphisms and the restriction--Gysin identities, do not formally
imply $A_0$.  This is a statement about those linear-algebra data,
not a geometric counterexample to a standard conjecture.
\end{proposition}
\begin{proof}
Take a Lefschetz module $V$ with vectors $v_0,v_1,v_2$ in degrees
$2m-2,2m,2m+2$ and $L_Vv_0=v_1$, $L_Vv_1=v_2$, $L_Vv_2=0$.
Take $U$ with vectors $u_0,u_1$ in degrees $2m-2,2m$, and
$L_Uu_0=u_1$, $L_Uu_1=0$.  Define
\[
 rv_0=u_0,\quad rv_1=u_1,\quad rv_2=0,
 \qquad gu_0=v_1,\quad gu_1=v_2.
\]
Then $gr=L_V$, $rg=L_U$, $rL_V=L_Ur$ and $gL_U=L_Vg$.
Both modules satisfy hard Lefschetz about their indicated centers,
and the low restriction and high Gysin maps are isomorphisms.
Declare their ``algebraic'' subspaces to be
\[
 V_{\mathrm{alg}}^{2m-2}=0,\quad
 V_{\mathrm{alg}}^{2m}=\mathbf Qv_1,\quad
 V_{\mathrm{alg}}^{2m+2}=\mathbf Qv_2,
 \qquad
 U_{\mathrm{alg}}^{2m-2}=0,\quad
 U_{\mathrm{alg}}^{2m}=\mathbf Qu_1.
\]
Every defined map preserves these subspaces.  Restriction detection
holds in the low degree because $r$ is injective and the target
algebraic subspace is zero.  The high Gysin map carries
$\mathbf Qu_1$ onto $\mathbf Qv_2$, so $G_0$ holds.  Nevertheless
$L_V^2$ takes the zero algebraic subspace in the low degree to zero,
not to $\mathbf Qv_2$; hence $A_0$ fails.  Equivalently,
$L_U^{-1}(u_1)=u_0$ is the missing algebraic preimage.
\end{proof}

\paragraph{Placement in Grothendieck's frameworks.}
The replacement \eqref{eq:corrected-scan050} is a statement about
rational algebraic correspondences and cycle classes, at precisely
cote~16 scans~034, 036, 040, 049--050.  Its restriction and Gysin
maps are the maps in weak Lefschetz (SGA~5, Expos\'e~VII,
Theorem~7.1, and the dual form), not information recovered from
Frobenius characteristic polynomials.  For a smooth model over an
open subscheme of $\operatorname{Spec}\mathbf Z$, the same maps
are restriction and proper pushforward on the corresponding
\(\ell\)-adic realizations, with the usual twists.  The
$\mathbf P^1\times\mathbf P^1$ counterexample and its smooth
$(1,1)$ diagonal are already defined over $\mathbf Z$, so the
rank failure occurs in every geometric fibre.  In rational
$K$-theory, the extra $L_T$ is multiplication by the first-Chern-class
operation in the associated graded for the gamma filtration;
forgetting it changes codimension by one.  In the category of
homological motives, $C_{\mathrm G}(X)$ supplies the inverse
Lefschetz morphisms.  None of SN-C01--SN-C05 alone asserts the
index inequality, positivity of an involution, or semisimplicity;
those are separate steps on scan~049.  Nor do these replacements
change the already-established sign-only purity theorem or any
local $L$-factor.


\section{The abelian cases of the index route}
\label{sec:abelian}

We use $C_{\mathrm G}(X)$ for the historical condition on cote~16:
the normalized lowering operator $\Lambda_X$ is algebraic.  This is
$B(X)$ in \cite[\S2]{Kleiman1968}.  It is stronger than the algebraicity
of the K\"unneth projectors, now usually called $C(X)$.
Throughout, an algebraic class means a rational cycle class modulo the
specified Weil homological equivalence.  The statements do not replace
homological equivalence by rational or algebraic equivalence.

\begin{proposition}[SN-10: explicit critical inverse maps in the abelian cases]
Let $A$ be an abelian variety of dimension $g$ over an algebraically closed
field, let $h$ be an ample rational divisor class, and use Betti cohomology
for a complex variety (or after spreading and embedding in $\mathbb C$),
or $\ell$-adic cohomology with $\ell$ invertible.
Then $C_{\mathrm G}(A)$ holds.  If $j:Y\hookrightarrow A$ is a smooth
hyperplane section for the chosen polarization, the following maps are
algebraic:
\[
 s=\Lambda_A j_*:H^k(Y)\longrightarrow H^k(A),\qquad
 t=j^*\Lambda_A:H^{k+2}(A)\longrightarrow H^k(Y).
\]
They satisfy $sj^*=1$ for $k\leq g-1$ and $j_*t=1$ for $k\geq g-1$.
The same statements apply to every power and product of abelian varieties.
Consequently the correctly indexed weak algebraicity conditions, as well
as the weighted critical Gysin conditions used in the corrected square
criterion, hold in these cases.  Moreover $C_{\mathrm G}(Y)$ holds.
\end{proposition}
\begin{proof}
The algebraicity assertion is the classical abelian Lefschetz theorem,
\cite[Theorem~2A11]{Kleiman1968}; in the stronger divisor-generated form it
is \cite[Theorem~5.9, with the author's correction to its proof]{Milne1999a}.
We cite that result rather than reprove it.  Here $\Lambda$ means unit
lowering on a Lefschetz string:
\[
 \Lambda(L^r v)=L^{r-1}v\ (r>0),\qquad \Lambda v=0
 \quad(v\text{ primitive}).
\]
No string terminates in degree less than $g$ and no string starts in degree
greater than $g$.  Therefore $\Lambda L=1$ on $H^k(A)$ for $k<g$, and
$L\Lambda=1$ on $H^{k+2}(A)$ for $k+2>g$.  The projection formula
$j_*j^*=L$ proves the two asserted identities, including the endpoints.
All factors are algebraic correspondences, so the identities give the
claimed descent and lifting of algebraic classes.  Explicitly, if $T$ has
even dimension $d$ and $j:W\hookrightarrow T$, a right inverse to
$L_Tj_*:H^{d-2}(W)\to H^{d+2}(T)$ is $j^*\Lambda_T^2$:
$L_Tj_*j^*\Lambda_T^2=L_T^2\Lambda_T^2=1$ in degree $d+2$.
Every Lefschetz string occurring there has at least two preceding steps,
which verifies the last equality.  Products remain abelian.  Finally,
$\Lambda_Y=j^*\Lambda_A^2j_*$ by cote~16 scan~040, equivalently
\cite[Proposition~2.12]{Kleiman1968}; this proves the last assertion.
\end{proof}

\begin{theorem}[SN-11: completion of the abelian index route]
Under the supplied CM Hodge theorem, every abelian variety $A$ over an
algebraically closed field satisfies $C_{\mathrm G}(A)$ and $I(A\times A)$.
In particular, for every rational ample $h$ on $A$ and every $2r\leq g$,
\[
 (x,y)\longmapsto (-1)^r\deg(xy h^{g-2r})
 \quad\text{on }\ker(h^{g-2r+1}:A^r_{\rm hom}(A)\to A^{g-r+1}_{\rm hom}(A))
\]
is positive definite over $\mathbb R$; homological and numerical
equivalence coincide.  The rational algebra of algebraic correspondences
modulo homological equivalence is semisimple, and its subalgebras stable
under the Lefschetz adjoint are semisimple.  Algebraic injections and
surjections between the cohomological summands of abelian varieties split
by algebraic correspondences.

For powers of a CM elliptic curve, including the hexagonal curve of
HODGE\_SIGN SG-04, the same conclusions use SG-04 and Milne's positivity
for Lefschetz classes and do not require the supplied general CM Hodge
theorem.
\end{theorem}
\begin{proof}
The first input is precisely the supplied 032, Theorem~1.1 (p.~1 and its
concluding proof on p.~47), equivalently van Hoften--Sempliner,
Hexagon~2610.00259v2, Theorem~I, p.~1.  Milne's published
Theorem~3.3(a),(b) gives the asserted equality of equivalence relations and
the index inequality in every characteristic; the attached 032,
Corollary~8.4, p.~50, states exactly this consequence, including rational
ample polarizations.  Apply it to $A\times A$ with $h\boxtimes1+1\boxtimes h$.
SN-10 supplies the other hypothesis of the classical positive-adjoint
theorem \cite[Theorem~3.11]{Kleiman1968}.  Its Corollaries~3.12 and~3.14
then give semisimplicity and splitting.  These are cited consequences,
not new proofs of those classical theorems.  All products of the abelian
varieties in a proposed correspondence are again abelian, so the product
positivity hypotheses in those corollaries are satisfied.

For elliptic powers, SG-04 identifies all algebraic class spaces with the
divisor-generated spaces, on the relevant complex varieties and their
good reductions.  \cite[Remarks~3.7--3.8]{Milne2002} supplies positivity on
these spaces.  The same application of SN-10 and Kleiman's corollaries
therefore applies.  This last argument cites SG-04 and does not reprove it.
\end{proof}

\paragraph{Exact disposition of Q2c.}
The intended positivity question Q2c is closed by the preceding cited
chain: 032 Theorem~1.1 $\Rightarrow$ Milne~2002 Theorem~3.3 $\Rightarrow$
032 Corollary~8.4.  The human deposit's \S1.1.2 uses
Theorem~4.6(a),(b), the numbering of Milne's original version
(the author's \texttt{2002aS.pdf}).  The later version corresponding to
the published paper has the needed conclusion as Theorem~3.3(a),(b)
(\texttt{2002aP.pdf}, pp.~8--9).  Both versioned locators have been
checked; this is a version crosswalk, not a mathematical discrepancy.

\subsection{A register correction: algebraic does not mean divisor-generated}

\begin{proposition}[SN-12: a concrete counterexample to the literal U-02 wording]
With the supplied 032 theorem, there is a simple complex CM abelian
fourfold with algebraic codimension-two classes that are not generated by
divisors.  Thus the predecessor's literal U-02 sentence
``Every algebraic class on an arbitrary CM abelian variety is a Lefschetz
class'' is false, even though the intended Q2c positivity question is
closed by SN-11.
\end{proposition}
\begin{proof}
Let $F=\mathbb Q(\alpha)$, where
\[
 f(T)=T^4-5T^2+T+1,\qquad K=\mathbb Q(i),\qquad E=FK.
\]
The signs of $f$ give four distinct real roots, respectively in
$(-3,-2),(-1,0),(0,1),(2,3)$.
Modulo $3$, $f=T^4+T^2+T+1$ has no linear root.  Its remainders modulo the
three monic irreducible quadratics over $\mathbb F_3$,
$T^2+1,T^2+T+2,T^2+2T+2$, are respectively
$T+1,1,2T+1$.  It is therefore irreducible modulo $3$.
Modulo $2$, its factorization is
$(T+1)(T^3+T^2+1)$, with the cubic irreducible and the factors distinct.
The Galois group of its splitting field $N$ is thus a transitive subgroup
of $S_4$ containing a $4$-cycle and a $3$-cycle, hence $S_4$.
The field $N$ is totally real, so $N\cap K=\mathbb Q$.

The only proper CM subfield of $E$ is $K$.  Here is the subgroup check,
to make this assertion explicit.  In
$\operatorname{Gal}(NK/\mathbb Q)=S_4\times C_2$ the subgroup fixing $E$
is $S_3\times\{1\}$.  A subgroup containing it projects to either $S_3$
or $S_4$.  In the first case it is $S_3\times\{1\}$ or $S_3\times C_2$.
In the second case it is the full group or the graph of a homomorphism
$S_4\to C_2$ trivial on $S_3$; that homomorphism is trivial, so the subgroup
is $S_4\times\{1\}$.  The intermediate fields are therefore just
$E,F,K,\mathbb Q$.

Choose a CM type of $E$ with two embeddings above each embedding of $K$.
It cannot be induced from $K$, so it is primitive.  The CM classification
\cite[Propositions~3.12--3.13]{MilneCM} produces a simple abelian fourfold
$A$ with $\operatorname{End}^0(A)=E$ of this type.  The action of $K$ on
$H^{1,0}(A)$ has signature $(2,2)$.  Hence the rational two-dimensional
Weil subspace
\[
 W=\bigwedge_K^4 H^1(A,\mathbb Q)\ \subset\ H^4(A,\mathbb Q)
\]
has Hodge type $(2,2)$.  The inclusion here is the canonical determinant
subspace: after splitting $K$, its two lines are
$\bigwedge^4 V_+$ and $\bigwedge^4 V_-$ for the two $K$-eigenspaces in
$H^1(A,\mathbb C)$.

Rosati on the CM field $E$ is complex conjugation.  The
N\'eron--Severi/Rosati identification
\cite[Proposition~13.1]{MilneAV} gives
$\dim_{\mathbb Q}\operatorname{NS}(A)_{\mathbb Q}=\dim_{\mathbb Q}F=4$.
In the decomposition of $H^1$ into the eight $E$-embedding lines,
the four divisor directions pair an embedding with its conjugate.
For clarity, this last description does not follow merely by a dimension
count: a polarization pairs only conjugate lines, because
$\psi(ax,y)=\psi(x,\bar a y)$, and multiplying it by $F$ gives the four
independent such pairs.  Every product of two divisor classes consequently
has $K$-bidegree $(2,2)$.  The two lines of $W$ have bidegrees $(4,0)$ and
$(0,4)$.  Their intersection with the divisor-generated subspace is zero.
Every nonzero rational element of $W$ is algebraic by 032 Theorem~1.1.
This gives the asserted counterexample.
\end{proof}

\section{The additional inputs for the positive-adjoint argument}
\label{sec:positive}

Let $R_X$ denote the image of rational algebraic correspondences in
$\operatorname{End}(H^*(X))$.  It is the image algebra, so classes acting
as zero have already been quotiented out.  Under $C_{\mathrm G}(X)$,
the Lefschetz star and degree projectors are algebraic.  The adjoint
$u\mapsto u^\dagger$ with respect to the Lefschetz pairings preserves
$R_X$, and
\[
  \tau(u)=\sum_i\operatorname{Tr}(u_{ii}\mid H^i(X))\in\mathbb Q.
\]
The sum here is the ordinary trace, not the alternating trace.  The degree
projectors convert the cohomological Lefschetz trace formula to this
rational trace.  The extra condition needed by scan~049 is exactly
\begin{equation}\label{eq:positive-trace}
             \tau(u^\dagger u)>0\qquad(0\ne u\in R_X).
\end{equation}
More invariantly, the associated symmetric form on $R_X\otimes\mathbb R$
must be positive definite.  Kleiman~1968, Theorem~3.11, derives precisely
this from $I(X\times X)$ with the product polarization.  Positivity on
only divisor classes or on only the scalar Frobenius traces is not this
condition.

In an $\ell$-adic realization, this positivity also supplies the
finite-dimensionality over $\mathbb Q$ needed below. Indeed, if
$u_1,\ldots,u_s$ are rationally independent in $R_X$, their rational
Gram matrix $\bigl(\tau(u_i^\dagger u_j)\bigr)$ is positive definite
and hence nonsingular. Any $\mathbb Q_\ell$-linear relation among
their cohomological operators, paired with each $u_i^\dagger$, is
therefore zero only with all coefficients zero. Thus
$s\le(\dim_{\mathbb Q_\ell}H^*(X))^2$. This avoids inferring
rational finite-dimensionality merely from an inclusion in a
finite-dimensional $\mathbb Q_\ell$-space.

\begin{proposition}[SN-13: constructive splitting from the positive trace form]
Let $R$ be a finite-dimensional rational algebra acting faithfully on
a vector space over a characteristic-zero field, equipped with an
involution $\dagger$ and a cyclic rational trace $\tau$ satisfying
$\tau(a^\dagger)=\tau(a)$ and such that
$\tau(a^\dagger b)$ is positive definite over $\mathbb R$.
Let $e,f$ be self-adjoint idempotents
and $u=fue$.  If $u:eV\to fV$ is injective, then
\[
 v=(u^\dagger u)^{-1}u^\dagger\in eRf,
 \qquad vu=e,
\]
where the inverse is taken in $eRe$.  If $u$ is surjective, the analogous
right inverse is $u^\dagger(uu^\dagger)^{-1}$.
\end{proposition}
\begin{proof}
On $R_{\mathbb R}$ use the inner product
$\langle a,b\rangle=\tau(a^\dagger b)$.
It is symmetric because $\tau(a^\dagger)=\tau(a)$, and is positive
definite by hypothesis; both hold for the geometric trace form above.
Left multiplication by $u^\dagger$ is the adjoint of left multiplication
by $u$.  If $a\in eR$ and $ua=0$, faithfulness of the action on $V$ and
injectivity on $eV$ imply $a=0$.  Thus $u^\dagger u$ acts invertibly on
$eR$, since
$\langle a,u^\dagger ua\rangle=\langle ua,ua\rangle>0$ for $a\ne0$.
Its minimal polynomial on this finite-dimensional rational vector space
has nonzero constant coefficient.  Bezout with $T$ gives an inverse which
is a rational polynomial in $u^\dagger u$ with identity $e$, hence belongs
to $eRe$.  Multiplication verifies $vu=e$.  For a surjection,
$a\mapsto au$ is injective on $Rf$: if $au=0$, then $a$ vanishes on
$fV=\operatorname{im}(u)$ and $a=af$, so $a=0$.
Cyclicity shows that right multiplication by $u^\dagger$ is the adjoint
of right multiplication by $u$.  Thus right multiplication by
$uu^\dagger$ is invertible on $Rf$.  The same polynomial argument gives
its inverse in $fRf$, and multiplication verifies the stated right inverse.
\end{proof}

\begin{theorem}[SN-14: trace and $\lambda$ identities do not supply scan-049 positivity]
Even if all coefficients from a specified correspondence algebra are
allowed in traces, pure Frobenius eigenvalues, effective Grothendieck-group
identities, and their tensor, exterior-power and dual consequences do not
formally imply semisimplicity of that algebra or
\eqref{eq:positive-trace}.  They also do not imply semisimplicity of a
pure sheaf's Frobenius at a point.  These are counterexamples to a formal
implication, not counterexamples to the standard conjectures for varieties.
\end{theorem}
\begin{proof}
Work in the rigid tensor category of rational representations of
$\mathbb G_a\times\mathbb G_m$.  Take $V=\mathbb Q^3$ and
\[
 N=\begin{pmatrix}0&1&0\\0&0&1\\0&0&0\end{pmatrix},\qquad
 J=\begin{pmatrix}0&0&1\\0&-1&0\\1&0&0\end{pmatrix},
\]
with action $(t,s)\mapsto s\exp(tN)$.  Its endomorphism algebra is
$R=\mathbb Q[N]\simeq\mathbb Q[\epsilon]/(\epsilon^3)$:
an endomorphism commutes with the action exactly when it commutes with
the single Jordan block $N$, whose centralizer consists of its
polynomials.  Put $F=c\exp(N)$ with $c=q^r>0$ for an integer $r\geq0$.
All Frobenius eigenvalues have pure weight $2r$, while both $F$ and
the algebra $R$ are nonsemisimple.  The $J$-adjoint satisfies
$N^\dagger=-N$, so $R$ is closed under it and $F^\dagger F=c^2I$.
Nevertheless $\tau(N^\dagger N)=-\operatorname{Tr}(N^2)=0$.
For every $a=a_0I+a_1N+a_2N^2$ and every integer $s\geq0$,
\[
 \operatorname{Tr}(aF^s)=3a_0c^s.
\]
Thus all these coefficient traces factor through $R/\operatorname{rad}R$.

More generally the invariant filtration
$0\subset\ker N\subset\ker N^2\subset V$
has three one-dimensional trivial unipotent quotients.  It is compatible with
$F$, and the standard filtrations of tensor products, exterior powers,
and duals have associated graded built from the same quotients.  Traces
are additive in these filtrations.  Consequently all the indicated
character identities, including insertions of the induced coefficient
operators, are identical to those on the associated graded.
In the Grothendieck $\lambda$-ring of this tensor category this is simply
$[V]=3[\mathbb Q_c]$ and
$\lambda_t([V])=(1+[\mathbb Q_c]t)^3$.
The notation $\mathbb Q_c$ denotes the one-dimensional torus character
evaluated at $F$, where it has eigenvalue $c$.
Even effectivity does not help: $V$ itself is an actual object.

For the sheaf assertion, on $\operatorname{Spec}\mathbb F_q$ take the
continuous $\ell$-adic representation of $\widehat{\mathbb Z}$ defined by
$z\mapsto\exp(z_\ell N)$, where $z_\ell$ is its projection to
$\mathbb Z_\ell$.  This is a continuous homomorphism into
$\operatorname{GL}_3(\mathbb Q_\ell)$, since the exponential is the finite
polynomial $I+z_\ell N+z_\ell^2N^2/2$.
Its Frobenius is $\exp(N)$ (or its inverse under the other
Frobenius convention).  At every extension point the eigenvalues are all
$1$, so this is a punctually pure sheaf of weight zero.  It is not
semisimple as a Frobenius representation and has
$\det(1-TF)=(1-T)^3$, identical to the split sheaf.
\end{proof}

\paragraph{Application to the sign-only theorem.}
HODGE\_SIGN SG-08--SG-11, together with GCT\_NSRH F03, is used as the
settled sign-only theorem stipulated in the prompt.  Its conclusion and
effective remainder are not reproved.  Applying its trace identities to
a product changes neither the preceding filtration argument nor the
distinction between a correspondence and a lisse coefficient sheaf.
A correspondence acting on $H^*(X)$ does not automatically preserve a
chosen generic hyperplane and its Gysin kernel.  Even when such
functorial compatibility is furnished, the trace-only information in
SN-14 remains insensitive to the radical.  To complete the
scan-049/Kleiman argument by this method, one must additionally establish
the algebraic adjoint and the nondegenerate positive form
\eqref{eq:positive-trace}; asserting either from the parity of an
effective remainder would assume the missing conclusion.  SN-14 does
not exclude a different semisimplicity argument.

\paragraph{Motivic and arithmetic formulations.}
For comparison, numerical motives with rational coefficients are already
a semisimple abelian category by Jannsen,
\emph{Motives, numerical equivalence, and semi-simplicity},
Theorem~1 \cite{Jannsen1992}.  That cited theorem concerns numerical
equivalence.  It does not identify numerical equivalence with the
homological equivalence used by the given realization, nor does it
provide the positive Lefschetz-adjoint form on that realization.
Those are the extra assertions supplied in the abelian case here.

SN-11 gives the semisimple polarized category of abelian motives modulo
$\ell$-adic homological equivalence, identified with its numerical
category on all products.  Over a finite field the Lefschetz adjoint makes
normalized Frobenius unitary, so it is semisimple; the classical exact
locator is \cite[Lemma~4.6 and Theorem~4.7]{Kleiman1968}.
The absolute values already come from the supplied sign-only theorem.
Splittings of motives split their local characteristic polynomials and
therefore multiply their local $L$-factors.  SN-14 explains why the
converse from these factors or their $K_0$ classes fails.

For a finite collection of the varieties, polarizations, cycles and
operators in SN-10--SN-13, choose a finitely generated ring over
$\mathbb Z$ carrying the data and shrink to a smooth projective model;
finite-presentation descent is EGA~IV$_3$, 8.8.2, and smooth proper
cohomological specialization is SGA~4, XVI, 2.2.
The realization consists of the lisse sheaves $R^i f_*\mathbb Q_\ell$ and
the corresponding morphisms of sheaves.  The trace/correspondence
operations are SGA~5, III, 3 and 5.1--5.3; the $\lambda$-ring identity in
SN-14 uses the exact-sequence rule SGA~6, V, 2.1--2.2, especially
Lemma~2.2.1.  A $K_0$ class
records that rule but does not carry a polarization or an integral
splitting.  This identifies the precise place of these findings in
cote~16 scans~040, 049--050 and in the predecessor's U-01/U-02, without
promoting them to claims about Chow groups.


% Contribution by /root/integral_fock, GPT-6 Astra Pro; effort not exposed.
% Findings IF-01--IF-11. This file is intended for inclusion after amsmath,
% amssymb and amsthm are loaded. No external macros other than standard LaTeX.
\section{The integral Fock comparison on the actual exterior lattice}
\label{sec:integral-fock}

Put $M_n=\bigoplus_{a=1}^n(\mathbb Z e_a\oplus\mathbb Z f_a)$,
$V_n=\bigwedge^*M_n$, $p_a=e_a\wedge f_a$, and
\[
 L=\Bigl(\sum_a p_a\Bigr)\wedge-,\qquad
 D=\sum_a\iota_{f_a^\vee}\iota_{e_a^\vee}.
\]
The lowering here is ordinary contraction, not the inverse-string lowering
of HODGE\_SIGN SL-08. All lattices below are specified over $\mathbb Z$.

\subsection{IF-01: an explicit integral Fock intertwiner}
For a subset $S\subseteq\{1,\ldots,n\}$ and choices
$\epsilon_a\in\{e_a,f_a\}$ for $a\in S$, let
$\mathcal P_{S,\epsilon}$ be the submodule of
$\mathbb Z[x_a,y_a:a\notin S]$ consisting of polynomials homogeneous of
column degree one in each pair $(x_a,y_a)$. Give its monomial basis the
factorial Fock pairing. This basis is orthonormal, since every exponent is
zero or one. Set
\[
 \mathcal F_n=\bigoplus_{S,\epsilon}\mathcal P_{S,\epsilon}.
\]
Assign a monomial containing $t$ of the $y$'s exterior degree $|S|+2t$.
On each summand set
\[
 \mathscr L=\sum_{a\notin S}y_a\partial_{x_a},\qquad
 \mathscr D=\sum_{a\notin S}x_a\partial_{y_a}.
\]

\paragraph{Statement IF-01.}
There is a graded integral isometry $\Phi_n:\mathcal F_n\longrightarrow V_n$
intertwining $(\mathscr L,\mathscr D)$ with $(L,D)$ and all their divided
powers. It is given, with factors in increasing pair order, by
\[
 \Phi_n\left(\prod_{a\notin S\cup T}x_a\prod_{a\in T}y_a\right)
 =\prod_{a=1}^n u_a,\qquad
 u_a=\begin{cases}\epsilon_a&a\in S,\\p_a&a\in T,\\1&\text{otherwise}.
 \end{cases}
\]
In particular this is a specified lattice-preserving Fock realization of
Lefschetz divided powers, not merely an equality of rational characters.

\begin{proof}
A standard exterior basis vector chooses, in each symplectic pair, one of
$1,e_a,f_a,p_a$. The displayed assignment is therefore a bijection of
integral bases and preserves their declared pairings and degrees.
The operator $y_a\partial_{x_a}$ replaces the unique factor $x_a$ by
$y_a$, if present, and otherwise kills the monomial. Under $\Phi_n$ this
inserts $p_a$ in an empty pair and vanishes on a singleton or complete
pair. Because $p_a$ has even degree, moving it to its prescribed pair
position introduces no sign. Similarly
$\iota_{f_a^\vee}\iota_{e_a^\vee}$ removes a complete pair with sign $+1$,
and kills the other three states. This proves the two intertwining
identities. The commuting operators for distinct pairs square to zero;
hence
\[
 L^r/r!=\sum_{|R|=r}\prod_{a\in R}(p_a\wedge-)
\]
is integral, and the corresponding expression holds for $D$ and for the
two polynomial operators. Equivalently the coefficient of $z^r$ in the
simultaneous row substitution $x_a\mapsto x_a+zy_a$ is
$\mathscr L^r/r!$. Thus the isometry also respects the actual coefficient
normalization of divided row transvections.
\end{proof}

The exact source comparison is Alp\"oge,
\path{Formalism/GradedWeylDivisibility/StatementA.lean}, declaration
\texttt{rowDiv}: it is coefficient extraction under a row transvection.
The present module is its explicitly specified multiaffine specialization.
It does not identify a global theta-lift integral with that transvection.
On the exterior side this is the model in cote 16, scan 014, paragraph 7;
its contraction must remain distinguished from the inverse identities on
scan 016, as HODGE\_SIGN SL-04 already demonstrates.

\subsection{IF-02: the divided-power order is saturated for every rank}
This is the normalization-specific type-$A_1$ coefficient-coalgebra
statement for the actual exterior lattice. We claim no novelty for
the general integral Schur-algebra framework.
Write
\[
 \mathcal O_n=\mathbb Z[L,D],\qquad
 \mathcal B_n=\mathbb Z[L^{(r)},D^{(r)}:r\geq0],\qquad
 \mathcal A_n^{\rm ext}=\mathbb Q[L,D]\cap\operatorname{End}_{\mathbb Z}(V_n),
\]
where $L^{(r)}=L^r/r!$ and similarly for $D$; all powers with $r>n$
vanish. The intersection is taken inside
$\operatorname{End}_{\mathbb Q}(V_n\otimes\mathbb Q)$.

\paragraph{Statement IF-02.}
For every $n\geq0$,
\[
 \mathcal B_n=\mathcal A_n^{\rm ext},\qquad
 \operatorname{rank}_{\mathbb Z}\mathcal A_n^{\rm ext}
 =\sum_{v=0}^n(v+1)^2=\frac{(n+1)(n+2)(2n+3)}6.
\]
All exterior degree projectors belong to $\mathcal B_n$. For $n\geq1$
the additive index
$[\operatorname{End}_{\mathbb Z}(V_n):\mathbb Z[L,D]]$ is infinite.
Consequently the finite order-index question in the intake has to be
posed in $\mathcal A_n^{\rm ext}$, not in the whole endomorphism ring.

\begin{proof}
We give a coefficient-lattice proof, including the integral saturation
step. Put $U=\mathbb Z^2$ with its standard $\mathrm{SL}_{2,\mathbb Z}$
action. IF-01 identifies $V_n$, for this action, with a direct sum of
copies of $U^{\otimes v}$, $0\leq v\leq n$: the multiplicity is
$\binom n{n-v}2^{n-v}$. Every $v$ occurs. Repeated identical blocks do
not change either the image order or its saturation, so we may use
$W_n=\bigoplus_{v=0}^n U^{\otimes v}$.

Let
\[
 R=\mathbb Z[a,b,c,d]/(ad-bc-1),\qquad
 R_n=\langle\text{monomials of total degree at most }n\rangle_{\mathbb Z}.
\]
The matrix coefficients of $W_n$ span exactly $R_n$: the coefficients
of $U^{\otimes v}$ are all degree-$v$ monomials in the four matrix
entries. The relation $ad=bc+1$ reduces every element to a linear
combination of monomials not containing both $a$ and $d$. These normal
monomials are independent, as will also follow from the triangular
calculation below. The number of normal monomials of degree $v$ is
\[
 \binom{v+3}{3}-\binom{v+1}{3}=(v+1)^2.
\]
The coproduct of $R$ carries $R_n$ into $R_n\otimes R_n$. Therefore
$S_n=\operatorname{Hom}_{\mathbb Z}(R_n,\mathbb Z)$ is an algebra, acting
faithfully on $W_n$ by evaluating its matrix coefficients. Faithfulness
uses that those coefficients span $R_n$.

We show that the divided root operators span $S_n$ integrally. On the
big cell write
\[
 g=\begin{pmatrix}1&0\\x&1\end{pmatrix}
   \begin{pmatrix}\tau&0\\0&\tau^{-1}\end{pmatrix}
   \begin{pmatrix}1&y\\0&1\end{pmatrix}
 =\begin{pmatrix}\tau&\tau y\\\tau x&\tau xy+\tau^{-1}\end{pmatrix}.
\]
The normal monomials have the following expansions:
\[
 a^r b^s c^t\mapsto x^t y^s\tau^{r+s+t}\quad(r\geq0),
\]
\[
 b^s c^t d^u\mapsto
 \sum_{j=0}^u\binom uj x^{t+j}y^{s+j}\tau^{s+t-u+2j}
 \quad(u\geq1).
\]
For the first family, the exponent of $\tau$ minus the sum of the
exponents of $x,y$ is $r\geq0$. For the second it is the fixed negative
integer $-u$. In the second family the $j=0$ term has coefficient one;
all other terms increase $s+t$ by $2j$. Thus the coefficient submatrix
on the $j=0$ monomials, ordered by $s+t$, is triangular with diagonal
one, separately for each $u$. The first family contributes an identity
matrix. This proves independence of the normal monomials and shows that
$R_n$ is a primitive sublattice of the Laurent-polynomial coefficient
lattice. More precisely, restriction of coefficient-extraction
functionals spans the full dual $S_n$. Every appearing power of
$\tau$ lies between $-n$ and $n$, and the powers of $x,y$ lie between
zero and $n$.

Let $E=D$, $F=L$, and let $H$ act on $U$ with weights $+1,-1$.
Coefficient extraction in $x^r y^s z^k$ after setting $\tau=1+z$
acts as
\[
 F^{(r)}\binom Hk E^{(s)}.
\]
For exponents $q=-n,\ldots,n$, the functions $q\mapsto\binom qk$,
$0\leq k\leq2n$, form an integral basis of all functions on this finite
set. Indeed, $\binom{q+n}{k}$ gives a lower triangular matrix with
diagonal one, and Vandermonde's identity changes it to $\binom qk$
by an integral triangular matrix with diagonal one. Hence extraction of
an individual Laurent coefficient $\tau^q$ is an integral combination
of these coefficient functionals.

For completeness, the toral operators $\binom Hk$ lie in the algebra
generated by the divided root operators. The matrix identity
\[
 \begin{pmatrix}1&u\\0&1\end{pmatrix}
 \begin{pmatrix}1&0\\v&1\end{pmatrix}
 =\begin{pmatrix}1&0\\v/(1+uv)&1\end{pmatrix}
  \begin{pmatrix}1+uv&0\\0&(1+uv)^{-1}\end{pmatrix}
  \begin{pmatrix}1&u/(1+uv)\\0&1\end{pmatrix}
\]
expanded as a formal power series at $u=v=0$ gives
\[
 E^{(r)}F^{(s)}=
 \sum_{j=0}^{\min(r,s)}F^{(s-j)}
 \binom{H-r-s+2j}{j}E^{(r-j)}. \tag{IF-02a}
\]
To check the coefficient, choose powers $s-j$ and $r-j$ from the
left and right unipotent factors; their two denominators leave
$(1+uv)^{H-r-s+2j}$, and its $j$th coefficient is the displayed
binomial. For $r=s=k$, the final term is $\binom Hk$. Every earlier
term involves a shifted binomial of degree $j<k$, an integral
combination of $\binom Hh$, $h\leq j$, by Vandermonde's identity
(including its integral negative-shift form). Induction on $k$ proves
the assertion. Thus all coefficient functionals just constructed belong
to $\mathcal B_n$, proving $\mathcal B_n=S_n$ as embedded orders.

Over $\mathbb Q$, ordinary and divided powers generate the same algebra,
so $\mathbb Q[L,D]=S_n\otimes\mathbb Q$. A rational functional on
$R_n$ acts integrally on $W_n$ if and only if it is integer-valued on
every matrix coefficient, equivalently on $R_n$. This proves
$\mathcal A_n^{\rm ext}=S_n=\mathcal B_n$. A degree projector is the
indicator of one weight of $H$, since exterior degree $j$ has weight
$n-j$. The same integral binomial interpolation puts it in
$\mathcal B_n$. Finally $\operatorname{End}_{\mathbb Z}(V_n)$ has rank
$(4^n)^2=16^n$, strictly larger than $\sum_{v=0}^n(v+1)^2$ for
$n\geq1$. The rank comparison proves that its index over $\mathcal O_n$
is infinite.
\end{proof}

This proof also identifies a precise integral object in Grothendieck's
language: the divided-power image is the dual of the finite coefficient
coalgebra $R_n$ of the split group scheme $\mathrm{SL}_{2,\mathbb Z}$.
The exterior module is $\bigoplus_i\lambda^i(M_n)$ in the representation
$\lambda$-ring. No rational decomposition into primitive strings was
used to choose an integral lattice.

\subsection{IF-03: the actual exterior order indices through rank four}
Let $\mathcal O_n^\pi=\mathbb Z[L,D,\pi_0,\ldots,\pi_{2n}]$.
The exact indices in the canonical saturation are
\[
\begin{array}{c|r|c|c|c}
 n&\operatorname{rank}\mathcal A_n^{\rm ext}
 &[\mathcal A_n^{\rm ext}:\mathcal O_n]
 &[\mathcal A_n^{\rm ext}:\mathcal O_n^\pi]
 &[\mathcal A_n^{\rm ext}:\mathcal B_n]\\\hline
 1&5&1&1&1\\
 2&14&2^7&2^2&1\\
 3&30&2^{27}3^4&2^{10}3^2&1\\
 4&55&2^{81}3^{33}&2^{32}3^{10}&1
\end{array} \tag{IF-03a}
\]
The nonzero embedded Smith factors, including units, are
\[
\begin{array}{c|l|l}
 n&\mathcal O_n&\mathcal O_n^\pi\\\hline
 1&1^5&1^5\\
 2&1^8,2^5,4&1^{12},2^2\\
 3&1^8,2^{17},4,12^4&1^{20},2^8,6^2\\
 4&1^8,2^{19},6^5,12^{14},24^4,72^4,288
   &1^{28},2^{17},6^7,12,24^2.
\end{array} \tag{IF-03b}
\]
Here exponents on the Smith factors mean multiplicities. These are
indices on the actual exterior lattice. They are not the indices
$1,2^5,2^{13}3^{10}$ on the different ordinary-power string lattice
of HODGE\_SIGN SL-09.

\begin{proof}[Exact computation certificate, with an elementary rank-two check]
We specify both the finite algorithm and its certificate. In IF-01,
retain one copy of each singleton block, that is, the direct sum of the
Boolean modules on subsets of $\{1,\ldots,v\}$, $0\leq v\leq n$.
There are $\sum_{v=0}^n4^v$ matrix coordinates. On the $v$th block,
\[
 L_{T,S}=\begin{cases}1&S\subset T,\ |T\setminus S|=1,\\0&\text{otherwise},\end{cases}
 \qquad D=L^{\mathsf t},\qquad
 (\pi_j)_{S,S}=\boldsymbol1_{n-v+2|S|=j}.
\]
Repeated singleton blocks have identical integer matrices and hence do
not affect integral saturation. Starting from the column vector of the
identity, repeatedly take the column Hermite normal form of
\[
 [B,\ R_LB,\ R_DB]
 \quad\text{or}\quad
 [B,\ R_LB,\ R_DB,\ R_{\pi_0}B,\ldots,R_{\pi_{2n}}B],
\]
where $R_g$ means right multiplication by the specified integral matrix.
By induction this is exactly the lattice spanned by words in the
generators of bounded length. At equality of two successive HNF bases
it is closed under every generator, contains the identity, and is
therefore exactly the generated order. Equality is equality of integer
matrices, not just equality of ranks. Applying Smith reduction to this
embedded lattice gives its index in its saturation as the product of the
nonzero Smith factors. IF-02 identifies that saturation independently.

The complete bases, the stabilization histories, and their Smith factors
are the files \path{order_basis_nN_KIND.json}; the algorithm is
\path{exterior_order_indices.py}. Each of the sixteen computed orders
is checked for closure under every generator. The independent program
\path{verify_exterior_orders.py} constructs the full signed exterior
matrices directly, verifies the generator normalization, and checks the
nontrivial local determinantal valuations by $p$-local elimination at
$2,3,5$, without using the Smith routine. Its output is
\path{independent_exterior_verification.json}. A separate check
by the root agent, \path{work/root/verify_left_closure.py},
starts from independently constructed full signed exterior matrices,
generates left words, and compares all sixteen complete HNF lattices
for $n=1,2,3,4$ with the right-word certificates. It passed; its
output is \path{work/root/left_closure_verification.json}. The
two closure computations share SymPy HNF infrastructure. The factors in
(IF-03b) are the exact integer outputs of the stated terminating
algorithm; multiplication gives (IF-03a). This is a computation
certificate, not a Lean formalization.

For an elementary check on the normalization-sensitive case $n=2$, take
the three-dimensional string with basis
$v_0=1$, $v_1=p_1+p_2$, $v_2=p_1p_2$, and the trivial vector
$a=p_1-p_2$. On this string
$Lv_0=v_1$, $Lv_1=2v_2$, $Dv_1=2v_0$, $Dv_2=v_1$.
Write $A=(a_{rs})_{0\leq r,s\leq2}$ for its rational operator and
$z$ for the operator on $a$. The independent trivial singleton
block forces $z\in\mathbb Z$. Preservation of the lattice with
$p_1=(v_1+a)/2$, $p_2=(v_1-a)/2$ is exactly
\[
 a_{01},a_{21}\in2\mathbb Z,\qquad
 a_{11}\equiv z\pmod2,
\]
with all other matrix entries integral. This is the even part of
$\mathcal A_2^{\rm ext}$; the odd part is a full $2$-by-$2$ matrix ring.
The degree projectors and $L,D$ generate this even order with the two
additional conditions $a_{02},a_{20}\in2\mathbb Z$. Indeed the endpoint
projectors generate
$E_{10},2E_{21},E_{12},2E_{01},2E_{20},2E_{02}$,
and their products give $2E_{11}$; the middle projector supplies
$(E_{11},1)$. These are precisely a basis subject to the displayed
conditions and the two additional divisibilities. Thus the index with
degree projectors is $2^2$, agreeing with the certificate and showing
explicitly why the string-model value $2^5$ does not apply.
\end{proof}

\subsection{IF-04: the full Weyl lattice and its dual in an exterior model}
The multiaffine construction can be extended to the actual polynomial
Weyl lattices in the source. This also locates their denominators.

Let $U=\mathbb Z^r$ and let $\lambda=(\lambda_1\geq\cdots\geq
\lambda_r\geq0)$ be a partition. Let $h_1,\ldots,h_q$ be its column
heights, where $q=\lambda_1$. Set
\[
 T^{\rm amb}_\lambda=\bigotimes_{j=1}^q\bigwedge^{h_j}U,
 \qquad t_\lambda=\bigotimes_{j=1}^q(e_1\wedge\cdots\wedge e_{h_j}).
\]
Give this tensor lattice the orthonormal exterior-basis pairing. In the
polynomial ring on an $r$-by-$r$ matrix let $\Delta_h$ be its upper-left
$h$-minor and put $p_\lambda=\prod_j\Delta_{h_j}$.
Let $T_\lambda$ be the lattice generated from $t_\lambda$ by all divided
row-root operators and let $\Delta_\lambda$ be the lattice generated
from $p_\lambda$ by those operators. The latter is exactly
\texttt{weyl} in \path{GradedWeylDivisibility/StatementA.lean},
with \texttt{hw} as its initial vector. A determinant twist is removed
as in \path{LocalRingCardinality/DataCoeff.lean}, declarations
\texttt{DomWeight.detShift} and \texttt{DomWeight.natW}.

\paragraph{Statement IF-04.}
Multiplying the initial minors gives an integral isomorphism
\[
 \pi:T_\lambda\xrightarrow{\ \sim\ }\Delta_\lambda.
\]
It is an isometry when the polynomial Fock pairing is divided by
$\langle p_\lambda,p_\lambda\rangle_{\rm Fock}$, precisely the source's
\texttt{normFockQ}. Fix a simple row root $i,i+1$ and put
$d=\lambda_i-\lambda_{i+1}$. Its cyclic string has basis
\[
 q_s=F_i^{(s)}t_\lambda\quad(0\leq s\leq d),
\]
with
\[
 \langle q_s,q_t\rangle=\boldsymbol1_{s=t}\binom ds,
 \quad F_i^{(k)}q_s=\binom{s+k}{k}q_{s+k},
 \quad E_i^{(k)}q_s=\binom{d-s+k}{k}q_{s-k},                 \tag{IF-04a}
\]
where out-of-range vectors are zero. The weight-line lattice dual to
$\mathbb Z q_s$ is $\mathbb Z q_s/\binom ds$. Moreover, this entire
root action has an explicit integral realization by ordinary Lefschetz
operators on a symplectic exterior lattice; its Weyl string is primitive
in that exterior lattice. The dual string generally is not contained in
that same exterior lattice.

\begin{proof}
Map $e_{i_1}\wedge\cdots\wedge e_{i_h}$ in a column factor to the minor
with rows $i_1,\ldots,i_h$ and the first $h$ columns, and multiply these
polynomials over all column factors. This defines an integral linear map
$\pi:T^{\rm amb}_\lambda\to\mathbb Z[X_{ab}]$, equivariant for row
substitutions, sending $t_\lambda$ to $p_\lambda$.

We prove injectivity on the specified cyclic lattice without presuming
an integral decomposition into irreducibles. For each initial set of
$k$ rows, a wedge of degree $h$ has at most $\min(k,h)$ entries in those
rows. The row weight of $t_\lambda$ attains the sum of these bounds for
every $k$. Equality for every $k$ forces each wedge factor to be its
initial wedge. Thus its weight line in $T^{\rm amb}_\lambda\otimes\mathbb Q$
is one-dimensional. Ordinary opposite row-root operators are adjoint
for the tensor exterior pairing and also for the factorial polynomial
Fock pairing. Since $\pi$ intertwines both roots, its rational adjoint
$\pi^*$ does also, and $A=\pi^*\pi$ commutes with every root and Cartan
operator. The one-dimensional weight line gives
$At_\lambda=c t_\lambda$, with
\[
 c=\langle p_\lambda,p_\lambda\rangle_{\rm Fock}>0,
\]
because $t_\lambda$ has norm one and $p_\lambda\ne0$.
Commuting with every divided root word shows that $A=c$ on
$T_\lambda\otimes\mathbb Q$. Thus $\pi$ is injective there, and it
multiplies the pairing by $c$. Its integral image is exactly
$\Delta_\lambda$ by the defining row-word generators. This proves the
integral isomorphism and the normalized isometry.

Exactly $d$ column factors have height $i$. The root $F_i$ can change
$e_i$ to $e_{i+1}$ only in these factors; it is zero on the initial
wedges of all other heights. Consequently $q_s$ is the sum, with
coefficient one, over the $s$-subsets of these $d$ columns. Its summands
are distinct orthonormal basis vectors. This proves the norm formula.
Counting the predecessors or successors of a fixed subset proves the
two binomial operator formulas. The sums for different $s$ have
disjoint supports and each has a coefficient one, so their span is a
primitive sublattice of the ambient tensor lattice.

We also check that this is the complete relevant weight-line lattice of
$T_\lambda$. Over $\mathbb Q$, root words can be reordered with negative
roots first, Cartan operators next, and positive roots last: exchanging
two factors uses their commutator, reducing the word length, so induction
performs the reordering. Positive roots kill $t_\lambda$. To reach the
weight $\lambda-s\alpha_i$, every remaining negative root must be
$\alpha_i$, because each positive root is a nonnegative sum of simple
roots. Hence that rational weight space is the one-dimensional span
of $q_s$. Primitivity in the ambient tensor lattice gives its integral
intersection $\mathbb Z q_s$. Different weights are orthogonal. The
normalized dual lattice therefore has intersection
$\mathbb Z q_s/\binom ds$ with this rational line. This is the Weyl/dual
Weyl pairing used in \texttt{StatementE.lean}, declarations
\texttt{NablaEqDualWeyl} and \texttt{WeylEqDualNabla}, and in
\texttt{DataCoeff.lean}, declaration \texttt{FactDelDualNabla}.

Here is the promised exterior construction. Embed
$T^{\rm amb}_\lambda$ into
$\bigwedge^*(\bigoplus_{j=1}^q U_j)$ by putting the $j$th wedge in the
$j$th summand. In each column order the root coordinates $e_i,e_{i+1}$
first. On their four exterior states apply the integral basis
permutation
\[
 e_i\mapsto1,\qquad e_{i+1}\mapsto e_i\wedge e_{i+1},\qquad
 1\mapsto e_i,\qquad e_i\wedge e_{i+1}\mapsto e_{i+1}.
\]
It carries the row operator $e_{i+1}\wedge\iota_{e_i^\vee}$ to wedge
multiplication by the complete root pair, and the opposite root to its
contraction. Thus the sum over columns becomes a Lefschetz $L$, and the
opposite sum becomes $D$. The remaining $(r-2)q$ exterior coordinates
are spectators. For a nondegenerate symplectic ambient space, replace
each spectator's two states by the two singleton states of a fresh
symplectic pair; both are killed by its $L,D$. This is an integral
isometric embedding into $V_{(r-1)q}$, and the full symplectic $L,D$
restrict to the desired root operators. The displayed $q_s$ still have
disjoint supports with coefficients one. Therefore they are primitive
in this actual exterior lattice. If $\binom ds>1$, the vector
$q_s/\binom ds$ is not integral there.
\end{proof}

In particular factorial Fock positivity by itself does not choose the
Weyl lattice instead of its dual. For one binary column of degree $d$,
the polynomial lattice has monomial basis $x^{d-s}y^s$ whereas the
row-divided orbit of $x^d$ has basis
$\binom ds x^{d-s}y^s$. Both inherit the same positive factorial
pairing. IF-04 identifies the latter with orbit sums on the exterior
side; the former becomes its normalized dual. Thus the sentence in
\path{INTAKE_HODGE_SIGN_20261121.md}, section 5, identifying a
factorial pairing alone with the binomial lattice needs this refinement.

The cyclic lattice is weight-graded over $\mathbb Z$: every divided-root
word on the highest tensor is weight-homogeneous. Thus orthogonality
computes the intersection with the full dual lattice, not just a dual
inside an isolated rational string. The source also permits
\texttt{rowDiv a a k}; on a homogeneous row weight $w_a$ this is
multiplication by $\binom{w_a}{k}$. These diagonal coefficient
operators therefore add nothing to the cyclic integral lattice.

\subsection{IF-05: the auxiliary row-factor test at two lifts}
The actual column-kernel calculation, with determinant shifts and
finite theta jets, is IF-08--IF-09 below. This subsection tests the
distinct auxiliary compact row factors used in the same source.
We use the actual numerical datum \texttt{numEx}, not an invented
rank-two theta tower. Its declarations are
\path{LocalToGlobal/Divisibility/NumbersLemmas.lean}, namespace
\texttt{Formalism.L2.DIV}, declaration \texttt{numEx};
\path{LocalRingCardinality/SpanPowerMembership/ExampleRanksRankComplement.lean},
\texttt{R0}; and
\path{ExampleWeightMiddleEntries.lean}, \texttt{hcV,hcB,hcE}.
The data are $m=2,k=8,t=15,nV=24$, $e=(100,118)$, $b=(91,109)$,
$v=(24,20,16,12,8,4,0)$, and the definite tower ranks are $8,24,40$.
The auxiliary signatures \texttt{sigW} are $(8,1)$ and $(9,0)$.
The row functions \texttt{aPlus,aMinus} are defined in
\texttt{Numbers.lean}; their actual values are
\[
\begin{array}{c|l|l}
 j&a_j^+&a_j^-\\\hline
 1&(109,106,103,100,97,94,91,1)&(1)\\
 2&(126,123,120,117,114,111,108,18,1)&()
\end{array}.
\]
The adjacent positive-row differences are
$(3,3,3,3,3,3,90)$ and $(3,3,3,3,3,3,90,17)$.
Constant determinant shifts do not change them.

\paragraph{Statement IF-05.}
Apply IF-04 to the polynomial Weyl factors prescribed by these row
weights. On an actual simple-root string of length $d+1$, two successive
ordinary lowering operators have Smith invariant $2$ from the top line
to the second descendant in the Weyl lattice, and $d(d-1)$ in its
normalized dual. The second divided operator has invariants $1$ and
$\binom d2$, respectively. In particular the actual available values are
\[
\begin{array}{c|c|c|c}
 d&\text{two ordinary steps: Weyl, dual}
   &\text{second divided step: Weyl, dual}&\binom d2\\\hline
 3&(2,6)&(1,3)&3\\
 90&(2,8010)&(1,4005)&3^2\cdot5\cdot89\\
 17&(2,272)&(1,136)&2^3\cdot17
\end{array}.                                                   \tag{IF-05a}
\]
The isometry of IF-04 intertwines the Weyl operators integrally with
Lefschetz operators on its explicitly constructed exterior sublattice.
It cannot simultaneously carry the normalized dual lattice into that
same exterior lattice. More generally, for $d>1$ the normalized dual
root string admits no isometric integral injection into any lattice
having an orthonormal integral basis: it contains a nonzero vector of
squared norm $1/d<1$.

\begin{proof}
In the Weyl basis $q_s$ of IF-04, the first two ordinary arrows have
coefficients $1,2$; their product is $2$. In the normalized dual basis
$r_s=q_s/\binom ds$, the first two coefficients are $d,d-1$. Division
of the composite by $2!$ gives the stated divided invariants. These are
maps between free rank-one lattices, so the absolute coefficient is
exactly the Smith invariant. Substitution of the three adjacent gaps
gives the table. The norm calculation gives
$\langle r_1,r_1\rangle=1/d$. A nonzero integral vector in an
orthonormal integral lattice has squared norm a positive integer. This
proves the obstruction. It also shows concretely why replacing a Weyl
lattice by its Fock dual changes an index despite preserving the
rational representation and the root operator.
\end{proof}

The fixed row-weight test therefore has both a positive and a negative
outcome: there is an explicitly constructed integral normalized Weyl
intertwiner, and there are explicitly computed nontrivial denominators
for the simultaneously required dual. The normalizations have not been
suppressed. This advances HODGE\_SIGN U-04 from unspecified Fock data to
an actual lattice realization and an exact denominator calculation.

For comparison with the arithmetic entries themselves, HODGE\_SIGN
SL-06's exact substitution in \texttt{Numbers.lean}, declarations
\path{Dplus,Dminus,FTypeG,FRamD,FInert,FjP}, gives
\[
 (D_1^+,D_2^+)=(5608,6704),\quad(D_1^-,D_2^-)=(8,0),
\]
\[
 (F_{G,1},F_{G,2})=(-2808,-3352),\quad
 (F_{{\rm ram},1},F_{{\rm ram},2})=(-1400,-1652),\quad
 F_{{\rm inert},j}(p)=\frac{56e_j}{p^2-1}.
\]
These particular substitutions are rechecked here; the orbit optimization
underlying the last formula is cited as SL-06, not proved again.
At $p=3$ the two arithmetic contributions are $700,826$, while at
$p=89$ both \texttt{FjP} values are zero: $89$ exceeds the datum's cutoff
$X''=8\exp(2^{1/300})<22$ and is neither a distinguished prime nor the
ramified prime $7$. Nevertheless the fixed Weyl/dual comparison for its
$d=90$ row string has a factor $89$. Thus an asserted equality of the
raw step-10 $F$ contribution with this Fock lattice index is false on
the actual $m=2$ data. The full step-10 theorem, which additionally
contains gains, errors and stabilizer denominators, is not refuted.
Here $89$ differs from $3697,7393$ (the distinguished primes),
$7$ (the ramified $d$ prime), and also $2,11$; the above-cutoff
lemma only requires the first three exclusions. The exact statement is
\texttt{NumbersLemmas.lean}, \texttt{FjP\_above}, lines 191--194.
Moreover \texttt{numEx} itself is not a full \texttt{NumFacts}
witness. The change to the valid \texttt{numBig} instance and its
column weights is recorded explicitly in IF-09; its larger cutoff
is not being substituted into this $89$ comparison.
A reproducible rational-arithmetic certificate is
\path{step10_fock_test.py} and its tables.

\paragraph{Full-module rank check.}
The positive compact factors determined by the two row shapes, after
any constant determinant twist, have Weyl dimensions
\[
 105068200743744433029120,
 \qquad17540857265584852885502309695488000,
\]
respectively. Both have odd prime divisors; for example the first has
factor $3$ and the second has factor $3^4$. Consequently neither factor
can be isomorphic to a full $V_n$, whose rank is $4^n$. These numbers
come from the Weyl dimension product
$\prod_{a<b}(\lambda_a-\lambda_b+b-a)/(b-a)$ (Etingof et al.,
\emph{Introduction to representation theory}, Theorem 4.63,
\url{https://math.mit.edu/~etingof/replect.pdf}), evaluated in exact
rational arithmetic in the same certificate. This rank test concerns
these finite compact factors. It does not assert that the rank of a
global automorphic lattice is the rank of this one factor. IF-04 gives
a suitable sublattice realization rather than a full-exterior
identification.

\subsection{IF-06: the full oscillator cannot be the finite exterior module}
The preceding restricted construction does not justify replacing the
whole oscillator representation by its multiaffine part. The actual
source makes the distinction visible.

In \path{LocalRingCardinality/DataTheta.lean}, the definitions
\path{lowB,raiB,omE,fockAct1} and theorem
\texttt{fockAct1\_Eab} imply the following. At a signature $(p,q)$ with
$p,q>0$, against a positive definite partner of rank $R>0$, choose one
positive and one negative row. The noncompact raising root on the
polynomial Fock space is multiplication by
\[
 Q=\sum_{c=1}^R x_c y_c,
\]
and the opposite root is $-C$, where
$C=\sum_c\partial_{x_c}\partial_{y_c}$. On these variables
\[
 [C,Q]=H=R+\operatorname{Euler},\qquad [H,Q]=2Q,\qquad[H,C]=-2C.
\]
This applies in particular to the actual auxiliary signature
$\texttt{sigW}(1)=(8,1)$ of the $m=2$ datum.

\paragraph{Statement IF-06.}
Over any characteristic-zero field, there is no nonzero intertwiner
from a finite-dimensional Lefschetz module to this polynomial Fock
space taking its raising operator to $Q$. If both roots and their
commutator are intertwined, there is also no nonzero intertwiner in
the opposite direction, even after replacing contraction by its
negative to match the $\mathfrak{sl}_2$ signs.

\begin{proof}
Multiplication by the nonzero polynomial $Q$ is injective, and every
power of it is injective. On a finite Lefschetz module $L^{n+1}=0$.
Thus an intertwiner $T$ into the polynomial space satisfies
$Q^{n+1}T=TL^{n+1}=0$ and is zero.

For the converse, first decompose the polynomial space into strings
$Q^a h$ with $Ch=0$. Here is a direct existence proof of that familiar
harmonic decomposition. If $h$ is homogeneous of degree $d$ and
$Ch=0$, induction from the commutators gives
\[
 C Q^a h=a(R+d+a-1)Q^{a-1}h.                       \tag{IF-06a}
\]
For a homogeneous polynomial of degree $j$, apply induction to its
image under $C$, which has degree $j-2$. Each string term $Q^a h$
in that image has a preimage
$Q^{a+1}h/((a+1)(R+\deg h+a))$; the denominator is a nonzero positive
integer. Subtracting these preimages leaves a harmonic polynomial.
This proves that the stated strings span, over the characteristic-zero
field. Other polynomial variables may be treated as fixed spectator
coefficients.

Let $T$ go from the oscillator to a finite Lefschetz module, with
$TQ=LT$ and $T(-C)=DT$. For a harmonic homogeneous $h$, the vector
$w=Th$ is killed by $D$ and has eigenvalue $-(R+\deg h)$ for
$H_{\rm Lef}=[D,L]$. In a finite Lefschetz module a nonzero vector
$w$ killed by $D$ cannot have negative $H_{\rm Lef}$ weight. To see
this without a decomposition theorem, let $L^a w\ne0$ be its last
nonzero translate and write $H_{\rm Lef}w=\delta w$. The commutator
formula gives
$0=DL^{a+1}w=(a+1)(\delta-a)L^a w$, so $\delta=a\geq0$.
Here the required weight is $-(R+\deg h)<0$, a contradiction.
Therefore $Th=0$, and $TQ^a h=L^aTh=0$ on every spanning string.
\end{proof}

This rules out the full-oscillator shortcut. It does not rule out
finite compact pieces or the explicitly constructed integral
intertwiner IF-04. In particular it is not a nonexistence theorem for
every possible geometric realization of the global definite tower
$8\to24\to40$.

\subsection{IF-07: what an exact period valuation can depend on}
For fixed integral linear maps $K_0,K_1$ and an integral terminal
functional $\ell$, put $P(v)=\ell(K_1K_0v)$. If $P(v)\ne0$, then
for every rational prime $p$ and $s\geq0$,
\[
 v_p(P(p^s v))=s+v_p(P(v)).
\]
This follows immediately from linearity, and the same identity holds
for linear theta lifts and their terminal periods when the allowed
input vector is scaled. Thus lift count alone does not determine a
nonzero period's exact valuation. With a fixed input vector this
argument says nothing about the valuation; it is not used to dismiss
that stronger problem.

There is already a fully constructed nonzero local comparison here:
for the actual root strings of IF-05 take $v=q_0$ and pair the result
with its normalized Fock dual $r_2$. The ordinary two-step period is
$\langle r_2,F_i^2q_0\rangle=2$, and the divided two-step period is
$\langle r_2,F_i^{(2)}q_0\rangle=1$. All their valuations are exact.
These are periods of the specified row-root maps, not evaluations of
the global theta integral. Identifying the latter with these maps
would require a further equality of actual kernels and their
stabilizer normalizations. \texttt{Inputs.lean}, structure
\texttt{ColChain}, supplies the rectangular global matrices and a
factorization, but no such equality is an axiom of that structure.

\paragraph{Algebraic-geometric placement and precise remaining assertion.}
IF-01--IF-05 live over $\operatorname{Spec}\mathbb Z$ and commute with
base change. For a principally polarized complex abelian variety,
$H^*(A,\mathbb Z)=\bigwedge^*H^1(A,\mathbb Z)$ realizes their exterior
operator calculations, with $L$ cup product by the polarization.
The same matrices act on prime-to-characteristic integral
$\ell$-adic cohomology of a principally polarized abelian scheme.
Passing from degree $i$ to $i+2$ uses a Tate twist; after the twists in
HODGE\_SIGN SL-07, the correspondence commutes with Frobenius. The
calculation does not make an integral inverse Lefschetz correspondence
algebraic, and it does not identify a Weyl lattice's Fock dual with an
integral cycle lattice. In $K$-theory the constructions use exterior
powers and tensor products; in the category of rational motives they
are operations on the appropriate realizations of $h^1(A)$, together
with the Lefschetz correspondence. No new identity of motivic
$L$-functions or of global theta periods follows from the lattice
isomorphisms alone.

The finite-kernel part of the predecessor's remaining assertion is
supplied next in IF-08--IF-09: each actual finite theta jet is expressed
by an explicit exterior-tensor matrix, including the dual lattices,
determinant shifts and stabilizer factors. What that does not assert
is that these varying tensor matrices are one common Lefschetz
$\mathfrak{sl}_2$ operator, or that an arbitrary global automorphic
period is the particular root period just computed. Those are distinct
operator and geometric assertions, not missing coefficient-model data.

% IF-08 and IF-09: actual coefficient kernels and actual m=2 columns.
% Included by contribution.tex. Full author proof reading recorded separately.
\subsection{IF-08: the actual finite theta jet has an exterior-tensor realization}
The finite coefficient map in the source is more general than a single
raising operator. It can nevertheless be realized explicitly and
integrally in exterior tensors, with exactly the same Smith indices.
Here the relevant source declarations are
\path{LocalToGlobal/KernelLattices/ThDistC.lean},
\path{subD,jetLin,jetMap,jetOp,JetInVV,derivAt0,jetC}
(lines 23--61), and
\path{LocalRingCardinality/DataKerI.lean},
\path{KM.Gxz,KM.applyAt0,KerEnt.kform} (lines 27--35 and 115--121).
The former is the coefficient map itself; the latter describes its
matrix entries tested against the dual product-minor lattice.

Fix a partition $\lambda$, with column heights $h_1,\ldots,h_q$,
and integers $a,b\geq\ell(\lambda)$. Put
\[
 A_\lambda(r)=\bigotimes_{c=1}^q\Lambda^{h_c}\mathbb Z^r,
 \qquad T_\lambda(r)\subset A_\lambda(r)
 \quad (r=a,b),
\]
where $T_\lambda(r)$ is the divided-root cyclic lattice of IF-04.
Write $\pi_r:A_\lambda(r)\longrightarrow\mathbb Z[X]$ for
multiplication of the corresponding initial-column minors. Use polynomial
matrices with $r$ rows and the common $\ell(\lambda)$ columns. This
makes rectangular substitution meaningful also when $b<a$; the
source \texttt{subD} sends unused source columns $s\geq b$ to zero,
and no such column occurs in these minors. Its cyclic
restriction is the isomorphism
$\Phi_r:T_\lambda(r)\xrightarrow\sim\Delta_\lambda(r)$ of IF-04.
The full tensor image is $\nabla_\lambda(r)$, by the supplied
\path{GradedWeylDivisibility/FinalC.lean}, theorem
\path{standardMonomialSpan}, whose exact statement is
\texttt{StatementA.lean}, \path{StandardMonomialSpan}, lines 221--226.
We use that supplied integral straightening result as background.

Let $Z=(z_{ui})$ be a $b$ by $a$ matrix of independent variables and set
\[
 A_\lambda(Z)=\bigotimes_c\Lambda^{h_c}Z
       =\sum_\beta Z^\beta A_\beta,
 \qquad A_\beta\in\operatorname{Hom}_{\mathbb Z}
       (A_\lambda(a),A_\lambda(b)).
\]
For a characteristic-zero coefficient field $K$ and a prescribed moment
function $J:\mathbb N^{ba}\to K$, define the concrete finite matrix
\[
 B_J=\sum_{|\beta|=|\lambda|}J(\beta)A_\beta.       \tag{IF-08a}
\]
Only finitely many entries of $J$ are used. If
$\Theta(Z)=\sum_\beta\theta_\beta Z^\beta$, the source convention is
\[
 J_\Theta(\beta)=\beta!\,\theta_\beta,
 \qquad \beta!=\prod_{u,i}\beta_{ui}!,              \tag{IF-08b}
\]
including the explicitly prescribed powers of the discriminant-period
generators in \texttt{jetC}. Thus $J$ consists of ordinary derivatives
at zero; it is not an unspecified divided-power normalization.

\paragraph{Statement IF-08.}
The lattices $T_\lambda(r)$ are saturated in $A_\lambda(r)$, and
\[
 \Phi_r:T_\lambda(r)\simeq\Delta_\lambda(r),\qquad
 \Phi_r:T_\lambda(r)^\#\simeq\nabla_\lambda(r)       \tag{IF-08c}
\]
are isometries for the tensor-exterior pairing on the left and the
normalized Fock pairing on the right. The matrix $B_J$ preserves the
rational cyclic spaces. If its moments lie in a subring $R\subset K$,
it carries $T_\lambda(a)\otimes R$ into $T_\lambda(b)\otimes R$.
On those spaces the actual polynomial substitution followed by moment
contraction satisfies
\[
 \Phi_b B_J=\operatorname{jet}_J\Phi_a.             \tag{IF-08d}
\]
This equation is an equality of explicit coefficient matrices. With
source and target lattices as displayed, it preserves every
Smith invariant over $\mathbb Z$ and every elementary-divisor exponent
over a discrete valuation ring. The transpose identity gives the same
statement for the normalized dual lattices.

For a determinant insertion of order $n\geq0$, multiplication by
$\det(X)^n$ is represented on the exterior side by insertion of $n$
copies of the oriented top wedge of $\mathbb Z^a$. Thus the complete
\texttt{jetMap J n}, including its determinant factor, is intertwined
by the same construction. Its rational retraction \texttt{retrVV}
has no effect under the supplied \texttt{JetInVV} hypothesis.

\begin{proof}
We first make the scalar normalization and saturation explicit. Let
\[
 c_\lambda=\langle p_\lambda,p_\lambda\rangle_{\rm Fock}
   =\sum_\alpha ([X^\alpha]p_\lambda)^2\alpha!>0.
                                                               \tag{IF-08e}
\]
This explicit finite sum is an integer independent of the ambient rank $r\geq\ell(\lambda)$.
IF-04 proves $\pi_r^*\pi_r=c_\lambda$ on $T_\lambda(r)_\mathbb Q$,
where the adjoint uses the unnormalized factorial Fock form and the
orthonormal exterior tensor form. The rational image of $\pi_r$ is
already the image of its cyclic subspace, by
\path{GradedWeylDivisibility/FinalB.lean}, theorem
\texttt{weylEqNablaQ} (the supplied \texttt{WeylEqNablaQ} statement).
Hence $T_\lambda(r)_\mathbb Q=(\ker\pi_r)^\perp$. On the whole tensor
space we therefore have
\[
 \pi_r^*\pi_r=c_\lambda P_r,\qquad
 P_r=\frac1{c_\lambda}\pi_r^*\pi_r,                 \tag{IF-08f}
\]
where $P_r$ is the orthogonal projection onto the cyclic rational
space. These formulas display, rather than discard, the scalar.

Suppose $t\in T_\lambda(r)_\mathbb Q\cap A_\lambda(r)$.
For every $v\in A_\lambda(r)$, (IF-08f) gives
\[
 \frac{\langle\pi_rt,\pi_rv\rangle_{\rm Fock}}{c_\lambda}
       =\langle t,v\rangle\in\mathbb Z.
\]
The full image $\pi_r(A_\lambda(r))$ is $\nabla_\lambda(r)$.
The supplied normalized duality theorem
\path{LocalToGlobal/Columns/CoefficientsModP.lean},
\texttt{weylEqDualNabla} (line 31),
therefore implies $\pi_rt\in\Delta_\lambda(r)$; IF-04 gives
$t\in T_\lambda(r)$. This proves saturation. The other normalized
duality, \path{nabla_eq_dual_weyl_holds} in
\path{Columns/CoefficientsModP/WeTransposeWordTransposeWord.lean}
(line 74), and the isometry of IF-04 yield
(IF-08c). Equivalently $P_r(A_\lambda(r))=T_\lambda(r)^\#$;
this also follows by extending an integral linear functional on the
primitive sublattice $T_\lambda(r)$ to the free ambient lattice.
The projection can have denominators, exactly those already present
in this specified dual lattice.

The exterior matrix $A_\lambda(Z)$ has entries
$\prod_c\det Z[I_c,J_c]$. Multiplication of minors gives
\[
 \pi_b A_\lambda(Z)=\rho_Z\pi_a,                  \tag{IF-08g}
\]
where $\rho_Z$ substitutes $X_{is}\mapsto\sum_u Z_{ui}X_{us}$.
For completeness, expand a determinant of these sums by multilinearity
in its rows. Terms with a repeated index vanish by alternation; grouping
all permutations of each increasing index set gives the determinant
$\det Z[I,J]$ times the corresponding minor of $X$. Multiplying these
identities over the columns gives (IF-08g), including all signs.

For a rational specialization of $Z$, (IF-08g) sends $\ker\pi_a$
into $\ker\pi_b$. Its transpose version sends $\ker\pi_b$ into
$\ker\pi_a$, since $A_\lambda(Z)^t=A_\lambda(Z^t)$.
Taking orthogonal complements shows
$A_\lambda(Z)T_\lambda(a)_\mathbb Q\subset
T_\lambda(b)_\mathbb Q$. This holds for every rational matrix $Z$.
Every coefficient $A_\beta$ consequently has the same containment:
apply any linear functional vanishing on the latter space and use
that a polynomial vanishing at all rational tuples is zero.
The coefficients $A_\beta$ are integer matrices; the saturation just
proved implies $A_\beta T_\lambda(a)\subset T_\lambda(b)$.
This proves the claimed integrality for arbitrary $R$ and $J$.

Apply the moment functional $Z^\beta\mapsto J(\beta)$ to (IF-08g).
It gives (IF-08d). For $J=J_\Theta$, direct differentiation of a
monomial proves (IF-08b); this is exactly the factorial in
\texttt{KM.applyAt0}. The test-minor version is
\[
 \operatorname{applyAt0}(\Theta,G_I(XZ))
   =\sum_J \Bigl(\sum_\beta\beta!\theta_\beta
                [Z^\beta]\prod_c\det Z[J_c,I_c]\Bigr)G_J(X),
                                                               \tag{IF-08h}
\]
with the direction of $Z$ transposed to match the source's test
convention. This is the same matrix, paired with the dual lattice.
Under (IF-08c), normalized Fock testing is ordinary tensor-exterior
testing. The $c_\lambda$ in (IF-08f) has canceled because the Fock
form is divided by that same $c_\lambda$; replacing normalized Fock
by unnormalized Fock would retain that factor.

Finally insert $n$ full-height columns in $A_\lambda(a)$. Their top
wedges have norm one and transform by $\det^n$, while every row root
annihilates them. Their cyclic sublattice is therefore exactly the
cyclic lattice of the shape $\lambda+n(1^a)$, and their minor product
is $\det(X)^n$. Thus determinant insertion also has no unspecified
scalar. In integral bases, (IF-08d) and its determinant variant are
conjugations by unimodular source and target matrices. Their images,
cokernels, saturated images and local lengths are correspondingly
isomorphic. This proves the assertions about Smith invariants,
including for rectangular maps and rank-deficient images, where the
finite index means the index in the saturation of the image.
\end{proof}

\paragraph{Assembly of the actual finite kernel.}
This is not merely a highest-weight-root comparison. For every
supplied theta jet $J_{j,c',c}$, (IF-08a) specifies its exterior-tensor
matrix. The complete block matrix retains the source's actual scalar
factors:
\[
 \widehat K_j(c',c)=\frac{\rho_j(c)\rho'_j(c')}{|\Gamma_{j,c}|}
        B_{J_{j,c',c}}\,D_{n_j}.                  \tag{IF-08i}
\]
Here $D_{n_j}$ inserts the determinant factors, and the target
one-dimensional determinant character is retained as a separate
rank-one lattice. The scalar entries and the jets include precisely
the period generators and frame changes specified in the source.
Formula (IF-08i) is the tensor version of
\path{SharpRestriction/GlobalSchurBounds/GL3Theta.lean},
\path{liftTie,lift_formula} (lines 575--596), and
\path{LocalRingCardinality/DataKernels.lean},
\path{wt,sharp,divDe,Kmap,K} (lines 52--68).

The coefficient isomorphisms commute with the relevant group actions;
hence they transport invariant submodules, local lattice tests and
finite direct sums over lattice classes. For an actual supplied
finite global matrix they give an isomorphic kernel matrix, with all
its stabilizer denominators still displayed. The test condition
\texttt{KerEnt.lat\_S} is pairing integrality against
\texttt{KM.nablaX}; its exterior version is integrality against
$T_\lambda^\#$, not against a second copy of $T_\lambda$.
When the source's local integrality assumptions hold, the transported
matrix has exactly the same integral elementary divisors. No new
integrality assumption is silently supplied here.

The precise global integral lattice is also identifiable in the source.
\path{Columns/Record.lean}, \texttt{stepH} (lines 1286--1287),
uses \texttt{intDel}; \path{Columns/CoefCol.lean},
\texttt{roundPk} (lines 59--64), uses \texttt{Del=delAtr} and
\texttt{Nab=nabAtr}. Thus the integral global column forms use the
Weyl lattice; the dual Weyl lattice supplies its testing lattice.
In \path{Divisibility/CoefMod.lean},
\texttt{transLat,delP,delK} (lines 196--211), the lattice is explicitly
the intersection of the inverse images of standard local lattices
under \texttt{actP} applied to $B^{-1}$ for local frames $B$.
Since \texttt{actP} is determinant twist times the transposed
polynomial row action (lines 69--76), its exterior action is
\[
 R_\lambda(B)=\prod_\tau\det(B_\tau)^{s_\tau}
       \bigotimes_{\tau,c}\Lambda^{h_{\tau,c}}B_\tau.
                                                               \tag{IF-08j}
\]
The round realization uses the source's inverse-transpose convention
\texttt{rhoVr}; its \texttt{delAtr} is formed from the dual frame
lattice (\texttt{CoefMod.lean}, lines 320--330). Applying exactly that
same convention to (IF-08j) transports the global intersection and
all its localizations; no free global basis of a nonprincipal
coefficient ideal is presumed.

In the following formula $B_{j,c}$ means the frame of the coefficient
lattice used by \texttt{transLat}. For an actual round column whose
geometric frame is $C_{j,c}$, this means $B_{j,c}=C_{j,c}^{-t}$,
with the dual-frame convention just specified; one must not insert
$C_{j,c}$ directly into $R_j(B_{j,c})$. For chosen local bases at a
prime $w$, the transported block matrix is therefore the explicit matrix
\[
 R_{j+1}(B_{j+1,c'})^{-1}
 \frac{\rho_j(c)\rho'_j(c')}{|\Gamma_{j,c}|}
 \Bigl(\sum_{\beta}J_{j,c',c}(\beta)A_{\beta}\Bigr)
 D_{n_j} R_j(B_{j,c}).                              \tag{IF-08k}
\]
One sums these blocks over the actual finite class representatives,
then restricts to the invariant and level submodules defining forms.
This is the double-coset sum \texttt{kerOpM}; its scalar analogue is
\texttt{thetaLiftAMF} in \path{ThetaLiftIndependence/Statement.lean},
lines 80--84. In a local field the scalar valuation in (IF-08k) is
\[
 v_w(\rho_j(c))+v_w(\rho'_j(c'))-v_w(|\Gamma_{j,c}|),
\]
and a jet additionally contributes the recorded period-generator
valuation $\sum_\tau |\beta_\tau|v_w(c_{D,\tau})$ to its ordinary
Taylor derivative valuation. They have not been dropped or replaced
by root factorials. For a resulting local matrix $M$ of rank $r$,
put $s_0=0$ and
\[
 s_k=\min_{|I|=|J|=k}v_w(\det M_{I,J}),\qquad 1\leq k\leq r.
\]
Its elementary-divisor exponents are $s_k-s_{k-1}$; its image index
in the saturated image, when $M$ is integral, has length $s_r$.
These statements follow directly from the Smith normal form over a
DVR, and the formulas are unchanged by (IF-08c). This gives an exact
finite algorithm for the actual global transport, including every
local denominator. A terminal period functional is transported by
the same dual isomorphism, so its value on the transported chain is
identical, not merely bounded by the original value.

The result does not identify (IF-08i) with a single $L$, $L^{(r)}$,
or a polynomial in one Lefschetz pair. The actual moments,
class numbers, characters and stabilizers remain part of the matrix.
It constructs the requested finite-kernel intertwiner and proves
index transport; equality with a separately defined Lefschetz order
index would still require an additional operator identity. IF-03 and
IF-05 show why discarding the operator and lattice distinctions gives
false numerical equalities.

\subsection{IF-09: the actual two-step column chain and a fixed kernel test}
The polynomial coefficient weights of the definite column chain must
be distinguished from the auxiliary $a^+$ row weights in IF-05.
For the actual \texttt{numEx} datum and its all-false word, define
\[
 u=(113/2,47/2,39/2,31/2,23/2,15/2,7/2,-1/2),\qquad
 \lambda_0=(53,21,18,15,12,9,6,3).
\]
The unshifted weights from \path{MemberForms.lam_eq} are
\[
 \lambda_1^N=(\lambda_0,-88^{[16]}),\qquad
 \lambda_2^N=(\lambda_0,-88^{[16]},-90^{[16]}).
\]
The actual \texttt{colVal} uses \texttt{Wts.lamTrue}, which subtracts
$8j$ from every entry at stage $j$. Thus the determinant exponents of
its three polynomial realizations are
\[
 (s_0,s_1,s_2)=(3,-96,-106),
\]
and their polynomial partitions are
\[
\begin{split}
 \alpha_0&=(50,18,15,12,9,6,3,0),\\
 \alpha_1&=(141,109,106,103,100,97,94,91,0^{[16]}),\\
 \alpha_2&=(143,111,108,105,102,99,96,93,2^{[16]},0^{[16]}).
\end{split}                                                     \tag{IF-09a}
\]
Here $x^{[r]}$ means $r$ repeated entries. These formulas follow by
substitution in \path{DataCoeff.hcToWeightQ},
\texttt{Params.wG}, \path{MemberForms.lam_eq}, and
\path{Columns/AbelianHostFieldConditions.lean},
\path{Wts.lamN, Wts.lamTrue} (lines 462--465).
\path{Columns/CoefCol.lean}, \texttt{colVal} (lines 89--91), and
\path{Divisibility/ColPins.lean}, \texttt{lamOfCol} (lines 21--26),
pin this distinction to the actual column map.

The kernel diagrams, before appending target zeros, are
\[
 \nu_1=(141,109,106,103,100,97,94,91),\qquad
 \nu_2=(143,111,108,105,102,99,96,93,2^{[16]}).       \tag{IF-09b}
\]
They are exactly the output of \texttt{NumDatum.nu} and
\texttt{DV.nuRow}; they are also
$\alpha_0+91(1^8)$ and $\alpha_1+2(1^{24})$.
The source has
$\mathrm{sSharp}=(88,98)$ and $\mathrm{sDiv}=(96,106)$.
The identities
$s_i+\mathrm{sSharp}_i=n_i$ and
$s_{i+1}=-\mathrm{sDiv}_i$ give
\[
 (n_0,n_1)=(91,2).                                  \tag{IF-09c}
\]
They are precisely \texttt{KerEntTied.WtTie} in
\path{KernelLattices/JetFieldLeKR.lean}, lines 108--111.
Consequently the actual finite coefficient chain of IF-08 is
\[
 T_{\alpha_0}(8)\xrightarrow{D_{91}}
 T_{\nu_1}(8)\xrightarrow{B_{J_1}}
 T_{\alpha_1}(24)\xrightarrow{D_2}
 T_{\nu_2}(24)\xrightarrow{B_{J_2}}
 T_{\alpha_2}(40),                                  \tag{IF-09d}
\]
The two determinant degree increments are $8\cdot91=728$ and
$24\cdot2=48$. Thus two lifts do not represent $L^2$.
The chain retains the determinant characters $s_i$ and the scalar block
factors (IF-08i). This specifies the whole finite polynomial kernel,
not just a simple-root summand.

The actual column simple-root gaps include $32,3,91,2$ (zero gaps are
trivial). Their top-to-second-descendant invariants are
\[
\begin{array}{c|c|c|c}
 d&F^2\text{ on Weyl}&F^2\text{ on normalized dual}
       &F^{(2)}\text{ on normalized dual}\\\hline
 2&2&2&1\\
 3&2&6&3\\
 32&2&992&496=2^4\cdot31\\
 91&2&8190&4095=3^2\cdot5\cdot7\cdot13
\end{array}.                                                   \tag{IF-09e}
\]
On the Weyl lattice every displayed $F^{(2)}$ invariant is one, by
IF-04. Thus the dual-lattice obstruction is present in the actual
column coefficients themselves. The auxiliary row gap $90$ in IF-05
is a separate, correctly labeled test.

The source's \texttt{numEx} is a numerical datum, not a full
\texttt{NumFacts} witness. The repaired \texttt{numBig} is given in
\texttt{NumbersLink.lean}, \texttt{numBig} (line 803) and
\texttt{numFactsBig} (line 1260). It retains the row parameters and
changes $u_7$ from $-1/2$ to $-7/2$ at one specified real place, besides
changing the size parameters. At that place $s_0=0$, $n_0=88$;
the last entries $91,93$ in (IF-09b) become $88,90$ respectively.
The other entries and the second determinant insertion $n_1=2$
are unchanged. The additional column gaps are $6,88$, with
$\binom62=15$ and $\binom{88}2=3828=2^2\cdot3\cdot11\cdot29$.
The normalized dual obstruction therefore persists for this valid
numerical-facts instance. It does not depend on treating
\texttt{numEx} as a nonexistent full global witness.

There is also a precise host obstruction to regarding
\texttt{numBig} itself as a full global \texttt{D7} record.
\texttt{NumbersLink.numBig} retains \texttt{numEx.H};
\path{NumbersLemmas.hostEx_q,hostEx_d} give $q=11,d=7$.
But \texttt{DataHost.lean}, \texttt{Host.qd\_big} (line 52), requires
$4tr+1<q,d$. Here $t=15,r=9$, so $4tr+1=541$, contradicting both
inequalities. Thus its \texttt{NumFacts} theorem certifies numerical
conditions only; it cannot instantiate that full global host with
these primes. The polynomial representations and kernel tests above
are nevertheless fully specified and valid on their stated coefficient
lattices.

\paragraph{An exact nonzero kernel test.}
For the rectangular coordinate inclusion $A:\mathbb Z^a\hookrightarrow
\mathbb Z^b$, define
\[
 \Theta_A(Z)=\exp\bigl(\sum_{u,i}A_{ui}z_{ui}\bigr),\qquad
 [Z^\beta]\Theta_A=A^\beta/\beta!.
\]
This is the precise series \texttt{ThE} in
\path{SpanPowerMembership/NZRowV.lean}, lines 77--79.
Its factorial cancellation is the supplied theorem
\texttt{fac\_ThE}, lines 89--97; its kernel evaluation is
\path{eval_applyAt0,eval_Gxz}, lines 99--146.
For this specified series $J(\beta)=A^\beta$, so
$B_J=A_\lambda(A)$ by (IF-08a). Each map in (IF-09d) is then an
inclusion with saturated image. Indeed it carries its ambient tensor
basis injectively into the larger ambient tensor basis; saturation of
the cyclic lattices was proved in IF-08. The determinant insertion
has the same property, since it tensors with a primitive norm-one
vector. Their composite is primitive for the same ambient basis
reason. Its nonzero Smith factors are all one.

Taking the highest vector at the beginning and the normalized dual
highest functional at the end gives period exactly $1$ after both
arrows. This assertion is fully fixed-data: the kernels, lattices,
vectors, determinant insertions and functional have been specified.
On the highest diagonal monomial the derivative factor is
$\prod_i\nu_{j,i}!$ and the coefficient of $\Theta_A$ is its inverse,
so the unit value is an explicit factorial cancellation, not an
assumption that the two normalizations agree.

\paragraph{A compatible common Lefschetz action on this fixed family.}
There is nevertheless a positive integral intertwining result for the
same specified kernel chain. Fix one simple root inside the first eight
rows, for example rows $2,3$. Every determinant insertion and coordinate
inclusion in (IF-09d) commutes with both of its root operators $E_i,F_i$:
$E_i$ and $F_i$ kill a full determinant vector, and the coordinate
inclusion is the identity on those two rows. These identities hold on
the entire cyclic coefficient lattice, not just the highest root string.

For each coefficient lattice apply the particle-hole embedding of IF-04
into $V_{N_j}$, with $N_j=(r_j-1)q_j$ and $q_j$ its number of columns.
Let $N=\max_jN_j$, and pad the embedding with a singleton $e$ in each
of the $N-N_j$ extra symplectic pairs. Singleton padding is an integral
isometry with primitive image and is annihilated by both the raising
and lowering on each extra pair. It therefore intertwines the same
$L,D$ on $V_N$. Primitivity of $T_\lambda$ in its column ambient lattice
(IF-08), followed by the coordinate-basis particle-hole embedding and
padding, makes each image a primitive $L,D$-stable sublattice of $V_N$.
Thus the fixed \texttt{ThE} arrows are specified integral isometric maps
between those sublattices, commuting with one common Lefschetz action,
and their nonzero Smith factors remain one. For (IF-09d), one may take
$N=39\cdot143=5577$ at each embedding factor. Tensoring the eight
factors gives the diagonal root action on $V_{8N}=V_{44616}$, with
the same integral and primitivity statements.

The exterior degree of a vector of simple-root weight $h$ is
$N_j-h$ before padding and $N-h$ after padding. Consequently these
intertwining maps have degree zero. They are not applications of $L$;
that stronger assertion is excluded below. No integral equivariant
extension of the arrows to all of $V_N$ is asserted. Nor is algebraicity
of these particular particle-hole image sublattices asserted; IF-10
provides a separate geometric coefficient realization with its stated
character normalization.

For an arbitrary moment block $B$ between the same standard coefficient
lattices, the exact additional common-root conditions are the finite
matrix equations
\[
 E_i' B=B E_i,\qquad F_i' B=B F_i.                 \tag{IF-09f}
\]
They suffice for all divided powers, by equality in characteristic zero
and their integral preservation. For the actual global form spaces one
must also check that the proposed root action preserves their invariant
and level submodules and their local lattices. These are explicit
operator/lattice conditions after (IF-08k), and are not consequences of
mere theta-kernel rationality. No such common-root invariance is assumed
for the unspecified global moments.

\paragraph{A fixed-normalization obstruction to two ordinary Lefschetz arrows.}
This particular two-step kernel chain cannot be intertwined with two
successive applications of the same fixed ordinary $L$ on one full
integral symplectic exterior lattice by integral embeddings whose
final target image is saturated. Indeed its highest vector is
sent to a primitive highest vector. A saturated integral embedding
preserves primitivity. But $L^2=2L^{(2)}$ on every such exterior
lattice, so the image of $L^2$ is contained in $2V_n$ and contains no
primitive vector. This is a contradiction. Even without saturation,
an isometric integral embedding for the fixed normalized pairings is
impossible: the resulting highest vector has squared norm one, whereas
a nonzero vector in $2V_n$ has squared norm at least four. The full
exterior monomial basis is orthonormal, so both assertions follow
immediately by inspecting its integer coefficients. This proves a
fixed-data failure of the ordinary-$L^2$ normalization, rather than a
failure caused by leaving the test vector unspecified. The same argument
also excludes two positive divided powers of that same fixed $L$:
for $a,b>0$,
\[
 L^{(a)}L^{(b)}=\binom{a+b}{a}L^{(a+b)},
\]
and the integer $\binom{a+b}{a}>1$ divides every coordinate of the
composite image. Thus a primitive or norm-one composite cannot occur.
This leaves tensor-functor realizations such as IF-09d, or operators
belonging to different specified Lefschetz pairs, outside the
obstruction.

These kernels are a concrete algebraic test family for the exact
source operator formula; the source also builds \texttt{KerEnt}
instances with \texttt{ThE} in \texttt{NZRowW.lean},
\path{levelRecordTwistEntries}, and \texttt{NZRowZ.lean},
\texttt{fin\_ent}. We do not assert that this exponential test family
is the globally automorphic theta function in Alp\"oge's construction,
or that it satisfies all the additional global theta hypotheses.
For that actual global function, (IF-08i) is still an explicit
intertwining formula in its actual finite jets; computing its Smith
numbers requires those finite jet values and class data. An
existential theta construction or a lower bound does not determine
those numbers. The finite-kernel realization has therefore been
constructed, while the distinct claim that every actual global index
equals one universal Lefschetz order index has not been proved.

\paragraph{Schemes, $\lambda$-operations, and realizations.}
Every $A_\lambda(Z)$ above is a polynomial matrix over
$\operatorname{Spec}\mathbb Z$, assembled from the exterior-power
operations $\Lambda^{h_c}$ and tensor products. Its class in $K_0$
is the corresponding product of $\lambda^{h_c}$ classes. The
saturated cyclic lattice and its explicit dual retain information
which $K_0$ alone forgets: a finite-index inclusion has the same
rational character and can have different integral cokernel.
Base change to $\mathbb Z_\ell$ preserves the displayed integral
intertwiners and carries a Smith factor $d$ to local length
$v_\ell(d)$. Tensoring with $\mathbb Q_\ell$ forgets that length.
When these spaces are realized in the cohomology of an abelian
variety, their tensor and exterior operations are motivic tensor
operations. The formal coefficient identification alone supplies no
new algebraic correspondence representing the theta moment matrix
or the Hodge sign operator. Such an assertion would require the
separate algebraicity comparison under discussion in Q6g and the
main question.


% IF-10: explicit geometric coefficient realization, with every normalization retained.
\subsection{IF-10: a geometric realization over a CM elliptic curve}
There is a geometric realization of the finite coefficient comparison,
but its integral and character normalizations must be stated. It is not
an identification of an automorphic motive with the cohomology of one
abelian variety, and does not replace the operators of IF-08 by one $L$.

Let $K\subset\mathbb C$ be a number field containing the entries of the
finite matrices under consideration. Enlarge it to a number field
$K'$ containing $i$, $K$, and their complex conjugates, and write
$R=\mathcal O_{K'}$. Denominators of the matrices remain allowed and are
recorded by fractional $R$-lattices, as in (IF-08k). Use the elliptic curve
\[
 E:y^2=x^3-x\quad\text{over }\mathbb Q(i),\qquad
 [i](x,y)=(-x,iy).
\]
We work with algebraic correspondences modulo homological equivalence,
with coefficients in $K'$, and their Betti or $\ell$-adic realizations.
All correspondences constructed below are explicit combinations of
coordinate maps, the origin, $[i]$, and linear combinations of their
graphs. No Hodge-conjecture input is needed for this coefficient model.

\paragraph{Statement IF-10.}
For every finite chain of the exterior coefficient spaces and operators
in IF-08 whose normalized matrix entries lie in $K$, there is a chain
of homological motives generated by powers of $E$ and rank-one CM
character twists with the following properties.

\begin{enumerate}
\item Its realization has exactly the matrices of IF-08, including the
moment coefficients, determinant insertions, class sums, stabilizers,
local frames, invariant subspaces and terminal functionals. The
coefficient lattices, with their specified duals, transport isomorphically.
Before coefficient extension every Smith invariant of the same specified
map is identical. After extension, determinantal ideals are their extended
ideals; local valuations change by the ramification index, as specified
below. No equality of absolute finite indices across a field extension
is asserted.
\item Before the character twists, the polynomial coefficient space of
partition $\lambda$ is a submotive of a power of $E$ whose realization is
$T_\lambda\otimes M^{|\lambda|}$, where $M$ is the $i$-eigenmotive in
$h^1(E)$. A determinant twist $s$ changes the exponent to
$w=|\lambda|+rs$. The underlying integral polynomial lattice is a
saturated intersection in integral cohomology after extending coefficients
to $R$.
\item The Fock isometry uses $2^{-|\lambda|}$ times the ambient
Hermitian form on that polynomial eigenspace. The CM projector has
denominator $2$, and the dual eigenline introduces the factor $-2i$.
At primes above $2$ the transported abstract dual lattice must therefore
be distinguished from the native saturated conjugate eigenlattice.
No assertion of an integral cycle projector, or of an integral cup
polarization with the normalized Fock Gram matrix, is made.
\end{enumerate}

\begin{proof}
Choose an oriented integral symplectic basis $e,f$ of $H^1(E,\mathbb Z)$
such that
\[
 [i]^*e=-f,\qquad [i]^*f=e,\qquad
 \psi(e,f)=1.
\]
This is the basis $dx,dy$ on the square analytic torus. Thus
$u=e+if$ satisfies $[i]^*u=iu$ and
\[
 U=R u=\ker([i]^*-i)\subset H^1(E,R).
                                                               \tag{IF-10a}
\]
This line is saturated: $u,f$ form an $R$-basis, since their coefficient
matrix in $e,f$ has determinant one. Set
$\pi_0=[\{0\}\times E]$, $\pi_2=[E\times\{0\}]$, and
$\pi_1=\Delta_E-\pi_0-\pi_2$, with the convention that the displayed
correspondences act by pullback followed by pushforward to the second
factor. They project to cohomological degrees $0,2,1$, respectively.
The transpose graph of $[i]$ acts by $[i]^*$. Hence
\[
 p_+=\frac{1}{2}\bigl(\pi_1-i\,[i]^*\pi_1\bigr)
                                                               \tag{IF-10b}
\]
is an algebraic idempotent in the homological category and its image is
$M$. Its realization has image $K'u$ and its intersection with integral
cohomology is precisely $U$. The idempotence follows directly from
$([i]^*)^2=-1$ on $H^1$; homological equivalence makes this calculation
an equality of correspondences in the stated category.

For the usual Hermitian extension of the orthonormal form on $e,f$,
\[
 (u,u)_{\rm Herm}=2,\qquad
 \psi(u,\overline u)=-2i,\qquad
 (u,u)_{\rm bilinear}=0.                         \tag{IF-10c}
\]
These are three different pairings. Define the normalized Hermitian
form on $U$ by dividing its ambient Hermitian form by $2$; then $u$ has
norm one. Equivalently its value on $U$ is
$\psi(x,\overline y)/(-2i)$. This last expression is restricted to
$U$, where it is integral and positive under the chosen complex
embedding; it is not asserted to be an integral pairing on all of
$H^1(E,R)$.

For a column of height $h$, use the actual cohomological space
$H^h(E^r,K')$. For each $h$-element subset
$I=\{i_1<\cdots<i_h\}$ of the $r$ factors, project to $M$ in those
factors and to $h^0(E)$ in the others. Its integral generator is the
ordered cup product $u_{i_1}\cdots u_{i_h}$. These generators occupy
distinct K\"unneth summands and each is primitive, so their direct sum
is a saturated lattice canonically identified with
$\Lambda^h R^r\otimes U^{\otimes h}$. Applying this construction in
separate blocks for all column heights $h_c$ realizes
\[
 A_\lambda(R)=\bigotimes_c\Lambda^{h_c}R^r
 \quad\text{as}\quad
 A_\lambda(R)\otimes U^{\otimes d}
 \subset H^d(E^{rq},R),\qquad d=|\lambda|,
                                                               \tag{IF-10d}
\]
where $q$ is the number of columns. We use actual ordered cup products
and the resulting permutation signs. In particular this is the ordinary
exterior power of $H^1$ in cohomology, not an unqualified categorical
exterior power of the odd motive $h^1(E)$.

Every summand indexed by a tuple of subsets in (IF-10d) is isomorphic
to $M^{\otimes d}$. The isomorphism orders the selected factors and
discards the degree-zero factors. Its inverse inserts the origin and
the same degree-zero factors. These are induced by projections,
coordinate permutations, and the degree projectors. Compose the
projection onto one such summand, its ordered identification with
$M^{\otimes d}$, and the inverse identification for another summand.
The resulting algebraic correspondences $E_{JI}$ act by
\[
 E_{JI}(v_I)=v_J,\qquad E_{JI}(v_{I'})=0\quad(I'\ne I),
 \qquad E_{JI}E_{LK}=\delta_{I,L}E_{JK}.           \tag{IF-10e}
\]
The necessary signs are fixed by the chosen increasing order of the
degree-one factors. This provides all matrix units on the coefficient
space; it also provides rectangular matrix units between two ambient
powers with the same $d$.

Choose an integral basis for the primitive sublattice $T_\lambda$
proved in IF-08 and let $C$ be its matrix in the orthonormal basis of
$A_\lambda$. The rational matrix
\[
 P_T=C(C^tC)^{-1}C^t                              \tag{IF-10f}
\]
is the projector to its rational span. Inserting its entries into
(IF-10e) makes it an algebraic projector in the homological category.
Its integral intersection is $T_\lambda\otimes R$: primitivity over
$\mathbb Z$ gives a free quotient and is preserved by scalar extension.
All denominators in (IF-10f) are explicit divisors of
$\det(C^tC)$, before cancellation. Together with the factors $2$ from
(IF-10b), they record the denominator of this chosen correspondence.
The ambient tensor basis in (IF-10d) has squared norm $2^d$; after
multiplying the restricted Hermitian form by $2^{-d}$ its Gram matrix
on $T_\lambda$ is the one in IF-04. Thus the normalized Fock isometry
is exact, including its pairing normalization.

Every coefficient matrix $A_\beta$ of IF-08, and hence every
$B_J=\sum_\beta J(\beta)A_\beta$ with $J(\beta)\in K'$, is now an
algebraic correspondence by (IF-10e), restricted on each side by
(IF-10f). Its matrix on the chosen lattice is unchanged. This argument
does not assert that the coefficients of the representing cycle are
integral; the matrix on the specified lattice is integral exactly
when the corresponding matrix of IF-08 is integral.

To handle determinant characters and arrows of different total
degrees, note that the determinant line of $U^{\oplus r}$ is
$M^{\otimes r}$ with generator $u^{\otimes r}$, using ordered cup
products. Negative tensor powers use the abstract dual of the chosen integral
line from the outset, including in the determinant twist. The polarization
identifies $M^\vee$ with the conjugate eigenmotive $\overline M(1)$;
the scalar in (IF-10c) specifies its integral normalization. These
rank-one motives are invertible in the homological category: evaluation
and its transpose, divided by their nonzero scalar, are inverse on the
one-dimensional realization, hence inverse there. Consequently the
coefficient motive for $(\lambda,s)$ is
\[
 \mathcal V_{\lambda,s}=T_{\lambda,K'}\otimes M^w,
 \qquad w=d+rs.
\]
Normalize it by the specified inverse CM character:
\[
 \mathcal V_{\lambda,s}\otimes M^{-w}
       \simeq T_{\lambda,K'}\otimes\mathbf1.       \tag{IF-10g}
\]
Use the abstract dual of the chosen integral rank-one lattice for this
evaluation, so that its generator evaluates to one. Formula (IF-10c)
says exactly what this means inside a conjugate geometric realization:
$U^\vee$ is represented by
$(-2i)^{-1}\overline U(1)$. In particular it is not the native
saturated lattice $\overline U(1)$ at primes above $2$. If one instead
uses native integral eigenlines in a pair of opposite powers
$M^w,M^{-w}$, evaluation has image the ideal generated by
$(2i)^{|w|}$, up to a unit and the ordering sign. This factor is retained,
or is removed by precisely the displayed abstract-dual normalization;
there is no unnoticed saturation at $2$. This $|w|$ formula concerns
the already collapsed pair $M^w,M^{-w}$. If $s<0$ and one begins
instead with the raw native tensor of $d$ copies of $U$ and
$r|s|$ copies of $\overline U(1)$, converting the latter to abstract
duals contributes $(-2i)^{-r|s|}$. Collapsing the raw native tensor by
pairings first contributes $(-2i)^{\min(d,r|s|)}$, up to the ordering
unit. These factors must be retained in any comparison with that
different native presentation.

After (IF-10g), every rectangular matrix, including the determinant
insertion and maps between different $w$, is an algebraic correspondence
between sums of the unit motive. This realizes exactly the original
coefficient matrices because the evaluation generators were fixed to
one. Finite direct sums, invariant subspaces and level subspaces are
handled by matrix projectors over $K'$ and their specified integral
intersections. An arbitrary projective global lattice is retained as
that lattice; a free global basis is not assumed. Thus the block sums
and local-frame conjugations of (IF-08k), and the terminal dual
functional, are unchanged.

Finally a finite torsion-free extension of rings of integers is
faithfully flat. The minors defining every determinantal ideal are
the same minors after extension to $R$; contracting the extended ideal
recovers the original one. Over a local extension of discrete valuation
rings, a valuation is multiplied by the ramification index, and this
known factor recovers the original length. This proves the index
assertion at every prime, with the abstract-dual and cycle-denominator
conventions just specified.
\end{proof}

\paragraph{The actual algebraicity input.}
The coefficient-field hypothesis is already part of a supplied
algebraic theta-chain datum, and is not proved by calling an arbitrary
complex derivative algebraic. The exact interfaces are
\texttt{ThDistC.lean}, \texttt{Amb.JetsRat} (lines 240--242), and
\path{KernelLocalFactors/ThDistC.lean}, \texttt{jetK\_spec}
(lines 615--621), which takes that rationality hypothesis and identifies
the descended moment with the period-normalized derivative.
For the actual tied kernel, \texttt{JetFieldLeKR.lean},
\texttt{KerEntTied.KerJets} (lines 113--116), includes the product
\texttt{rhoDet*rhoDet'} in the algebraic coefficient. One therefore
uses the combined moment $\widetilde J=\rho\rho'J$ when only that
combination is certified algebraic; algebraicity of its separate
factors is unnecessary. Alternatively,
\texttt{DataForms.lean}, \texttt{KRat,KL,KL\_spec} (lines 252--259),
gives the algebraic global map together with its exact complex
realization. These declarations are hypotheses and conditional descent
theorems; they do not by themselves construct a new closed global
instance. IF-10 applies to every instance supplied with this data. In
particular the fixed \texttt{ThE} test of IF-09 is unconditional over
$\mathbb Q$ at the level of its finite moments.

\paragraph{Actual weights and the limit of Tate normalization.}
For one embedding factor $\tau$ of the three actual column coefficients
in IF-09, direct summation gives
\[
 (|\alpha_0|,|\alpha_1|,|\alpha_2|)=(113,841,889),
 \qquad (w_0,w_1,w_2)=(137,-1463,-3351).           \tag{IF-10h}
\]
At the modified real place of \texttt{numBig}, all three $w_j$ decrease
by $3$; their successive differences remain $-1600,-1888$.
The actual coefficient module is a tensor product over $\tau$; apply
(IF-10g) to each factor and then tensor. If the same $M$ is used at
every place, its total character exponent is $\sum_\tau w_{j,\tau}$,
not a single entry of (IF-10h). The CM twists in (IF-10g) are therefore
substantial and are not pure
Tate twists. Indeed, under the selected embedding $M^w$ has Hodge type
$(w,0)$. A nonzero algebraic correspondence between
$M^w(a)$ and $M^{w'}(b)$ would preserve Hodge type, requiring
\[
 (w-a,-a)=(w'-b,-b).
\]
It follows that $a=b$ and $w=w'$. Hence in this natural single-eigenline
model the nonzero arrows of the actual coefficient chain cannot become
algebraic by Tate twists alone. The explicit CM-character normalization
is necessary for this construction. This is an obstruction to this
specified model, not a claim about every conceivable realization.

On an $\ell$-adic realization, if the eigenline has character $\eta$,
the untwisted coefficient motive has character $\eta^w$ on its rank-one
factor (with the sum of exponents when tensoring over $\tau$). Formula
(IF-10g) cancels that character. Thus this construction
does not preserve an unspecified automorphic $L$-function or supply its
motivic realization. It does give algebraic correspondences realizing
the fully specified finite coefficient matrices and their exact local
indices. In the language of cote 16, scans 014--018, it constructs
operators in a geometric correspondence ring after the stated
normalizations; IF-09 prevents replacing all such two-step kernels by
positive powers of one fixed Lefschetz element. The $\lambda$-operations
occur in their actual cohomological exterior powers, while $K_0$ alone
forgets the retained lattice indices and the prime-$2$ dual-lattice
factor.

The equation for $E$ and these coordinate correspondences spread over
$\mathbb Z[i,1/2]$. Statements about the transported lattice at rational
prime $2$ concern integral Betti cohomology or the specified $2$-adic
representation lattice. They do not assert good reduction of this
elliptic curve at $2$, or an integral Chow projector there.


\subsection{IF-11: an inhabited two-step algebraic theta-jet specialization}
The numerical examples in IF-09 are not the only source of coefficient
data. A separate supplied construction gives a genuine theta family
with $m=2$ and all the algebraicity needed by its finite kernels. This
does not assert the additional global size or nonzero-period hypotheses.

All paths in this subsection beginning with \path{Development/} are
paths in Alp\"oge's Lean archive. Short paths beginning with \path{Proofs/}
or \path{Definitions/} below are relative to
\path{Development/AbelianVarietyTheory/}. The theorem
\path{GlobalTheory.KerT.JetsInhab.exists_colIn_jets}, in
\path{Development/AbelianVarietyTheory/Proofs/GlobalInterpolation/}
\path{ExistsHypJetsAll.lean}, lines 73--115, supplies
\[
 \exists B,X,I,\quad
 \texttt{HasLetters}(B)\ \wedge\ I.\texttt{Ok}\ \wedge\ I.\texttt{JetsIn}.
                                                               \tag{IF-11a}
\]
We use this as a cited theorem. Its proof builds its base from
\path{ExistsNatBoundPre.lean}, \path{exists_pre_two_words}
(lines 29--132), and then fixes the distinguished word in
\texttt{exists\_pre\_sized} (lines 19--71 of
\path{ExistsHypJetsAll.lean}). Its rank datum is
\path{Formalism/SharpRestriction/InertJacobiLevels/HostFam.lean},
\texttt{Inhab.Rt 0} (lines 14--29). Thus
\[
 k=8,\quad m=2,\quad t=15,\quad f=8,\qquad
 (N_0,N_1,N_2)=(8,24,40).
\]

\paragraph{Statement IF-11.}
The witness constructed in (IF-11a) has the following coefficient
specialization, constant over its eight embedding factors:
\[
\begin{split}
 u&=(49,45,37,29,21,13,5,-1)/2,\\
 v&=(24,20,16,12,8,4,0),\\
 b&=(1,37),\qquad e=(19,55),\\
 \lambda_0&=(21,20,17,14,11,8,5,3).
\end{split}                                                     \tag{IF-11b}
\]
Its actual column determinant exponents and polynomial parts are
\[
 (s_0,s_1,s_2)=(3,-15,-43),
\]
\[
\begin{split}
 \alpha_0&=(18,17,14,11,8,5,2,0),\\
 \alpha_1&=(28,27,24,21,18,15,12,10,0^{[16]}),\\
 \alpha_2&=(48,47,44,41,38,35,32,30,20^{[16]},0^{[16]}).
\end{split}                                                     \tag{IF-11c}
\]
The two finite kernels have shapes
\[
 \nu_1=(28,27,24,21,18,15,12,10),\qquad
 \nu_2=(48,47,44,41,38,35,32,30,20^{[16]}),          \tag{IF-11d}
\]
with determinant insertions $(10,20)$, of degrees $80,480$ per
embedding. All required normalized moments lie in the finite number
field supplied by (IF-11a). Hence IF-08 and IF-10 apply to this actual
theta-jet family, not only to the exponential test series.

The distinct nonzero simple-root gaps in these three coefficient
lattices are $1,2,3,10,20$. The gap $1$ has no second descendant; for
the others the exact second-descendant invariants are
\[
\begin{array}{c|c|c|c}
 d&F^2\text{ on Weyl}&F^2\text{ on normalized dual}
       &F^{(2)}\text{ on normalized dual}\\\hline
 2&2&2&1\\
 3&2&6&3\\
 10&2&90&45=3^2\cdot5\\
 20&2&380&190=2\cdot5\cdot19
\end{array}.                                                    \tag{IF-11e}
\]
On every nonzero displayed Weyl second-descendant line,
$F^{(2)}$ has invariant one. These are coefficient-root indices,
not the values or Smith invariants of the full global theta composition.

\begin{proof}
The parameter construction
\path{Proofs/IwasawaDecomposition/ExistsHarishChandraWeights.lean},
\path{exists_preCh_of_prescribed} (lines 12--33), gives
\[
 v_l=(k-2-l)4\Delta,\qquad
 b_j=b_1+2hj,\qquad e_j=b_j+h,
\]
and endpoint values $u_0=u_F+1/2$, $u_{k-1}=u_L+1/2$, with
$u_{l+1}=v_l-1/2-\Delta$ for the distinguished all-true word.
The invocation in \path{exists_pre_two_words} has
\[
 \Delta=1,\quad h=t+3=18,\quad u_F=24,\quad u_L=-1,
 \quad b_1=1.
\]
Its distinguished word is $0\in\operatorname{Fin}2$, on which
\texttt{decide}(u=0) is true. These substitutions prove (IF-11b).
The passage to the one-word base keeps this weight family. The field
enlargement retains the preliminary rank and weight data, and the
subsequent change of ambient theta data retains the resulting base $B$. Thus the actual witness in (IF-11a) retains these
parameters; no selection of an unrelated existential witness is used.

Subtracting the rank-eight half-sum of positive roots gives the
displayed $\lambda_0$. From \texttt{Params.wG} one obtains
$w_G=(7,27)$, so the unshifted weights are
\[
 \lambda^N_1=(\lambda_0,-7^{[16]}),\qquad
 \lambda^N_2=(\lambda_0,-7^{[16]},-27^{[16]}).
\]
Subtract $8j$ at stage $j$, as required by
\texttt{Wts.lamTrue}, and then subtract the last entry to obtain
(IF-11c). Here
\[
 \mathrm{sSharp}=(7,35),\qquad
 \mathrm{sDiv}=(15,43),\qquad
 n_j=s_j+\mathrm{sSharp}_j=(10,20).
\]
Adding these $n_j$ to each source polynomial part gives (IF-11d),
which also agrees by direct substitution with \texttt{nuRow}.
The determinant degrees are $8\cdot10=80$ and $24\cdot20=480$.
Successive differences in (IF-11c) give the stated root gaps.
The invariant table is now exactly IF-04, with the specified Weyl
or normalized-dual lattice; no integral identification of those two
lattices is assumed.

It remains to check that the cited algebraicity is exactly the one
used by the kernel. In
\path{Definitions/AutomorphicMembership/JetsInIffStepField.lean},
\texttt{ColIn.jetNum} (lines 99--103) is
\[
 J(\beta)=\chi_F^{-1}
       \left(\prod_{\tau,u,i}c_{D,\tau}^{\beta_{\tau,u,i}}\right)
       (\partial^\beta\Theta)(0),                 \tag{IF-11f}
\]
where $\Theta$ is the actual \texttt{thInd} family. This is precisely
\texttt{ThDistC.jetC}, with ordinary derivatives and no division by
$\beta!$. Its \texttt{JetsIn} definition (lines 109--113) asserts
membership in the coefficient field for multidegree
$\sum_r\nu_{j,r}$ at each $\tau$, and
\path{JetsIn.exists_eq,isAlgebraic} (lines 142--150) identifies
these numbers with their field-valued versions.

Every coefficient $A_\beta$ used by the polynomial kernel of shape
$\nu_j$ has exactly this multidegree: each exterior column contributes
its height, so the total is $|\nu_j|$ at that embedding. Consequently
the degree-restricted algebraicity hypothesis proves algebraicity of
every moment that actually occurs in the finite operator. An
all-degree hypothesis is unnecessary. The degrees here are
$155,635$ per embedding, or $1240,5080$ over the eight embeddings.

The field existence is not an assertion about the default-valued
function \texttt{desc}. Its supplied proof is
\path{Proofs/IwasawaDecomposition/DegreeEqTotDeg.lean},
\texttt{finAlg\_step} (lines 193--226), followed by
\texttt{exists\_jetsIn} (lines 283--313): a finite algebraic generating
set of the required theta jets is placed in the coefficient field,
using the supplied normalized-vacuum algebraicity theorem
\path{GlobalTheory.T.vacuum_entire_and_normalizedJets_algebraic}
in \path{VacuumEntireAndNormalizedJetsAlgebraicStatement.lean},
and theta equivariance and lattice reduction.
\path{ExistsColInJetsIn.exists_bound_jetsIn} (lines 34--40)
specifies the finite algebraic generating set to adjoin; the
\path{BaseExt.exists_base} call in \path{exists_colIn_jets}
(lines 85--96) performs that adjunction while retaining the preliminary
rank and weight data.
Finally \path{Definitions/GlobalTwisting/JetsAllOfColIn.lean},
\texttt{jetsAll\_of\_colIn} (lines 108--116), identifies these
degree-restricted jets with the actual assembled kernel's
\texttt{JetsAll}. This proves exact compatibility with IF-08 and
discharges IF-10's coefficient-field input for this inhabited family.
The cited theorems are used as supplied results; we have not run a
new Lean kernel check of their proofs.
\end{proof}

\paragraph{A concrete derivative test and a full-module invariant.}
Write $\Delta_h(Z)$ for the leading $h$-row, $h$-column minor.
At one embedding factor the two highest flag polynomials for (IF-11d)
are the explicit polynomials
\[
 P_1=\Delta_1\Bigl(\prod_{h=2}^6\Delta_h^3\Bigr)
            \Delta_7^2\Delta_8^{10},\qquad
 P_2=P_1\Delta_{24}^{20}.                         \tag{IF-11g}
\]
For the full coefficient tensor put
$P_j^{\rm tot}=\prod_\tau P_j(Z_\tau)$. In standard coefficient
frames, before the scalar factors in (IF-08k), the normalized
highest-to-highest block coefficient is the moment functional
$J(P_j^{\rm tot})=\sum_\beta[Z^\beta]P_j^{\rm tot}\,J(\beta)$.
For other frames one uses the complete matrix conjugation (IF-08k);
one does not replace it by a scalar change of this one entry.
To prove the standard-frame formula,
apply the target highest functional to $A_{\nu_j}(Z)t_{\nu_j}$.
Cauchy--Binet selects the leading minors at each embedding, which gives
$P_j^{\rm tot}(Z)$; then contract its coefficients with $J$. Thus
(IF-11g) specifies an actual differential-operator test of degrees
$155,635$ per embedding, or $1240,5080$ in total. For the
coordinate exponential moments it equals one; no such value is
asserted for the actual analytic theta family.

The source coefficient representation has the wrong rank to be a
whole symplectic exterior algebra. The Weyl dimension
formula cited in IF-05 gives
\[
 \operatorname{rank}T_{\alpha_0}(8)
  =2^{28}\cdot3\cdot5^2\cdot7\cdot11^2\cdot19\cdot23.
                                                               \tag{IF-11h}
\]
This follows by multiplying
$\prod_{i<j}(\alpha_{0,i}-\alpha_{0,j}+j-i)/(j-i)$; the supplied
exact prime-exponent computation records each factor. Since a full
$V_n$ has rank $4^n$, even rational isomorphism to a whole $V_n$ is
impossible. The eightfold coefficient tensor retains the odd prime
factors. This does not obstruct the primitive embeddings or tensor
realizations constructed above, and says nothing about dimensions
of global invariant-form subspaces.

The per-embedding CM exponents of IF-10 are now
$(99,-205,-1085)$, with sums $(792,-1640,-8680)$ over all embeddings.
The same Tate-only obstruction and explicit CM-character untwist apply.
For the coordinate exponential family at these partitions, the positive
common-root construction of IF-09 fits inside $V_{39\cdot48}=V_{1872}$
per embedding, and inside $V_{14976}$ after tensoring all eight factors.

\paragraph{Exact scope.}
The input in (IF-11a) certifies a genuine algebraic theta-jet family.
It does not assert a nonzero composed terminal period, positive full
\texttt{sg}, or the full \texttt{SizeData}$_2$ hypotheses.
Indeed the separate source check SF-Q6-02 shows that the named witness
has top endpoint $b_T=-3$, incompatible with that full size condition.
No full step-10 divisibility bound is inferred from (IF-11e).
The movement here is that IF-08 and IF-10 now apply to an explicitly
traced, inhabited source theta family with the exact derivative
normalization, rather than to a hypothetical list of algebraic moments.
The algebraic geometry and $\ell$-adic translation are precisely those
of IF-10: specified coefficient correspondences and lattices, with the
CM character and all local denominators retained.


\subsection{IF-12: genuine CM theta kernels need not commute with the fixed root}
The common-root construction at the end of IF-09 is a positive result
for the specified coordinate-exponential test family. It is not a
root-invariance theorem for every theta jet. We give both a completely
specified CM theta counterexample and an existence statement for the
torsion-translation family of an actual nonzero source CM theta vector.
Neither assertion identifies an individual translate with the prescribed
global automorphic class sum.

\paragraph{The source operator and its normalization.}
For a positive Hermitian form $H$, linear in its first argument, put
\[
 U_x f(z)=f(z+x)\exp\{-\pi H(z,x)-\pi H(x,x)/2\}.       \tag{IF-12a}
\]
This is exactly \path{Formalism.L2.Th1.H.opU} in
\path{Formalism/LocalToGlobal/TorsionIntegrality/}
\path{QuadraticSemicharacters/IsSemiChar.lean}, lines 21--22,
and \texttt{TaylorLogMetric.U} in
\path{Formalism/LocalToGlobal/KernelLattices/ColumnData.lean},
lines 101--102. The latter uses
$H(z,x)=\operatorname{area}(\Lambda)^{-1}\sum z_i\overline{x_i}$.
We retain the Gaussian in (IF-12a); replacing this operator by bare
translation would change the jets.

For a polynomial $P=\sum_\beta p_\beta Z^\beta$ write
\[
 J_f(P)=\sum_\beta p_\beta\partial^\beta f(0).
                                                               \tag{IF-12b}
\]
These are ordinary derivatives, with no factorial division. This is
the convention of \texttt{ThDistC.derivAt0} and \texttt{jetC}
(\texttt{ThDistC.lean}, lines 56--61). If $P$ is homogeneous of degree $D$ in one
embedding factor, the source normalized functional is
$\chi_F^{-1}c_D^D J_f(P)$. With several embedding factors it is
$\chi_F^{-1}\prod_\tau c_D(\tau)^{D_\tau}J_f(P)$.
For nonzero $\chi_F,c_D(\tau)$ these common factors cannot turn a
nonzero commutator coefficient into zero. This statement concerns the
fixed multidegree, and does not discard factors between different
degrees.

\paragraph{Statement IF-12.}
Use the genuine first-arrow shapes from IF-11:
\[
 \alpha=(18,17,14,11,8,5,2,0),\qquad
 \nu=(28,27,24,21,18,15,12,10).
\]
Let $D_{10}:T_\alpha(\mathbb Z^8)\longrightarrow
T_\nu(\mathbb Z^8)$ be the determinant insertion, and let
\[
 B_f=\sum_\beta J_f(Z^\beta)A_{\nu,\beta}D_{10}:
 T_\alpha(\mathbb C^8)\longrightarrow T_\nu(\mathbb C^{24})
                                                               \tag{IF-12c}
\]
be the exact transported finite kernel of IF-08. Take $F$ to be the
row root $E_{21}$, on both spaces. There is a genuine CM theta vector
$f$ in the source's \texttt{Th1.H.Th} model, given explicitly below,
such that
\[
                 F B_f-B_f F\ne0.                            \tag{IF-12d}
\]
Its period ratio is $i10^{400}$, and its translation point has order
dividing $12$. It is the unitary translate of the source's normalized
product theta \texttt{F0}, not a formal exponential series. It has
the original semicharacter on the common refined period lattice
$12\Lambda^{192}$.

More generally, let $f_0\ne0$ be any entire theta vector for a full
integral polarized lattice $\Lambda\subset\mathbb C^q$, with the
operator (IF-12a), and let $P\ne0$ be any polynomial. For every
integer $\ell\ge2$ there exist $n\ge0$ and
$x\in\ell^{-n}\Lambda$ such that
\[
                         J_{U_xf_0}(P)\ne0.                  \tag{IF-12e}
\]
The translate retains the original semicharacter on
$\ell^n\Lambda$. Applying this to the polynomial in the proof
below gives (IF-12d) in the torsion-translation family of an actual
nonzero source CM theta vector, even if denominators are restricted
to powers of one chosen allowed odd prime. It does not assert failure
of (IF-12d)'s opposite equality for a particular preassigned translate
or global class sum.

\begin{proof}
First isolate an exact coefficient, including the determinant
insertion. Let $t_\alpha$ be the highest tensor and
$t=D_{10}t_\alpha$. In the target, use the first eight coordinates.
The column-height multiplicities of $\nu$ are
\[
          (a_1,\ldots,a_8)=(1,3,3,3,3,3,2,10).
\]
Thus $r=Ft$ is the single tensor obtained by changing $e_1$ to
$e_2$ in the unique $\bigwedge^1$ factor. All factors of height
at least two contain both $e_1,e_2$ and are killed by $F$.
Consequently $r$ is an ambient unit basis vector and $E_{12}r=t$.
The insertion $D_{10}$ commutes with $F$ because a full determinant
has trivial $\mathfrak{sl}_8$ action. With $\Delta_h(Z)$ the leading
$h$-minor of the $24$ by $8$ matrix $Z$, define
\[
 \begin{split}
 Q(Z)&=\prod_{h=2}^{6}\Delta_h(Z)^3\,
                         \Delta_7(Z)^2\Delta_8(Z)^{10},\\
 P(Z)&=(Z_{11}-Z_{22})Q(Z).
 \end{split}                                                  \tag{IF-12f}
\]
The tensor exterior form is the normalized Fock form by IF-04.
Taking the coefficient of $r$ in the commutator on $t_\alpha$
therefore gives the exact identity
\[
       \langle r,(F B_f-B_fF)t_\alpha\rangle=J_f(P).            \tag{IF-12g}
\]
Indeed the first summand has polynomial coefficient $Z_{11}Q$ and
the second has coefficient $Z_{22}Q$. No scalar $c_\nu$ remains:
it has already been included in the normalized Fock/exterior
identification. The functional $\langle r,-\rangle$ is an integral
dual-Weyl test, since $r$ is an integral unit tensor. The polynomial
$P$ has degree $D=155$ and is nonzero.

We construct an explicit genuine theta vector. Put $T=10^{400}$,
$\tau=iT$, $\Lambda_T=\mathbb Z+iT\mathbb Z$, and
\[
 \vartheta_T(z)=\sum_{n\in\mathbb Z}
                   e^{-\pi T n^2}e^{2\pi i n z},\qquad
 g_T(z)=\frac{e^{\pi z^2/(2T)}\vartheta_T(z)}{\vartheta_T(0)}.
                                                               \tag{IF-12h}
\]
Here $\vartheta_T(0)>0$. Let $L_T$ be the period pair with lattice
$\Lambda_T$ and let $\omega_0=-(1+\tau)$. Then
\[
                g_T=\texttt{Th1.F.th0}(L_T,\omega_0/2).
                                                               \tag{IF-12i}
\]
For completeness, the exact source definition is
$\texttt{th0}(L,h)=U_h\operatorname{sgStar}(L)/
(U_h\operatorname{sgStar}(L))(0)$ in
\path{EisensteinTranslationProducts/SemiChar.lean}, line 34.
The supplied identity
\path{GlobalTheory.SgStarOddTheta.exp_mul_oddTheta_eq_deriv_mul_sgStar}
in \path{Development/AbelianVarietyTheory/Proofs/GlobalInterpolation/}
\path{GaussianOddJacobiThetaEqDerivMulSgStar.lean}, lines 100--124,
identifies the Gaussian-gauged odd Jacobi theta with
$\operatorname{sgStar}$ times its derivative at zero. Its odd theta
is defined in \path{DifferentiableOddTheta.lean}, lines 19--20, by
\[
 e^{\pi i\tau/4+\pi i(z+1/2)}
       \vartheta_T\bigl(z+(\tau+1)/2\bigr).
\]
Substitute $h=-(\tau+1)/2$ into (IF-12a). The linear exponential
in $z$ cancels, leaving a nonzero constant times
$e^{\pi z^2/(2T)}\vartheta_T(z)$. Division by its value at zero
gives (IF-12i), including its scalar. The two hypotheses of that
source identity are supplied by \texttt{Th1.F.thetaLine} in
\path{EisensteinTranslationProducts/EllipticOnePole.lean}, line 27,
and \texttt{Th1.F.sgStarBasic} in
\path{EisensteinTranslationProducts/QuasiPeriodConstUnique.lean},
lines 21--22. We use these supplied results as cited theorems.
Equivalently, their
quasiperiod law and normalized value at zero give the same identity.

Set $q=24\cdot8=192$ and
$G_T(Z)=\prod_{i,j}g_T(Z_{ij})$. Thus $G_T$ is exactly the source
$\texttt{Th1.H.F0}(L_T,\omega_0,\operatorname{Fin}24\times
\operatorname{Fin}8)$, by
\path{QuadraticSemicharacters/ProductThetaMem.lean}, lines 10--11.
Let $a$ be the $24$ by $8$ real matrix with
\[
 a_{11}=1/4,\quad a_{22}=1/3,\quad
 a_{jj}=1/4\ (3\le j\le8),\quad a_{ij}=0\ (i\ne j),
                                                               \tag{IF-12j}
\]
and put $f_T=U_{iTa}G_T$, with
$H(Z,W)=T^{-1}\sum Z_{ij}\overline{W_{ij}}$.
Direct substitution into (IF-12a), without a gauge change, gives
\[
 f_T(Z)=\vartheta_T(0)^{-q}e^{\pi\sum Z_{ij}^2/(2T)}
      \sum_{n\in\mathbb Z^q}
       e^{-\pi T\|n+a\|^2}e^{2\pi i\langle n+a,Z\rangle}.
                                                               \tag{IF-12k}
\]
All these series and their derivatives converge normally on compact
sets: a fixed derivative contributes only a polynomial in $n$,
whereas the Gaussian exponent is negative quadratic in $n$.

This is an allowed analytic theta function in the source's exact
theta model. Its original product theta has semicharacter
$\prod(-1)^{m_{ij}n_{ij}}$ at lattice points $m+iTn$; substitution
in (IF-12h) verifies the law. Since $12a$ is integral, the Heisenberg
commutator for $iTa$ and any vector in $12\Lambda_T^q$ is one.
Consequently $f_T$ has the original semicharacter, here trivial,
on $12\Lambda_T^q$, and belongs to the source's
$\texttt{Th}(H,1,12\Lambda_T^q)$. The latter definition requires
analyticity and the stated translation law, not self-duality of the
refined lattice. Also $iT$ is imaginary quadratic, so the underlying
elliptic curve and its products are CM; $iTa$ is a torsion point of
order dividing $12$. We do not identify this refined lattice with
the restricted self-dual lattice of a particular \texttt{D7.Amb}.

There is a useful exact finite-degree heat formula. Write
$\Delta=\sum_{i,j}\partial_{X_{ij}}^2$. For a homogeneous polynomial
of degree $D$, ordinary differentiation of the Gaussian gives
\[
 \begin{split}
 S_T&:=\vartheta_T(0)^q e^{\pi T\|a\|^2}(2\pi i)^{-D}
                                                    J_{f_T}(P)\\
 &=\sum_{n\in\mathbb Z^q}
    e^{-\pi T(\|n+a\|^2-\|a\|^2)}
       \left(e^{-\Delta/(8\pi T)}P\right)(n+a).
 \end{split}                                                  \tag{IF-12l}
\]
The heat exponential is a finite sum on this polynomial. To verify
its sign and factor, in one variable divide each derivative by
$2\pi i$. Pairing two derivatives on
$e^{\pi z^2/(2T)}$ contributes
$(\pi/T)/(2\pi i)^2=-1/(4\pi T)$; the exponential differential
operator therefore has coefficient $-1/(8\pi T)$. The multinomial
expansion proves (IF-12l) in all variables. This also verifies that
there is no missing multi-index factorial.

We now give an explicit error bound; no numerical evaluation of an
exponentially small theta value is used. Let $M=\|P\|_1$ be its
coefficient $\ell^1$ norm. Expansion of determinants gives
\[
 M\le 2\prod_{h=2}^{6}(h!)^3(7!)^2(8!)^{10}
                  \le2(8!)^{27}<2\cdot10^{135}.                \tag{IF-12m}
\]
At the diagonal matrix (IF-12j),
\[
           P(a)=-4^{-128}3^{-28},\qquad |P(a)|>10^{-156}.
                                                               \tag{IF-12n}
\]
For $R=\max(1,\|x\|_\infty)$ and
$u=qD^2/(8\pi T)$, termwise differentiation gives
\[
 |(e^{-\Delta/(8\pi T)}P)(x)|\le M R^D e^u,
 \quad |(e^{-\Delta/(8\pi T)}P)(a)-P(a)|\le M(e^u-1).
\]
For example, at each differentiation step a coefficient is
multiplied by at most $D^2$ for each of the $q$ summands of $\Delta$;
the degree decreases, so the displayed bound applies to every term.
Here $qD^2=4{,}612{,}800<10^7$, so $u<10^{-393}$ and
\[
                  M(e^u-1)<4\cdot10^{-258}.                   \tag{IF-12o}
\]
For $n\ne0$, using $|a_{ij}|\le1/3$ and integer coordinates gives
\[
 \|n+a\|^2-\|a\|^2\ge\|n\|^2/3,
 \qquad
 (1+\|n\|_\infty)^D\le e^{D\|n\|^2}.
\]
Since $T\ge2D$ and $\pi\ge3$, the absolute tail in (IF-12l) is at
most
\[
 2M\sum_{n\ne0}e^{-T\|n\|^2/2}
 \le16qM e^{-T/2}
 \le128qM/T^2
 <5\cdot10^{-661}.                                            \tag{IF-12p}
\]
To check the middle inequality, set $r=e^{-T/2}$. The
one-coordinate theta sum is at most $1+4r$; thus the sum over
nonzero vectors is at most $(1+4r)^q-1\le8qr$ when $4qr\le1$.
The last bound uses $e^{T/2}\ge T^2/8$. All of these conditions
hold for the specified integer $T$. The bounds (IF-12o)--(IF-12p)
are smaller in total than $|P(a)|/2$, proving $S_T\ne0$ and hence
$J_{f_T}(P)\ne0$. Equations (IF-12g) and (IF-12l) prove the explicit
counterexample. A homogeneous change to CM period coordinates, or
the common source factors in (IF-12b), multiplies this nonzero
coefficient by a nonzero scalar and does not affect the conclusion.

Finally prove (IF-12e), which does not vary the initial CM datum.
For simplicity choose coordinates on $\mathbb C^q$ and write
$H(z,x)=z^{\mathsf t}C\overline x$, with $C$ invertible. If
$d=\deg P$, multiplying $J_{U_x f_0}(P)$ by
$e^{\pi H(x,x)/2}$ gives a polynomial in $\overline x$, of degree
at most $d$, whose coefficients are holomorphic functions of $x$:
\[
 R(x)=\sum_\beta p_\beta\sum_{\gamma\le\beta}
    {\beta\choose\gamma}(-\pi C\overline x)^{\beta-\gamma}
                                      (\partial^\gamma f_0)(x).
                                                               \tag{IF-12q}
\]
Its part of top antiholomorphic degree is
$f_0(x)P_d(-\pi C\overline x)$, where $P_d$ is the leading
homogeneous part of $P$. Since $C$ is invertible and $P_d\ne0$,
choose a multi-index $\eta$ of degree $d$ with nonzero coefficient
$b_\eta$ in $P_d(-\pi C\overline x)$. Wirtinger differentiation
then gives
\[
             \partial_{\overline x}^{\eta}R(x)
                        =\eta!\,b_\eta f_0(x).                \tag{IF-12r}
\]
Every lower antiholomorphic term vanishes, and holomorphic
coefficients have zero antiholomorphic derivatives. Thus $R$ is not
identically zero. Its nonzero set is nonempty and open. A real basis
of the full lattice identifies
$\bigcup_n\ell^{-n}\Lambda$ with a dense subset of
$\mathbb R^{2q}$, so this set contains a point $x$ proving
(IF-12e). If $\lambda\in\ell^n\Lambda$, integrality of
$\operatorname{Im}H$ on $\Lambda$ makes the commutator of
$U_x$ and $U_\lambda$ equal to one. Hence the original theta law
holds on that refined lattice with its original semicharacter.
\end{proof}

\paragraph{Application to the actual eight-factor source family.}
Choose the nonzero source vector
$f_0=B.\texttt{thZero}(8,24)$.
The supplied theorem \path{WeilVac.thZero_ne_zero}, in
\path{Development/AbelianVarietyTheory/Proofs/IwasawaDecomposition/}
\path{ExistsWeilOThZeroTriv.lean}, lines 38--42, proves its
nonvanishing under \texttt{B.StdOK}; \texttt{thZero\_mem}
(lines 80--103) places it in the actual CM theta space. Its period
lattice is full: it is built from a nonzero fractional ideal in the
CM field under its CM embeddings, and coordinate tensor powers of
that lattice. The CM embeddings identify its real span with all the
ambient complex coordinates. The source polarization is integral on
this lattice and has positive scalar matrix $\operatorname{area}^{-1}I$.
For a chosen factor $\tau_0$, use
\[
 P_{\mathrm{all}}(Z)=P(Z_{\tau_0})
          \prod_{\tau\ne\tau_0}\prod_{h=1}^{8}
                         \Delta_h(Z_\tau)^{a_h}.               \tag{IF-12s}
\]
It is nonzero and has degree $155$ in each of the eight factors.
The coefficient of the corresponding target tensor in the root
commutator is exactly $J_{U_xf_0}(P_{\mathrm{all}})$.
The last part of the proof therefore supplies a torsion translate
with a nonzero commutator for the actual IF-11 coefficient shapes.
It does not require the CM theta vector to split over embeddings.

For the algebraicity assertion, assume in addition the source
arithmetic hypothesis \texttt{Crv D cc} $\delta$: this consists of
\texttt{IsCurveIdeal} and
\path{IsDiscriminantTwelfthRootOfIdealLattice}. The preceding
analytic nonvanishing assertion requires no such arithmetic
hypothesis. Under \texttt{Crv}, the normalized coefficients of
these finitely many jets are algebraic: use
\texttt{thZero\_coefIn} (lines 236--257) and
\texttt{ta\_thMoved} (lines 311--316), in
\path{VacuumEntireAndNormalizedJetsAlgebraic.lean}; the latter
combines the former with \texttt{coefIn\_translate} (lines 53--85
of \path{CoefInTranslate.lean}). The general
\texttt{TA.transl} (lines 281--290) also transports an already
established \texttt{TA} property under rational translation.
These are properties of the actual normalized CM vector, not a
consequence of the default map \texttt{desc}. A finite number field
contains the finitely many resulting coefficients. Clearing their
denominators gives an integral map with the same nonzero commutator;
no assertion about its Smith factors is needed for this obstruction.

The relation to the actual characteristic kernels is explicit but
one-directional. The source's
\path{GlobalTheory.ColNorms.IndPres.xiInd_expand}, in
\path{Proofs/DivisorProofs/CosetIndicatorThetaTorsionTranslates.lean},
lines 102--114, writes each normalized \texttt{xiInd} kernel as a
finite sum of these exact unitary translates. Its
\texttt{pres\_thInd} (lines 182--248) retains the coefficients,
volume factors, character phases, and rational translation points.
Thus their commutators are the same explicit linear combinations of
(IF-12g). This does not say that a nonzero term survives a particular
combination, nor that every torsion translate is an individual
\texttt{thInd} satisfying all fixed global data. In particular,
(IF-12d) refutes automatic common-root invariance of the genuine
theta-translation family; deciding invariance or period valuations
for a prescribed global automorphic block still requires its actual
combination and global invariant lattices from IF-08.



\section{Three exact arithmetic admissibility checks}
\label{sec:arithmetic-admissibility}

The preceding coefficient constructions are actual constructions.
They must still be distinguished from the additional hypotheses of
a complete step-10 divisibility application. The following three
checks prove incompatibilities of specific supplied specializations;
they are stronger than unsuccessful searches for an input file.

\begin{proposition}[SF-Q6-01: the numerical examples are not global records]
There are no global datum $D:\texttt{D7}$, size datum
$L:\texttt{SizeData}_2(D)$, and nonnegative integers
$\Delta,h_0$ for which \path{ofRecord D L Delta h0} has
the rank and host fields of \texttt{numEx} or \texttt{numBig}.
\end{proposition}
\begin{proof}
The definitions \texttt{hostEx}, \texttt{numEx} and
\texttt{ofRecord}, with the equalities \texttt{hostEx\_q} and
\texttt{hostEx\_d}, occur in
\path{Formalism/LocalToGlobal/Divisibility/NumbersLemmas.lean}.
The definition \texttt{numBig} in the adjacent
\path{NumbersLink.lean} copies these rank and host fields:
\[
 (k,m,t,r,f)=(8,2,15,9,8),\qquad(q,d)=(11,7).
\]
An actual global datum has \texttt{D.host : Host D.R}.
The \texttt{qd\_big} field of \texttt{Host}, in
\path{Formalism/LocalRingCardinality/DataHost.lean}, requires
$4tr+1<q$ and $4tr+1<d$. These become $541<11$ and
$541<7$, a contradiction. The definition \texttt{ofRecord}
copies these global fields, so none of the other choices repairs
them. This does not invalidate the numerical inequalities proved
by \texttt{numFactsBig} on its declared numerical domain.
\end{proof}

The supplied \texttt{Examples}\(_2\)\texttt{.H0} in
\path{Formalism/LocalRingCardinality/SpanPowerMembership/ExampleLayerExistsGenerators.lean}
has the permissible host values $(q,d)=(571,599)$. That is a full
host, not a complete global row-and-period datum. In particular,
the later level-prime congruences of \texttt{Pre} are additional
conditions; the host alone is not a substitute for them.

\begin{proposition}[SF-Q6-02: the inhabited small-weight jets fail the size condition]
Suppose a global datum has $k=8$, a nonempty top set
\texttt{host.TT}, and last Harish--Chandra coordinate
$u_7=-1/2$ at every top embedding. Then it admits no
$\texttt{SizeData}_2$. In particular this applies to the
specific constructed witness of IF-11.
\end{proposition}
\begin{proof}
Choose $t\in\texttt{host.TT}$ and the corresponding top embedding,
using \texttt{D.eTau} and \texttt{tau1}. The field
\texttt{SizeData}\(_2\)\texttt{.top\_u}, in
\path{Formalism/LocalRingCardinality/} followed by
\texttt{DataParams}\(_2\)\texttt{.lean}, lines 10--101,
would give
\[
 -\tfrac12=-b_T(t)-\tfrac72,\qquad b_T(t)=-3.
\]
The \texttt{bT\_ge} field instead gives $b_T(t)\ge b'>0$.
Here the scale definitions are
$b=\exp(x^\varepsilon)$, $a=xb$ and
$b'=a/\exp(x^{\varepsilon/3})$, with $x>1$, so the strict
positivity is immediate. This is a contradiction. In the selected
IF-11 construction the top set is $\{0\}$ and the endpoint is
$-1/2$. The claim concerns that traced witness; the unrefined
existential theorem does not require every possible witness to
have those parameters.
\end{proof}

\begin{proposition}[SF-Q6-03: the asymptotic rank window does not repair these ranks]
For $(k,m,t,r,f)=(8,2,15,9,8)$, no scale satisfies
\path{GlobalTheory.EX.RankVal}. More generally,
$m=2$ together with \texttt{RankVal} implies $x<81$.
\end{proposition}
\begin{proof}
The last condition of \texttt{RankVal}, in
\path{Development/AbelianVarietyTheory/Definitions/GlobalTwisting/PairRowIffQdRows.lean},
lines 17--32, is the equality of nonnegative integers
\[
 r+t=8fr+4+\texttt{TOv}(s)\,\texttt{nWv}(R,s).
\]
The fixed ranks give $24\ge580$, which is impossible.
For the second assertion, the same definition requires
\[
 \sqrt{x\delta_1}\le m,\qquad
 \delta_1=x^{-\varepsilon/8},\qquad
 x>1,\quad 0<\varepsilon<1/52.
\]
Since $1-\varepsilon/8>1/2$,
$x\delta_1=x^{1-\varepsilon/8}\ge\sqrt{x}$, hence
$m\ge x^{1/4}$. If $x\ge81$, then $m\ge3$,
contradicting $m=2$. The supplied general implication is
\texttt{NodeGap.le\_m}, in
\path{Development/AbelianVarietyTheory/Proofs/IwasawaDecomposition/EventuallyNodeGapGeRankPow.lean},
lines 20--35. The short argument here records its specified
two-step consequence and the separate fixed-rank contradiction.
\end{proof}

The constructor \path{exists_sizeData_explicit} in
\path{Proofs/GlobalInterpolation/ExistsExplicitParameters.lean}
(relative to \path{Development/AbelianVarietyTheory/}) assumes
\texttt{RankVal}; the assembled route in
\path{Proofs/IwasawaDecomposition/SizeParametersNonempty.lean}
uses \texttt{Ctx.ok.rank}. It therefore cannot supply the
missing size datum for IF-11's fixed ranks. An unbounded-scale
family with that rank window cannot keep $m=2$ either.
The abstract types \texttt{SizeData}\(_2\) and
\texttt{DivInputs} do not impose \texttt{RankVal} universally.
A direct construction with a different host or weights and all
of its bounds verified remains a distinct approach.

\paragraph{Arithmetic and geometric meaning.}
These restrictions precede the intended arithmetic application.
The numerical polynomial representations still define exterior and
$\lambda$-functors over $\operatorname{Spec}\mathbb Z$, and the
IF-11 family still supplies actual CM-normalized algebraic jets.
The excluded host and size values do not enter the abelian
$C+I$ theorem or the corrected geometric height theorem. They
are additional requirements of the intended arithmetic lattice
construction, with its local denominators. Neither a rational
$K_0$ character, an $\ell$-adic weight, nor a motivic local
$L$-factor supplies these missing numerical equalities or
inequalities. Their relationship to cote 16 scans 014--018 is
therefore through the explicitly specified coefficient and
operator models, not through an identification of those
arithmetic constraints with a Lefschetz conjecture.


% Source-level normalization and conditional local formula.
% GPT-6 Astra Pro; reasoning effort not exposed. 2026-11-21.
\subsection{SF-Q6-04: the actual stabilizer denominator and a full local expression}

\paragraph{Statement SF-Q6-04(a).}
Let $E$ be the host CM field and let $R$ be the rank datum used by
the genuine column construction.  For the denominator
\path{GlobalTheory.Mem.Col2.DGam E R} used in that construction,
every prime divisor satisfies
\[
 p\le [E:\mathbf Q]N_m+1.
\]
For $R=R_0$, this bound is $p\le641$.  In particular its
$863$-adic valuation is zero.

\begin{proof}
The source \path{PredMemGamRanks.lean} defines
\texttt{gamRanks R} as
\[
 \{N_j:1\le j\le m\}\cup
 \{\operatorname{rank}X',\operatorname{rank}Y_{m-1},n_V\}.
\]
Every rank in this set is at most $N_m$; the complete elementary
proof is also supplied as \path{le_N_of_mem_gamRanks} in
\path{DGamDvdOfMinkMult.lean}.
For each such rank $n$, \texttt{autOrd E n} is the set of orders
of the finite unitary automorphism groups of unimodular lattices
in $E^n$.  The source proves that the resulting set of orders is
finite and defines \texttt{DGam} to be their least common multiple.
It proves its positivity; the zero alternative in its definition
does not occur for the positive host involution.

Suppose a prime $p$ divides one of these finite group orders.
Cauchy's theorem gives an element $g$ of order $p$ in
$\mathrm{GL}_n(E)$.  Restriction of scalars realizes $g$ faithfully
on a rational vector space of dimension $[E:\mathbf Q]n$.
Its minimal polynomial divides $T^p-1$.  As $g\ne1$ and
$\mathbf Q$ has characteristic zero, the cyclotomic factor
$\Phi_p(T)$ divides that minimal polynomial.  Consequently
\[
 p-1=\deg\Phi_p\le [E:\mathbf Q]n\le[E:\mathbf Q]N_m.
\]
Every prime divisor of a finite least common multiple divides
one of its entries, which proves the assertion.  For $R_0$,
$[E:\mathbf Q]=2f=16$, $N_m=40$ and $16\cdot40+1=641$.

This is the actual denominator in the column chain:
\texttt{chainAt} sets \texttt{DG := DGam X.E R}.
The proof of \texttt{ColIn.Ok.chIn\_ok} uses this same denominator,
including the free initial stabilizer and all subsequent decorated
lattice stabilizers.  Thus the conclusion does not replace an
arbitrary denominator by a smaller one without justification.
\end{proof}

\paragraph{Statement SF-Q6-04(b): full expression at an inert prime.}
Let $D$ have the fixed ranks $(k,m,t,f)=(8,2,15,8)$ and let
$L:\texttt{SizeData}_2(D)$.  Fix a word $u$ and a coefficient
place $w$ above a rational prime $p>641$.  Assume the following
\emph{ties to the actual data}:
\begin{enumerate}
\item $p$ is neither a type prime $g_i$ nor the level prime $T_u$;
      \path{IsResidueClassZero D p}, $h_P(D,p)=0$, and
      $p\le X''$.
\item $L.F$ has the defining values \texttt{Fval}, $L.E$ has the
      defining values \texttt{Eval}, and the actual
      \texttt{ParametersFill} equality for the gain holds at $w$.
\item The chain denominator is the \texttt{DGam} of part (a),
      and the error constant is $c_{\mathrm{dig}}=3$.
\item At every CM embedding,
      $e_1=M+18$ and $e_2=M+54$.
\end{enumerate}
Then the full source expression, including the error term, is
\[
 \boxed{\quad
 \operatorname{sg}(u,w)
   =\frac{96M+3456}{p^2-1}
       -76800\log_p(a).
 \quad}
\]

\begin{proof}
The defining first branch of \texttt{gainP} is zero when
\texttt{IsResidueClassZero} and $h_P=0$ hold.  Hence
$\mathrm{gainD}_j=0$, whether or not $p$ lies in the drop set.
If $w$ is outside the support, the support axiom already gives
the same conclusion.  In the stated branch of \texttt{Fval},
\[
 F_j=\frac{(k-1)\sum_\tau e_{j,\tau}}{p^2-1}.
\]
The defining branch of \texttt{Eval} is
\[
 E_j=v_p(D_\Gamma)
       +c_{\mathrm{dig}}fN_m^2\log_p(a)
       +\frac{\sum_\tau e_{j,\tau}}{p^2-1}.
\]
Its stabilizer term is zero by (a).  There are $f=8$ CM
embeddings, so
$\sum_{j,\tau}e_{j,\tau}=8(2M+72)=16M+576$.
Subtracting the two errors from the two forced terms gives
\[
 (k-2)\frac{16M+576}{p^2-1}
       -2\cdot3\cdot8\cdot40^2\log_p(a),
\]
which is the asserted expression.
\end{proof}

\paragraph{A strictly positive numerical specialization of this expression.}
Set
\[
 p=863,\qquad x=10^{1000},\qquad
 \varepsilon=\frac1{100},\qquad J=1000,\qquad
 M=\lceil\exp(10^{10})\rceil.
\]
Writing $B=10^{10}$, the source scale gives
$b=\exp(B)$ and $a=10^{1000}\exp(B)$.
Since $M\ge B^2/2$, $\log 863>1$ and $\log10<3$, the displayed
full expression is strictly greater than
\[
 \frac{48B^2}{744768}-76800(B+3000)>0.
\]
The last assertion is an exact rational inequality, independently
checked by the accompanying script.  This uses the full
$F-E$ formula and the actual vanishing stabilizer term; it is not
a conclusion from positivity of $F$ alone.

\paragraph{Remaining realization condition.}
Part (b) is a conditional identity, and the final paragraph proves
positivity of its concrete right hand side.  It is not yet a
certificate that the complete actual global row, twist, integral
lattice and period-valuation data realizing all four ties have
been assembled.  In particular, the source's abstract
\texttt{RowsRestModel} completion of actual columns must not be
substituted for the actual Poincare map required by a
\texttt{RealPack}.

The starting full host $H_0$ has $(q,d)=(571,599)$,
$TT=\varnothing$, $T_2=\{0\}$, $L_w=50$ and split bound $2$.
It must first be relevelled: the complete proof of
\path{exists_relevel_primes_one_mod} preserves its raw host
field and these combinatorial fields while imposing the stronger
\texttt{Pre.level\_cond}.  Choosing its bound above $863$ makes
the new type primes and distinguished level primes exceed $863$.
The concrete scale has
$\lfloor x^{\varepsilon/(5J)}\rfloor
 =\lfloor10^{1/500}\rfloor=1\le2$,
so its split-bound condition does not require changing the raw host.
The remaining construction must be checked independently of the
\texttt{RankVal} route excluded by SF-Q6-03.

\paragraph{Exact source locators and checks.}
Paths are relative to \path{lean/src/}.
\begin{itemize}
\item \path{Development/AbelianVarietyTheory/Definitions/AutomorphicMembership/
PredMemGamRanks.lean}: definitions \texttt{autO}, \texttt{autOrd},
\texttt{gamRanks}, \texttt{gamOrd}, \texttt{DGam} and proofs
\texttt{autO\_finite} through \texttt{DGam\_pos} (lines 322--389);
\texttt{chainAt} (lines 955--966).
\item \path{Development/AbelianVarietyTheory/Proofs/IwasawaDecomposition/
DGamDvdOfMinkMult.lean}: \path{le_N_of_mem_gamRanks}
(lines 13--27); \path{TrueColumnChainsSatisfyConditions.lean}:
\texttt{ColIn.Ok.chIn\_ok} (lines 55--75).
\item \path{Formalism/LocalToGlobal/Columns/
DecorationStabilizerTrivialUnipotent.lean}:
\texttt{primeOrder\_le}, \path{prime_le_of_dvd_card}
(lines 345--389), the same restriction-of-scalars/Cauchy argument.
\item \path{Formalism/LocalRingCardinality/DataForced.lean}:
\texttt{FInert}, \texttt{FvalP} (lines 112--135);
\texttt{DataErr.lean}: \texttt{EatC} (lines 121--123),
\texttt{Eval} (lines 159--170); \path{DataParams$_2$.lean}:
\texttt{gainD}, \texttt{sg} (lines 212--216).
\item \path{Formalism/SharpRestriction/GrowthBounds/
IsResidueClassZero.lean}: \texttt{gainP} (line 355),
\texttt{ParametersFill} (lines 401--411).
\item \path{Development/AbelianVarietyTheory/Proofs/IwasawaDecomposition/
ExistsHostCongruencePrimes.lean}:
\path{exists_relevel_primes_one_mod} (lines 16--102).
\end{itemize}
The full cited stabilizer proofs and the complete defining local
branches were read and reconstructed.  No Lean compilation is claimed.


% Source-level obstruction to a proposed arithmetic completion.
% GPT-6 Astra Pro; reasoning effort not exposed. 2026-11-21.
\subsection{SF-Q6-07: the constant-twist row model cannot realize the forced term}

\paragraph{Statement SF-Q6-07.}
Suppose $D$ is completed from actual columns by
\path{RowsRestModel.model}, as in
\texttt{TrueColsTest.later}.  Let $w$ be a coefficient place above
a prime $p$ and let $1\le j\le m$.  Assume
\[
 p\notin\{g_0,g_1\},\qquad
 \texttt{IsResidueClassZero}(D,p),\qquad h_P(D,p)=0,
\]
and
\[
 k>1,\qquad \sum_{\tau}e_{j,\tau}>0.
\]
There is no $L:\texttt{SizeData}_2(D)$ simultaneously satisfying
$p\le X''$ and \texttt{L.FDef}.

\begin{proof}
The definition \path{RowsRestModel.model} sets
\texttt{muAt w := C.mu0} at every place.  The assembled record
retains this value.  Therefore the \texttt{untwisted} field of
any $L:\texttt{SizeData}_2(D)$ gives $F_j(w)=0$.

On the other hand \texttt{L.FDef}, the stated prime class, and
$p\le X''$ select the residue-class-zero branch of
\texttt{FvalP} and \texttt{FInert}, giving
\[
 F_j(w)=\frac{(k-1)\sum_\tau e_{j,\tau}}{p^2-1}>0.
\]
The denominator is positive because $p$ is prime.  This is a
contradiction.
\end{proof}

\paragraph{Application and scope.}
The proposed large-weight specialization in SF-Q6-04 has positive
$e_1=M+18$ and $e_2=M+54$ and is intended to use $p=863<X''$.
Thus the generic row-model completion cannot supply its arithmetic
record with the required forced-term definition.  This remains true
even though the column theta maps and their compact class-sum
periods inside that completion are actual maps: the auxiliary
row model supplies the wrong place-dependent twist support.

The obstruction does not apply to a different, actual row stage
whose \texttt{MuIs.on} field realizes the chosen twists on a finite
prime set containing $863$.  The datum
\texttt{TwDat.untw B pl} itself exists for every finite set
\texttt{pl}; constructing that datum is only one part of the
required row-stage and period-valuation realization.

\paragraph{Exact source locators.}
All paths are relative to \path{lean/src/}.
\begin{itemize}
\item \path{Development/AbelianVarietyTheory/Definitions/AutomorphicMembership/
ModelLamPDef.lean}, \path{RowsRestModel.model} (lines 130--135);
the actual compact-period formula is retained in its
\texttt{perCol} field (lines 162--166).
\item The same directory, \texttt{ExistsOkPins.lean},
\texttt{TrueColsTest.later} (lines 194--203), which selects
\path{RowsRestModel.model (fits l) 1 1 1};
\path{RowsRestWordColumnUnitBaseTuple.lean}, the assembly
retaining \texttt{muAt} (line 125).
\item \path{Formalism/LocalRingCardinality/DataParams$_2$.lean},
\path{SizeData$_2$.untwisted} (line 108).
\item \path{Formalism/LocalRingCardinality/DataForced.lean},
\texttt{FInert}, \texttt{FvalP}, \texttt{FDef}
(lines 112--135).
\item \path{Development/AbelianVarietyTheory/Definitions/AutomorphicMembership/
CmQuadraticFormEqAlternatingForm.lean}, \texttt{MuIs}
(lines 140--148) and \texttt{TwDat.untw} (lines 196--206).
\end{itemize}
The author read the complete defining model and the complete
\texttt{later} proof, and checked the assembly and forced-term
branches.  This is a human source-level argument; no Lean build is
claimed.


% SF-Q6-05. Human source-level derivation; no Lean compilation claimed.
\subsection{A single-family two-step implication without a common pin}
\label{sec:single-family-m2}

\begin{proposition}[SF-Q6-05: the common pin is unnecessary for one family]
\label{prop:single-family-m2}
Retain the base, column, size, extra-data, and actual
$W:\mathrm{RealPack}(\mathrm{ex})$ hypotheses of the source theorem
\path{famAllQ_of_allTrue}, together with its four named published
theorem inputs (cusp rapid decay, discrete-series cover,
Harish--Chandra modules, and strong approximation for $\mathrm{SU}$),
$\mathrm{AllTrue}(W)$, and $\mathrm{OntoAll}(D)$.
The hypothesis $1\leq r_S$ can be omitted if the conclusion is replaced
by the following assertion: for each fixed word $u$ there are an actual
family $f$ and bottom datum $d$ such that
\[
 \begin{gathered}
 \mathrm{RowOK}(f),\quad \mathrm{LawsX}(f.t),\quad
 \mathrm{LinkZero}(f,d),\quad \mathrm{LinkTail}(f),\\
 \mathrm{LinkDown}(f),\quad \mathrm{LinkOnto}(f),\quad
 \mathrm{LinkComm}(f),\quad \mathrm{RatioQ}(f).
 \end{gathered}
\]
No common central pin for these bottom data is asserted.

Suppose in addition that $m=2$ and that a bottom value for this datum
is nonzero:
\[
 d.\mathrm{per}(\phi,T;y,v)\ne0,
 \qquad T\in\mathrm{ThetaLiftSpaceZero}(W).
\]
Then an actual complex terminal tail of the composed column operator
is nonzero.  More explicitly, writing $t_0$ for the base family member,
there exist $G\in(f.t(t_0)).\mathrm{CT}$ and
$\phi'\in D.\mathrm{what}_2.\mathrm{Mu}(u)$ such that
\[
 (f.t(t_0)).\mathrm{tail}_2
 \bigl((f.t(t_0)).\mathrm{col}_2(\mathrm{col}_0(f,t_0,\phi')),
       G\bigr)\ne0.
\]
The conclusion is a complex nonvanishing statement.  It does not
assert that $G$ is already a prescribed number-field coordinate vector
$\mathrm{topOf}(F)$, nor that its integral preimage has unit output.
\end{proposition}

\begin{proof}
Follow the proof of
\path{TrueGeometryFamilyRowData.famAllQ_of_allTrue}, beginning with
\path{all_true_family_datum}, then
\path{exists_consistent_row_choices},
\path{true_geometry_family_invariance}, and
\path{true_geometry_family_column_data}.
These calls supply a geometric datum $a$, row and column choices, and
the family $f=a.\mathrm{fam}(\mathrm{row},\mathrm{col})$, including
the stated row validity, laws, chart-zero assertion, tail,
commutation, and ratio assertions.  They have no $1\leq r_S$
premise.  Set $d=a.d$.

The exact \texttt{GeoAll.Is.bot} clause in
\texttt{IsStdFr.lean} gives
\[
 \forall h_0,\quad a.d.B.\mathrm{lev}=f.c.B\,0\,r_S\,h_0,
 \qquad a.d.b=f.g.b\,0.
\]
These are precisely the two bottom-identification hypotheses in
\path{link_zero_of_chart_zero}; applying it gives
$\mathrm{LinkZero}(f,d)$.
The unchanged calls to \path{family_link_downward} and
\texttt{family\_link\_onto} give the other two links, using
$\mathrm{AllTrue}(W).\mathrm{pk}$, row validity, and
$\mathrm{OntoAll}(D)$.  The omitted call is exactly
\path{exists_central_bottom_datum}, which had used $1\leq r_S$
to replace $a.d$ by a datum with a common central pin.  Thus every
asserted property is retained without that call.

The supplied theorem
\path{seesawTransportRelations_of_familyLaws} gives
$\mathrm{MainX}(f.t)$ from these laws and the named analytic inputs.
For $m=2$, put $n=0$.  By $\mathrm{LinkOnto}$ there is $G$ with
\[
 \mathrm{down}(f,0,t_0,G)=T.
\]
By $\mathrm{LinkZero}$, $\mathrm{col}_0(f,t_0,\phi)$ belongs to
the required ramified subspace.  Apply $\mathrm{MainX}$ with a
singleton index set and coefficient one.  If
\[
 (f.t(t_0)).\mathrm{tail}_2
 \bigl((f.t(t_0)).\mathrm{up}_1(\mathrm{col}_0(f,t_0,\phi)),G\bigr)=0,
\]
the top hypothesis of $\mathrm{MainX}$ holds.  Its bottom conclusion,
the preceding down identity, and $\mathrm{LinkZero}$ imply
\[
 c_{\mathrm{Bot}}(t_0)\,d.\mathrm{per}(\phi,T;y,v)=0.
\]
Here $c_{\mathrm{Bot}}(t_0)=\prod_{j<m}c_I(t_0,j)$, and each factor
is nonzero by $\mathrm{LawsX}.\mathrm{tup}.c_I$.  This contradicts
the assumed bottom value, so the displayed tail is nonzero.

Finally $\mathrm{LinkComm}$ at $n=1$ supplies a linear equivalence
$S$ and a nonzero scalar $s$ with
\[
 \mathrm{tail}_2(\mathrm{up}_1(\mathrm{col}_0\phi),G)
 =s\,\mathrm{tail}_2(\mathrm{col}_2(\mathrm{col}_0(S\phi)),G).
\]
Taking $\phi'=S\phi$ proves the assertion.  No rationality or
integrality of the arbitrary complex preimage $G$ was used.
\end{proof}

\paragraph{Source and scope.}
The source calls just used are in
\path{Development/AbelianVarietyTheory/Proofs/DecoratedLatticeProofs/}.
The definitions of the links are in
\path{Definitions/LatticeFrames/RatioOKOfRootOfUnity.lean},
lines 28--50; $\mathrm{MainX}$ and $c_{\mathrm{Bot}}$ are in
\path{Definitions/LatticeFrames/ValLastX.lean}, lines 244--252
and 174.  This is a mathematical derivation from the supplied complete
source proofs, not a newly compiled Lean theorem.  It removes the
uniform-pin gate for this conclusion; it retains an actual valid
$\mathrm{RealPack}$, the actual bottom nonzero value, and every
other displayed hypothesis.


% SF-Q6-06. Exact diamond calculation and genuine refined-level period.
\subsection{Diamond characters and the bottom period}
\label{sec:diamond-fourier-anchor}

Fix a word $u$ in an actual geometric column datum.  Put
$A_v=\mathbf F_v^{\times}$, $m_v=|A_v|=T-1$, and
$A=\prod_v A_v$.  On the actual column space $M_0$, diamond operators
at distinct places commute: in the decorated description they rescale
different decoration functionals, and in the adelic description their
lifts are supported at distinct places.  This commutation is part of
the geometric specialization used below; it is not a field of the
bare abstract record \texttt{ColumnsAtT}.  Write
\[
 \rho(a)=\prod_v\rho_T(v,a_v),\qquad
 N_v=\sum_{a\in A_v}\rho_T(v,a),\qquad
 M_{\rm ram}=\bigcap_v\ker(N_v|_{M_0}).
\]

\begin{proposition}[SF-Q6-06(a): the exact finite Fourier test]
\label{prop:diamond-fourier-test}
Over $\mathbf C$, define
\[
 E=\prod_v\left(1-\frac{N_v}{m_v}\right),\qquad
 e_\chi=\frac1{|A|}\sum_{a\in A}\chi(a)^{-1}\rho(a)
 \quad(\chi\in\widehat A).
\]
Then $E$ is the projection onto $M_{\rm ram}$ and
\[
 E=\sum_{\substack{\chi\in\widehat A\\
                  \chi|_{A_v}\ne1\ \text{for every }v}}e_\chi.
\]
For any linear map $P:M_0\to W$, including the genuine
Poincar\'e map after the column dictionary, the following are
equivalent:
\[
 P|_{M_{\rm ram}}\ne0,
 \qquad PE\ne0,
 \qquad Pe_\chi\ne0\text{ for some such }\chi.
\]
For a scalar functional $\ell$ and characters $\eta_v\in\widehat A_v$, if
$\ell\rho_T(v,a)=\eta_v(a)\ell$ for every $v,a$, then
\[
 \ell E=
 \begin{cases}
 \ell,&\eta_v\ne1\text{ for every }v,\\
 0,&\eta_v=1\text{ for some }v.
 \end{cases}
\]
Thus a trivial character at even one place is an actual vanishing
obstruction, provided this eigenfunctional identity has been proved.
\end{proposition}
\begin{proof}
The group law gives $N_v^2=m_vN_v$.  The averages $N_v/m_v$ are
commuting projections onto the corresponding fixed spaces.  Therefore
$E$ is idempotent, $N_vE=0$ for every $v$, and $E$ is the identity
on the intersection of the norm kernels.  Simultaneous decomposition
of this finite abelian group representation into character spaces
gives the asserted sum: $N_v$ acts on the $\chi$-space by $m_v$
if $\chi|_{A_v}=1$, and by zero otherwise.  The equivalences follow
by composing with this direct-sum decomposition.  The same character
sum applied to $\ell N_v$ proves the final formula.
\end{proof}

This test has a denominator-free lattice version.  If the diamonds
preserve a full lattice, then
\[
 Q=\prod_v(m_v-N_v)=\left(\prod_v m_v\right)E
\]
preserves it and has image in its norm kernels.  Consequently a
nonzero value of $PQ$ on a lattice vector supplies an integral
nonzero input, although its value need not be a unit.  One can also
test the scalar functional on the explicitly finite sum
\begin{equation}
 \ell(E x)=\sum_{a\in A}
 \left(\prod_v\left({\bf1}_{a_v=1}-\frac1{m_v}\right)\right)
 \ell(\rho(a)x).
 \label{eq:diamond-finite-test}
\end{equation}

\paragraph{The genuine map is vector-valued.}
In \texttt{PoincareInHH.lean}, lines 144--210,
$P_u:\mathrm{Cu}\to\mathrm{HH}$ is a linear map into an automorphic
Hilbert space.  It is not a scalar period functional.  The scalar
period is obtained after pairing with a row.  In particular the
eigenfunctional hypothesis in the last part of the proposition
cannot be imposed on $P_u$ merely because the row has a level character.

\begin{proposition}[SF-Q6-06(b): the actual diamond translation]
\label{prop:diamond-period-translation}
Let $\delta\in H(\mathbf A_f)$ be a diamond lift of $a\in A_v$
in the source's column dictionary.  Use right translation
$R_x f(g)=f(gx)$, and let $\iota:H\hookrightarrow G$ be the bottom
Bessel embedding.  For the ordinary bottom period with a fixed
row translate $b$,
\[
 \ell_{b,f}(\phi)=B(\phi,R_bf),
\]
one has the exact identity
\begin{equation}
 \ell_{b,f}(\rho_T(v,a)\phi)
   =B(\phi,R_{\iota(\delta)^{-1}b}f).
 \label{eq:diamond-exact-translation}
\end{equation}
If $f$ has right level $(J,\chi)$, with the convention
$f(gk)=\chi(k)f(g)$ for $k\in J$, and
$\kappa=b^{-1}\iota(\delta)b\in J$, then and only using this
additional membership condition the right side simplifies to
$\chi(\kappa)^{-1}\ell_{b,f}(\phi)$.

For a nonidentity diamond, its transported lift is outside the
inherited column level $K_u$.  Equivalently, in the analytic
coordinates of \texttt{ColDict.LevelMatch},
\[
 x_u^{-1}\delta_Hx_u\notin K_2.
\]
Thus the host character $\chi_2$ on $K_2$ supplies no scalar
eigenvalue for this lift by level invariance alone.
\end{proposition}
\begin{proof}
The exact dictionary is
\texttt{FormsTied.rhoT\_eq} and \texttt{ColDict.ColDictDia}:
$\rho_T(v,a)$ corresponds to $R_\delta$ on the column form.
The supplied theorem \path{Core.Mi.per_rTrans} proves
\[
 B(R_\delta\phi,R_{\iota(\delta)}g)=B(\phi,g)
\]
for the actual automorphic Bessel period.  Apply it with
$g=R_{\iota(\delta)^{-1}b}f$; the product order follows from
$R_xR_y=R_{xy}$.  This proves
\eqref{eq:diamond-exact-translation}.  If the stated $\kappa$
belongs to $J$, then
$R_{\iota(\delta)^{-1}b}f=R_bR_{\kappa^{-1}}f
=\chi(\kappa)^{-1}R_bf$, giving the scalar formula.  In the
source the fixed $J$ and $\chi$ are precisely \texttt{PairForm.rJ}
and \texttt{PairForm.rChi}, and $b$ is the actual bottom translate;
neither is removed from this calculation.

For the last claim, let $S$ be the source's \texttt{StepLevel}.
Its local column group at $v$ is
\[
 U_v=\operatorname{Stab}(L_v)\cap\operatorname{Stab}(f_T),
 \qquad f_T\ne0.
\]
The \texttt{IsDia} condition says
$k f_T(\delta x)=f_T(x)$, where the integer $k$ maps to $a$ in
$\mathbf F_v$.  If $\delta\in U_v$, it also gives
$f_T(\delta x)=f_T(x)$.  Transport the additive target of $f_T$
to $\mathbf F_v$ by the supplied additive isomorphism
\texttt{FormsTied.fMe}.  Choose $x$ with nonzero image.
Then $(a-1)f_T(x)=0$ in a field, so $a=1$, a contradiction.
Finally \texttt{ColDict.LevelMatch} identifies this column group
with $H(\mathbf A_f)\cap x_uK_2x_u^{-1}$.  This proves the claimed
nonmembership.
\end{proof}

\paragraph{What the exact formula tests.}
For several places choose the product lift $\delta_a$.
Substituting \eqref{eq:diamond-exact-translation} into
\eqref{eq:diamond-finite-test} tests the actual values
\[
 B(\phi,R_{\iota(\delta_a)^{-1}b}f).
\]
These are a finite orbit of translated row periods, not a host
character evaluated on points known to belong to its domain.
Global rational cosets can collide even when a lift is outside
$K_u$; the preceding local nonmembership assertion does not
exclude such collisions.

\begin{proposition}[SF-Q6-06(c): a genuine period after level refinement]
\label{prop:refined-level-ram-period}
Let $H\subset G$ be a pair of reductive groups over a number field
$F$ whose archimedean groups are compact.  Let
$K_0\subset G(\mathbf A_f)$ be compact open,
$U=H(\mathbf A_f)\cap K_0$, and assume
$H(F)\cap U=\{1\}$.
Let a finite abelian group $A=\prod_vA_v$, with $|A_v|>1$, act
on the column space by lifts $\delta_a\in H(\mathbf A_f)$,
normalizing $U$ and its finite level character.  Assume the orbit
of the identity class in
$H(F)\backslash H(\mathbf A_f)/U$ is free.
Let $V_H,V_G$ be the actual finite-dimensional archimedean coefficient
representations, with a nonzero $H(F_\infty)$-invariant bilinear
pairing $\beta:V_H\times V_G\to\mathbf C$.

Then there is a compact open $K'\subset K_0$ and an actual algebraic
modular row form $f$ of coefficient type $V_G$ and level $K'$,
and a column $\phi\in M_{\rm ram}$ of the prescribed $U$ level,
whose ordinary definite period is nonzero.  The construction retains
the given diamond norm kernels; it does not replace them by an
arbitrary chosen subspace.
\end{proposition}
\begin{proof}
Choose $v\in V_H$, $w\in V_G$ with $\beta(v,w)\ne0$.
Let $T_0$ be the column supported on the identity class
$H(F)U$, with value $v$ at the identity and the prescribed right
$U$ character.  This is well-defined because $H(F)\cap U=\{1\}$.
The support of $R_{\delta_a}T_0$ is
$H(F)U\delta_a^{-1}$.

Put $\Gamma=G(F)\subset G(\mathbf A_f)$ and
$S=\bigcup_{a\ne1}U\delta_a^{-1}$.
The set $S$ is compact.  Freeness of the global class orbit gives
$S\cap H(F)=\varnothing$.  Since $H$ is defined over $F$,
$H(\mathbf A_f)\cap G(F)=H(F)$, so $S\cap\Gamma=\varnothing$.
The arithmetic group $\Gamma$ meets every compact subset of
$G(\mathbf A_f)$ in a finite set: the corresponding subset of
$G(F_\infty)\times G(\mathbf A_f)$ is compact, and the diagonal
$G(F)$ is discrete and closed.

Consequently one can choose a compact open $K'\subset K_0$ so that
\[
 \Gamma\cap SK'=\varnothing,
 \qquad \Gamma\cap UK'=\{1\},
 \qquad \Gamma\cap K'=\{1\}.
\]
Indeed $SK_0$ and $UK_0=K_0$ meet $\Gamma$ in finite sets;
each unwanted rational point is outside respectively $S$ or $U$,
which are compact closed sets.  A sufficiently small identity
neighborhood separates all of these finitely many points, and it
contains a compact open subgroup.  Shrink $K'$ further so that
the column character is trivial on
$U'=H(\mathbf A_f)\cap K'$.  This is possible because its kernel
is open in $H(\mathbf A_f)$ and the topology on that subgroup is
the subspace topology.  If a compatible finite row character on
$K_0$ has already been prescribed, $K'$ can also be placed in its
kernel.

Define the row form $f$ to be supported on $\Gamma K'$, with value
$w$ at the identity and the usual coefficient covariance under
$G(F)$.  This is an actual algebraic modular form: the two
descriptions of any point in its support are unique modulo
$\Gamma\cap K'=\{1\}$, and it is locally constant in the finite
variable.  On the identity column support $U$, it is supported
exactly on $U'$.  The invariant pairing there is the constant
$\beta(v,w)$.  Thus, for the fixed positive Haar normalization,
\[
 B(T_0,f)=\operatorname{vol}(U')\,\beta(v,w)\ne0.
\]
For every $a\ne1$, the support separation
$\Gamma\cap U\delta_a^{-1}K'=\varnothing$ gives
$B(R_{\delta_a}T_0,f)=0$.
Apply the genuine ram projection $E$ from part (a).  Its identity
coefficient is $\prod_v(1-|A_v|^{-1})$, which is positive, so
\[
 B(ET_0,f)=
 \left(\prod_v(1-|A_v|^{-1})\right)
 \operatorname{vol}(U')\,\beta(v,w)\ne0.
\]
The column $ET_0$ belongs to every required norm kernel.
\end{proof}

\paragraph{Source freeness and stabilizers.}
The actual geometric freeness hypothesis has a supplied theorem:
\path{SiegelMassRatios/DiamondFree.diaFreeCols_holds}, under
the source's field and prime validity hypotheses and
$2fk<g_0-1$, proves freeness on the global decorated class set.
For $k=f=8$ this numerical hypothesis is $128<g_0-1$.
This is stronger than the local nonmembership in part (b).
If a scalar permutation orbit is instead $A/S$, its occurring
characters are exactly $S^\perp$.  A sufficient condition for the
ram projection of its identity point to be nonzero is
\[
 \sum_v\frac{|S\cap A_v|}{|A_v|}<1;
 \qquad\text{in particular }\quad
 \sum_v\frac{|S|}{|A_v|}<1\ \text{suffices}.
\]
To check this, among the $|A|/|S|$ characters trivial on $S$, the
proportion also trivial on $A_v$ is
$|S\cap A_v|/|A_v|$.  The union bound leaves a character nontrivial
on every factor.  This scalar-permutation statement must not be
applied without adjustment to a coefficient fibre carrying a
nontrivial stabilizer representation.

\paragraph{The prescribed-level boundary.}
Part (c) is a genuine geometric nonzero-period construction at a
refined level for a definite ambient group.  The actual Hida ambient
at $m=2$ has signature $(k,1)$, not $(k+1,0)$: the source theorem
\path{InHHDef.card_epsS_false} gives $r_S+1=1$ negative
direction.  Its bottom stabilizer has compact archimedean group,
but the ambient group does not.  Thus the finite-support proof in
part (c) does not directly apply to that ambient group even before
the level condition is considered.  An actual nonzero Poincar\'e
series would additionally require an analytic separation or
nonvanishing argument for the prescribed archimedean type.

There is also a separate prescribed-level condition.  The exact
source definition is
\[
 K_2=K_{\rm lattice}\cap
 K_{2,\mathrm{cong}}(g_0g_1,e_0,z_0,\mathrm{bw}).
\]
\texttt{HidaSetting.L\_eq}, \texttt{g\_eq}, and
\texttt{e\_eq} fix the lattice, modulus, and characteristic
projector; they do not contain a freely variable auxiliary compact
open subgroup.  A smaller $K'$ need not be another allowable $K_2$,
and averaging the constructed row back to the original level can
annihilate its ramified period.  One must therefore either realize
the refined level by an actual admissible datum and its theta-chain
comparison, or evaluate the finite test
\eqref{eq:diamond-finite-test} in the existing row space.  Neither
the nonzero column projection nor complex surjectivity of a separate
theta map supplies that last identification by itself.

\paragraph{Exact source locators.}
The norm kernels are
\path{Formalism/LocalRingCardinality/DataKernels.lean},
lines 146--150.  The local level and diamond formulas are
\path{HermitianForms/IsLocBasis.lean}, lines 84--100 and 137--144;
the dictionary is \path{HermitianForms/FormsTied.lean},
lines 133--141, and
\path{SiegelLatticeTheory/FlagStabilizerTensors/ColDict.lean},
lines 143--157 and 190--196.  The proved period covariance is
\path{Proofs/DecoratedLatticeProofs/BesselPeriodRightTranslationInvariant.lean},
\path{Core.Mi.per_rTrans}.  The fixed row and translate occur in
\path{Definitions/LatticeFrames/PairForm.lean}, lines 41--47 and
75--78.  The prescribed $K_2$ and projector occur in
\path{SiegelLatticeTheory/HaarTransversals/HidaSetting.lean},
lines 128--141, 189, 199, 258--259, and 350--362.


% GPT-6 Astra Pro; reasoning effort not exposed. SIGN_NEXT, 2026-11-21.
\subsection{SF-Q6-08: a pointwise inert estimate with the actual input gates}

\paragraph{Status and statement.}
The following is a conditional theorem about the \emph{actual} column
chain and terminal period of an actual datum $D:\texttt{D7}$.
It does not assert that the numerical candidate below supplies all its
hypotheses.  Its purpose is to remove the irrelevant global size and
split-bound hypotheses from one fixed inert-prime estimate.

Fix a word $u$, a coefficient place $w$ of residue characteristic
$p\ge5$, and a used twist $\mu$.  Assume
\path{IsResidueClassZero D p}, $h_P(D,p)=0$, and
\path{LevelFree Dx u mu w}.  In particular $p$ is distinct from the
column level prime and both type primes, and does not belong to the
exceptional twist set $\texttt{Dx.Qp mu}$.
Use the actual twisted forms $T:\texttt{TwForms D}$ and exact kernel
entry data $E:\texttt{OrbitPlaces.Ent D}$.  Write
$s_{i,z}=|\nu_{i,z}|$ for the sizes of the partitions in $E$, and
$e_w$ for the ramification index normalizing the source's
$\texttt{ur D w}(b)=e_wb$.  The required gates are:
\begin{enumerate}
\item For every step $i<m$ and every actual pair of column classes,
      the series $E.\texttt{Th}$ at this fixed $w$ satisfies
      \path{ThetaExpansionRealPlaces}.  This includes the integral
      remainder and the logarithmic theta factorization with conditions
      B1--B4; coefficient algebraicity alone is insufficient.
\item The actual transfer statement \texttt{KForm E Dx T 1} holds at
      this word, twist and place, for all steps and all bounds.  This
      transfers the Schur coefficient estimate to the actual column
      lattices, including class sums and stabilizer denominators.
\item The actual terminal pairing satisfies \texttt{PairInt D Dx T}
      at this word, twist and place, and \texttt{DGamDef D} holds.
      The row vector lies in the specified lattice
      $\widetilde M(\mu)$ and space $V_0(\mu)$.
\item $N_iN_{i+1}$ is even for each step, and a real number $A$ satisfies
      $p\le A$ and $s_{i,z}\le A$ for all $i,z$.
\end{enumerate}
Let $d_w=\texttt{vDG D u w}$, with the source normalization.
For every input $\phi\in\texttt{MramL D u}$ and every indicated row
vector $F$, the actual period is the image of some $y$ in the actual
coefficient field.  Either $y=0$, or
\begin{equation}\label{eq:SF08-pointwise}
 \frac{v_w(y)}{e_w}\ \ge\
 \sum_{i<m}\left(
   \frac{\sum_z s_{i,z}}{p^2-1}
       -fN_iN_{i+1}(\log_p A+1)\right)-m d_w.
\end{equation}
The same conclusion applies to any specified input lattice locally
contained in \texttt{MramL}, using that containment explicitly.
No nonvanishing of the period is asserted.

\begin{proof}
The residue and $h_P=0$ assumptions give $n_P=1$.
For each $z$, the one-orbit local theorem
\texttt{Loc1.ThmOne 1} applies.  Its entire
\texttt{Loc1.HidaSetting A} hypothesis is
$p\le A$ and $|\nu_{i,z}|\le A$.
In particular it does not require a logarithmic lower bound on the
matrix rank.  Its estimate is
\[
 \frac{s_{i,z}}{p^2-1}
       -N_iN_{i+1}(\log_p A+1).
\]
The closed source theorem
\path{split_bound_eleven_one} supplies this constant-one local
bound.  The theorem \path{orbitsAddAtConjugate_ent} adds these
bounds over the $f$ real places using the stated theta expansion;
thus it gives \texttt{EntGe} with
\[
 B_i=\frac{\sum_zs_{i,z}}{p^2-1}
                   -fN_iN_{i+1}(\log_p A+1).
\]
The actual \texttt{KForm} gate with $\texttt{sp}=1$ gives
\texttt{KGe} with $B_0$ at the first step and $B_i-d_w$ at
all subsequent steps.  The proof of \path{OrbitPlaces.compose}
iterates these actual lattice estimates.  Its twisted branch uses
\path{base_map_mem_local_lattice}; consequently the twist
and its integral base-lattice map have not been discarded.
The sum after $m$ steps is $\sum_i B_i-(m-1)d_w$.
The theorem \texttt{nuIn\_mem\_Sram} places the indicated input in
that initial local lattice.  Finally \texttt{period\_ge}, using
\texttt{PairInt}, represents the actual terminal period over the
coefficient field and subtracts one further $d_w$.  Its elementary
valuation proof is the multiplication identity
$v_w(ay)=v_w(a)+v_w(y)$ for nonzero factors, with separate zero cases.
This proves \eqref{eq:SF08-pointwise}, including the ramification
factor and every stabilizer loss.
\end{proof}

\paragraph{Why the stripped hypotheses are legitimate.}
The packaged \texttt{Shape.rankA} condition would impose
$2\log_5(a^2)+2\le k-1$.  It fails for the candidate below, but
it is not used by the one-orbit proof just given.  Similarly the
packaged true-column theorem assumes $3\le\texttt{splitBound}$.
In its fixed inert-prime branch that assumption is used to obtain
$p\ge5$; here $p=863$ supplies this directly.
The actual standard-series theorem
\path{thetaExpansionRealPlaces_stdSer_of_adm} has hypotheses
\texttt{HostAt}, $5\le p$, the residue-class and $h_P=0$ conditions,
admissible lattice pairs, standard bases, and the stated normalized
combination identity.  It has no \texttt{rankA} or
$3\le\texttt{splitBound}$ hypothesis.  The fixed-prime part of the
adapted-kernel proof can therefore be repeated using that theorem
and the actual identification of $E.\texttt{Th}$ with its standard
series.  This observation reduces an input gate; it does not
manufacture the required actual identification or integral bases.

\paragraph{The explicit two-step partitions, including a variable gap.}
Take $k=8$, $m=2$, $t=15$, $f=8$, so
$(N_0,N_1,N_2)=(8,24,40)$.  At every embedding take
\[
 u=\tfrac12(113,45,37,29,21,13,5,-1),\qquad
 (b_1,b_2)=(M,M+G),\qquad
 (e_1,e_2)=(M+18,M+G+18).
\]
The source formula \texttt{nuRow}, with its ordinary derivative
normalization, gives
\begin{align*}
 \nu_0={}&(M+59,M+26,M+23,M+20,M+17,M+14,M+11,M+9),\\
 \nu_1={}&(M+G+43,M+G+10,M+G+7,M+G+4,\\
          &\hspace{8mm}M+G+1,M+G-2,M+G-5,M+G-7,(G-16)^{16}).
\end{align*}
For integers $M\ge0$ and $G\ge16$ these are partitions, and their sizes are exactly
\[
 |\nu_0|=8M+179,\qquad |\nu_1|=8M+24G-205.
\]
These identities follow by substituting into the first eight rows
of \texttt{srcPar} and its sixteen added rows; in the latter rows
the two half-integer rank shifts cancel, leaving $G-16$.
Thus $\sum_{i,z}s_{i,z}=128M+192G-208$ and
$f(N_0N_1+N_1N_2)=9216$.
For $p=863$, SF-Q6-04(a) gives $d_w=0$ for the genuine denominator.
With $A=a^2$, \eqref{eq:SF08-pointwise} becomes
\begin{equation}\label{eq:SF08-explicit}
 R_G=\frac{128M+192G-208}{744768}
                      -9216(2\log_{863}a+1).
\end{equation}

If in addition the actual F/E/gain ties of SF-Q6-04(b) are supplied,
their full expression for this variable gap is
\[
 \operatorname{sg}_G
     =\frac{96M+48G+1728}{744768}-76800\log_{863}a.
\]
For $M\ge0$, $G\ge16$ and $\log_{863}a\ge1$ one has
\[
 R_G-\operatorname{sg}_G
 =\frac{32M+144G-1936}{744768}
                   +58368\log_{863}a-9216>0.
\]
Consequently the pointwise theorem implies this full weaker bound
when all its actual arithmetic ties hold.

For $B=10^{10}$, $M=\lceil e^B\rceil$,
$a=10^{1000}e^B$, and any $16\le G\le M$, the size hypotheses with
$A=a^2$ hold: $M\le2e^B$, $M\ge179$, and both partition sizes
are at most $32M$, whereas $a^2>32M$ and $a^2>863$.
For $G\ge36$, $\operatorname{sg}_G$ is at least the positive
expression of SF-Q6-04(c), whose explicit rational lower bound is
$22020927106278400000/3879$.
The old illustrative gap $G=36$ is sufficient for this pointwise
calculation; it is \emph{not} asserted to satisfy the separate
\texttt{NodeGap} hypotheses of theta-chain surjectivity.
A large gap, for example $G=10^9$, retains this local positivity.
For the exact \texttt{NodeGap} gate, $\operatorname{rank}X'=39$,
so its existential auxiliary dimension can be chosen as $N=40$.
The required bound is $15+2\cdot40^3+1=128016<10^9$.
Thus the large gap satisfies this particular gate.  This does not
supply the other geometric and analytic hypotheses of surjectivity.
The exact definition is
\path{Development/AbelianVarietyTheory/BlockTheory/}
\path{TwistedDiscreteSeries/ThetaSourceAction.lean}, lines 209--212,
with a run copy in the critical reviewer's \texttt{pu\_source\_cache}.

\paragraph{Remaining actual input gates.}
An actual global specialization still requires the relevelled host
and its actual column datum; the prescribed kernel and normalized
standard-series identification; integral frames supplying
\texttt{KForm}; the actual row spaces, twists, and integral
\texttt{PairInt}; and, to apply the complete DivInputs theorem,
its cutoff, row, gain, F/E, input-lattice and size-data ties.
A coefficient-field realization of a finite kernel does not prove
these theta-factorization or integral-pairing statements.
The constant-twist model is separately rejected by SF-Q6-07.
The theorem above therefore removes a real obstruction in the
\emph{packaged proof}, but is not a certified nonzero global period
or an unconditional global divisibility example.

\paragraph{Exact archive and run locators.}
All paths below are below \path{lean/src/} in the supplied archive;
run copies are below \path{lean_selected/} or
\path{SIGN_NEXT/work/singular_fibres/arithmetic_review/source_cache/}.
\begin{itemize}
\item \path{Formalism/SharpRestriction/Divisibility/}
\path{ProductZerothFirstParts.lean}, lines 263--272:
\texttt{Loc1.Inp}, \texttt{HidaSetting}, \texttt{ThmOne}.
\item \path{JetPolynomialLogs/DworkPowerSeriesWeights.lean},
\texttt{thmOne\_of}, lines 235--266; and
\path{JetPolynomialLogs/LocalBoundElevenComparison.lean},
\path{split_bound_eleven_one}, lines 64--68, below that
Divisibility directory.
\item \path{Formalism/SharpRestriction/Divisibility/ChainOperators/OrbitEntryBounds.lean},
\path{orbitsAddAtConjugate_ent}, lines 23--24.
\item \path{Formalism/SharpRestriction/Divisibility/OrbitPlaces.lean},
\texttt{Ent}, lines 34--52; \texttt{KForm}, lines 76--79;
\texttt{compose}, lines 806--814; $n_P=1$, lines 1129--1137.
\item \path{Formalism/SharpRestriction/GrowthBounds/}
\path{WeilJacquetShalikaLaw.lean}, \texttt{PairInt}, lines 383--395;
\texttt{period\_ge}, lines 419--442;
\texttt{nuIn\_mem\_Sram}, beginning line 444.
\item \path{Development/AbelianVarietyTheory/Proofs/SatakeParameters/}
\path{B0StandardSeriesAdmissible.lean}, lines 26--39;
the full fixed-prime proof in
\path{FrobeniusConjugationConditionsAdaptedKernelDatum.lean};
and the packaged wrapper \path{AtOfTrueColumns.lean}.
\item \path{Formalism/LocalRingCardinality/DataDIVDefs.lean},
\texttt{srcPar} and \texttt{nuRow}, lines 60--73.
\end{itemize}
The displayed proof rechecks the complete valuation-composition,
terminal-pairing, and one-orbit wrapper proofs.  Deeper analytic
prerequisites of the imported theta expansion and one-letter estimate
are used as source theorems; no complete Lean build or independent
reproof of all transitive dependencies is claimed.


\section{An explicit bounded formula for every exterior-order index}
\label{sec:bounded-indices}

The tables of IF-03 can be extended without constructing a matrix on
the full exterior space and without an unbounded word-closure loop.
The following gives one finite integer matrix, of explicitly bounded
size, whose Smith factors are exactly all the desired order indices.
It is a formula for every rank, not a claim that a small-rank table
determines an unspecified theta period.

\begin{theorem}[SN-15: all-rank index matrices]
Let $V_n=\Lambda^*\mathbb Z^{2n}$, with its standard integral
symplectic basis, let $L$ be multiplication by its symplectic form,
and let $D$ be ordinary contraction as in IF-01. Put
\[
 R_n=\mathbb Q[L,D],\qquad
 A_n=R_n\cap\operatorname{End}_{\mathbb Z}(V_n),\qquad
 O_n=\mathbb Z[L,D],\qquad
 O_n^\pi=\mathbb Z[L,D,\pi_0,\ldots,\pi_{2n}],
\]
where $\pi_j$ is the exterior-degree projector. Let
\[
 \mathcal R_n=\{(v,k,l,t):0\le v\le n,\ 0\le k,l\le v,
 \ \max(0,k+l-v)\le t\le\min(k,l)\},
\]
and index columns by $(a,c,b)$ with $0\le a,c\le n$ and
$0\le b\le2n$. For $s=l-c$, define integer matrices $M_n$ and
$M_n^\pi$ by
\[
 (M_n)_{(v,k,l,t),(a,c,b)}=
 \begin{cases}
 a!c!\binom ts\,(v-2s)^b,&0\le s\le t,\quad k=s+a,\\
 0,&\text{otherwise},
 \end{cases}                                      \tag{SN-15a}
\]
and
\[
 (M_n^\pi)_{(v,k,l,t),(a,c,b)}=
 \begin{cases}
 a!c!\binom ts\,\mathbf1_{b=n-v+2l},
       &0\le s\le t,\quad k=s+a,\\
 0,&\text{otherwise}.
 \end{cases}                                      \tag{SN-15b}
\]
Here $0^0=1$. Each matrix has $\binom{n+4}{4}$ rows and
$(n+1)^2(2n+1)$ columns, and rank
$r_n=(n+1)(n+2)(2n+3)/6$.

The nonzero Smith factors of $M_n$ are precisely the invariant
factors of $A_n/O_n$ (including unit factors), and those of
$M_n^\pi$ are precisely the invariant factors of $A_n/O_n^\pi$.
In particular, if $\delta_j(M)$ is the nonnegative gcd of all
$j$-rowed minors, with $\delta_0=1$, then
\[
 [A_n:O_n]=\delta_{r_n}(M_n),\qquad
 [A_n:O_n^\pi]=\delta_{r_n}(M_n^\pi),\qquad
 d_j=\delta_j/\delta_{j-1}.                         \tag{SN-15c}
\]
These are explicit finite binomial-and-determinant formulas for
every rank. All their prime divisors are at most $n$.
\end{theorem}

\begin{proof}
For each symplectic pair choose one of its four exterior states
$1,e_i,f_i,e_i\wedge f_i$. Fix the pairs in singleton states
$e_i$ or $f_i$. On the remaining $v$ pairs, the basis is indexed
by subsets $S\subset\{1,\ldots,v\}$: a member of $S$ is in the
full state and a nonmember in the empty state. The actions are
\[
 L[S]=\sum_{i\notin S}[S\cup\{i\}],\qquad
 D[S]=\sum_{i\in S}[S\setminus\{i\}].             \tag{SN-15d}
\]
The signs are positive because the pair state has degree two;
the two contraction signs cancel. This is an integral direct-sum
decomposition, and every $v=0,\ldots,n$ occurs. Repeated copies
have identical operator matrices. We may therefore test matrix
integrality on one copy of each of these Boolean modules.

Set $H=DL-LD$. On $[S]$ it acts by $v-2|S|$; equivalently it
acts on exterior degree $j$ by $n-j$. Directly from (SN-15d),
\[
 [H,L]=-2L,\qquad [H,D]=2D,\qquad DL=LD+H.
\]
These integral commutation rules reorder every word in $L,D,H$
as an integral sum of $L^aH^bD^c$. For a detailed termination
argument, order the three letter types as $L,H,D$ and move each
out-of-order adjacent pair past one another. The exchanged term
has fewer inversions of that word length, while the commutator
terms have smaller word length. Induction on length and then
inversions terminates and introduces only integer coefficients.
Since $H\in O_n$, this gives a spanning set for $O_n$.
The identities $L^{n+1}=D^{n+1}=0$ and
\[
 \prod_{j=-n}^{n}(H-j)=0
\]
reduce it to $0\le a,c\le n$, $0\le b\le2n$; the polynomial
in $H$ is monic with integer coefficients. Thus the columns in
(SN-15a) will span the whole order, not just a finite selection
of its short words.

Fix a source subset $S$ of size $l$ and a target subset $T$ of
size $k$, with $|S\cap T|=t$. The intermediate subset after
$D^c$ has size $s=l-c$. For a nonzero contribution it must be a
subset of $S\cap T$, and $k=s+a$. There are $\binom ts$ such
intermediate subsets. Ordinary removal contributes $c!$, ordinary
addition contributes $a!$, and the intermediate action of $H^b$
contributes $(v-2s)^b$. This proves (SN-15a).

For the order with projectors, move every degree projector to
the right using $\pi_jL=L\pi_{j-2}$ and
$\pi_jD=D\pi_{j+2}$, with out-of-range projectors zero.
Orthogonality collapses their products. Reorder the remaining
word as above; $H^bD^c\pi_j=(n-j+2c)^bD^c\pi_j$.
Consequently $L^aD^c\pi_j$, $0\le a,c\le n$,
$0\le j\le2n$, span $O_n^\pi$ integrally. The source subset
has exterior degree $n-v+2l$, including the fixed singletons.
The preceding count proves (SN-15b).

All entries of a matrix from $R_n$ on a given Boolean module
depend only on $(|T|,|S|,|T\cap S|)$: permutations of the $v$
pairs commute with both operators and act transitively on pairs
of subsets having these three invariants. Every admissible
quadruple in $\mathcal R_n$ actually occurs. Recording one entry
per quadruple is therefore injective on $R_n$, and an element
of $R_n$ has an integral full exterior matrix if and only if
these recorded entries are integral. The count of quadruples
for a fixed $v$ is $\binom{v+3}{3}$: they are the four
nonnegative sizes of $S\cap T$, $T\setminus S$,
$S\setminus T$, and the complement. Summing gives the claimed
row count.

Thus in the recorded coordinate space, $A_n$ is exactly the
saturation of the column lattice of $M_n$. The same holds for
$M_n^\pi$: its rational span is still $R_n$, because IF-02
puts the degree projectors in $A_n\subset R_n$. IF-02 also
gives the stated rank. Smith normal form of an integer matrix
identifies the quotient of its saturated column span by its
column lattice with the direct sum of the cyclic groups of
its nonzero Smith factors. This proves all the index formulas.

Finally IF-02 gives
$A_n=\mathbb Z[L^{(a)},D^{(c)}:0\le a,c\le n]$.
After inverting $n!$, each generator is already in $O_n$,
so $A_n[1/n!]=O_n[1/n!]=(O_n^\pi)[1/n!]$.
The finite quotients can therefore have no prime divisor
larger than $n$. For $n=0$, all orders are $\mathbb Z$ and
the same formulas give index one.
\end{proof}

\paragraph{Integral row compression and independent shift blocks.}
For evaluation one may retain only the rows satisfying
\[
 t=0\quad\hbox{or}\quad v-k-l+t=0.                \tag{SN-15e}
\]
There are exactly $r_n$ of these rows. Removing the others preserves
all nonzero Smith factors, for the following integral reason.
Whenever $t>0$ and $v-k-l+t>0$, Pascal's identity in (SN-15a)
gives the row identity
\[
 R(v,k,l,t)=R(v,k,l,t-1)+R(v-2,k-1,l-1,t-1).      \tag{SN-15f}
\]
Both rows on the right are admissible. The intermediate powers of
$H$ agree, because replacing $(v,l)$ by $(v-2,l-1)$ leaves
$v-2(l-c)$ unchanged. The same argument works in (SN-15b), since
$n-(v-2)+2(l-1)=n-v+2l$. Recursive application strictly decreases
both positive quantities $t$ and $v-k-l+t$ along each nonterminal
branch. Every deleted row is therefore an integer combination of
retained rows, so subtracting these combinations gives zero rows
by unimodular row operations.

The four-size count in the proof shows that the number of retained
rows for a fixed $v$ is
$\binom{v+3}{3}-\binom{v+1}{3}$, where the subtracted term counts
the tuples with both the intersection and complement nonempty.
Summing gives
$\binom{n+4}{4}-\binom{n+2}{4}=r_n$. Since the original matrix
has rank $r_n$, the retained matrix has full row rank. Its
saturated column lattice is therefore the whole $\mathbb Z^{r_n}$.

Each nonzero entry further satisfies $k-l=a-c$. Sorting rows and
columns by this shift gives independent blocks for
$-n\le k-l\le n$, with no interactions between them. One may
compute each block's Smith factors, take the direct sum of the
resulting finite abelian groups, and sort the prime-wise valuations
to recover the global invariant factors. This is an integral block
decomposition, not a rational decomposition into irreducible
Lefschetz strings that might change the lattice.

\paragraph{Reproducible arithmetic and scope.}
The program \path{work/root/finite_index_formula.py} builds these
bounded matrices directly from (SN-15a)--(SN-15b) and records all
nonzero Smith factors. For $n=1,2,3,4$ it gives exactly the
indices of IF-03, for both orders; this is a third construction,
independent of the right-word and left-word closure programs.
It uses the same exact integer linear-algebra library and is
not described as an independent formal proof kernel.

For larger matrices, \path{work/root/finite_index_padic.py} performs
unit-pivot elimination modulo $p^c$ for each prime $p\le n$.
When no unit remains it divides the remaining matrix and modulus by
$p$; precision is doubled if exhausted, and no undecided exponent is
accepted. A unit pivot splits off one free summand over the local
ring; the division increases every remaining elementary-divisor
exponent by one. These two operations justify the computation.
The program \path{work/root/verify_finite_index_padic.py} checks the
new ranks instead by exact rational Schur complements, always
choosing a pivot of minimum $p$-adic valuation. All elimination
ratios are $p$-integral, so these are invertible row and column
operations over $\mathbb Z_p$. It uses neither modular precision
nor a Smith routine. The 144 block-and-prime computations at
$n=5,6$ agree, and the local method also reproduces every Smith
factor at $n=1,2,3,4$. The additional exact indices are
\[
\begin{array}{c|c|c|c}
n&r_n&[A_n:O_n]&[A_n:O_n^\pi]\\\hline
5&91&2^{201}3^{78}5^4&2^{80}3^{28}5^2\\
6&140&2^{420}3^{149}5^{20}&2^{168}3^{62}5^{10}.
\end{array}                                      \tag{SN-15g}
\]
All invariant factors, including unit factors, are recorded in the
JSON certificates. The prime-support proof above excludes any
unexamined prime; these are exact indices, not just lower bounds
from selected primes.

\paragraph{Geometric formulation.}
All matrices and formulas are defined over
$\operatorname{Spec}\mathbb Z$, before any rational or
$\ell$-adic realization. They concern the exterior model for
the Lefschetz algebra in cote 16 scans 014--018, with ordinary
contraction explicitly distinguished from normalized lowering.
Exterior powers give the $\lambda$-operations of SGA 6 V 2.2.1;
the Smith factors retain the integral information absent from
their rational $K_0$ character. Over $\mathbb Z_\ell$ the
quotients have lengths $v_\ell(\delta_{r_n})$. In the abelian
cohomological realization the rational Lefschetz operators are
algebraic by SN-10, but these arithmetic order indices alone
do not change a weight, determine a theta moment, or alter a
motivic local $L$-polynomial.


% Contribution owned by /root/singular_fibres.  Finding IDs SN-44-01--06.
% Uses only standard amsmath/amssymb/amsthm commands; bibliography entries
% are in bibliography.tex in this directory.

\section{Cote 44: the corrected height formula and the three nodal terms}
\label{sec:singular-fibres}

The text to be corrected is Grothendieck, cote 44, scan 004, Theorem 5
(manuscript page 2).  HODGE\_SIGN, AR-01, already disproves its displayed
formula under its explicitly transcribed hypotheses.  We do not repeat
that counterexample.  We construct the missing corrections, including a
normalization of the height, and compute them individually.  The
qualification ``geometrically integral fibres'' does not remove these
terms: their carriers can be non-Cartier cycles inside integral fibres.

\subsection{Conventions and the cited height comparison}

Let $k$ be an algebraically closed field of arbitrary characteristic.
Let $f:M\to B$ be a projective morphism over $k$,
where $M$ is smooth projective of dimension $2m$ and $B$ is a smooth
projective connected curve.  Assume that $f$ is flat and its generic
fibre $Y/K$, $K=k(B)$, is smooth and geometrically connected.  All Chow
groups in this subsection have rational coefficients.  An extension
$\alpha\in\operatorname{CH}^{m}(M)_{\mathbb Q}$ of a class
$z\in\operatorname{CH}^{m}(Y)_{\mathbb Q}$ is \emph{admissible} if
\[
 \deg(\alpha v)=0
 \quad\hbox{for every codimension-$m$ cycle $v$ supported on a fibre.}
 \tag{44.1}\label{eq:44-admissible}
\]
For algebraically trivial $z$, such extensions exist.  The degree of the
intersection of two admissible extensions is independent of the
extensions and of the smooth model.  These are the cited results of
Kahn, \emph{Refined height pairing}, Definition 2.2, Proposition 2.8,
Definition 2.10 and (2-9), Propositions 3.6 and 3.8, and Theorem 5.6(a)
\cite{KahnHeight2024}.  We write this \emph{unsigned geometric height} as
$h_{\mathrm{geom}}(z,z')$.  Existence uses algebraic equivalence in the
present application; it does not assume that all homologically trivial
cycles have numerically trivial extensions.

For the next comparison only, assume $k=\mathbb C$.
Here is an independent identification of this height, so that the
correction formula below is not a definition masquerading as a
comparison.  Put $P=\operatorname{Pic}^{m}_{Y}$, the higher Picard variety
for incidence equivalence, and let
$\theta:\operatorname{CH}^{m}_{\mathrm{alg}}(Y)\to P(\overline{\mathbb C(B)})$
be its Abel--Jacobi map.  Choose a Poincare $m$-cycle on $P\times Y$
whose induced endomorphism of $P$ is multiplication by $k_Y^m>0$.
Its transpose induces an isogeny $\lambda_Y^m:P\to P^\vee$.
The rational homomorphism $\lambda_Y^m/k_Y^m$ is the normalized
incidence duality; the integer is necessary when a universal cycle
does not exist with multiplicity one.  With the geometric Poincare
biextension normalization of the bilinear Neron--Tate pairing,
\[
 h_{\mathrm{geom}}(z,z')
 =\frac{1}{k_Y^m}
    \big\langle\theta z,\lambda_Y^m\theta z'\big\rangle_{\mathcal P}.
 \tag{44.2}\label{eq:44-nt-comparison}
\]
This is Chen, \emph{Degeneration of the archimedean height pairing of
algebraically trivial cycles}, arXiv:2512.22788v1, Theorem 7.9
(also Theorem 7.19(a),(b)); see Lemma 3.16, Definition 3.17,
Lemma 3.19 and Remark 7.8 for the normalization and descent
\cite{ChenHeight2025}.  It is the geometric form of Kunnemann's
\emph{Higher Picard varieties and the height pairing}, Theorem 8.2.
The comparison in \eqref{eq:44-nt-comparison} is a cited theorem, not a
new result of this run.  It does \emph{not} use the Griffiths incidence
conjecture: that conjecture enters Chen's further comparison with the
Lear extension of the Hodge-theoretic biextension, Theorem 7.19(c).

We use the sign convention of cote 44 by putting
\[
 \langle z,z'\rangle_G=(-1)^m h_{\mathrm{geom}}(z,z').
 \tag{44.3}\label{eq:44-signed-height}
\]
The pairing $\langle\ ,\ \rangle_{\mathcal P}$ is the degree of the
extended Poincare biextension; there is no additional factor $1/2$.
If a convention for the polarization of a quadratic height includes
that factor, the whole identity, including its right side, is divided
by $2$.  The assertion here is an identity for the canonical incidence
duality.  Its identification with a positive polarization requires the
primitive signs noted by Grothendieck in cote 44, scan 008; no
positivity of every unmodified higher Picard autoduality is assumed.
For curves, $k_Y^1=1$ and $\lambda_Y^1$ is the negative of the usual
principal polarization (Chen, Example 3.20), which checks the sign.
In arbitrary characteristic (44.3) is the signed version of Kahn's
geometric height.  Formula (44.2), as cited here, is a comparison
over complex function fields.  SN-44-06 proves its extension to
arbitrary characteristic separately, by reduction to curves and
with the finite-extension normalization stated there.  In particular
the geometric statement below does not require descent of a higher
Picard variety over an imperfect function field.

\subsection{SN-44-01: the corrected Theorem 5}
\label{sn:44-01}

\paragraph{Statement.}
Let $X/k$ be smooth projective of dimension $2m$, with a pencil
whose smooth base locus has codimension two.  Let
$\pi:M\to X$ be its blowup and $f:M\to\mathbb P^1$ the resulting
morphism.  Assume the generic member $Y$ is smooth and geometrically
connected.  Let $Z,Z'\in\operatorname{CH}^{m}(X)_{\mathbb Q}$ have
geometrically algebraically trivial restrictions $z,z'$ to $Y$.
Here $k$ is algebraically closed of arbitrary characteristic and the
smooth model just specified is part of the hypotheses; no general
resolution-of-singularities assertion in positive characteristic is
being used.  If algebraic triviality is initially witnessed over a
finite extension of $K$, transfer gives the same admissibility
assertions with rational coefficients over $K$ (Kahn, Theorem 5.6(a)
and Proposition 3.9).

Let $V_s\subset N^m(M)_{\mathbb Q}$ be the image of cycles supported
on the fibre $F_s$, where $N^m$ denotes cycles modulo numerical
equivalence on $M$.  Let $R_s$ be the radical of the intersection
pairing on $V_s$.  For each singular fibre choose actual vertical
cycles $D_{s,1},\ldots,D_{s,r_s}$ whose numerical images give a basis
of $V_s/R_s$, and put
\[
 G_s=(\deg D_{s,i}D_{s,j})_{i,j},\qquad
 a_s(Z)=(\deg(\pi^*Z\,D_{s,i}))_i.
\]
An empty matrix contributes zero.  Then
\[
 \boxed{\quad
 \langle z,z'\rangle_G
 =(-1)^m\left(\deg_X(ZZ')
       -\sum_s a_s(Z)^{\mathsf t}G_s^{-1}a_s(Z')\right).
 \quad}                                                     \tag{44.4}
 \label{eq:44-corrected}
\]
The summands and their sum are independent of all the indicated
choices.  This is the corrected geometric Theorem 5 with specified
height and specified extensions.  Formula (44.2) over $\mathbb C$,
and its separately proved extension in SN-44-06 in arbitrary
characteristic, identify its height with the indicated Neron--Tate
pairing, using the full finite-extension degree normalization.
It allows integral singular fibres and does
not require that their vertical cycles be Cartier.

\begin{proof}
The space $N^m(M)_{\mathbb Q}$ is finite dimensional, with dimension
at most the $2m$-th $\ell$-adic Betti number
($\ell\ne\operatorname{char}k$).  To see the bound precisely, take
any finite numerically independent set of rational cycle classes.
The definition of numerical equivalence supplies complementary
cycles whose rational intersection matrix has full rank.  Its
nonzero minor remains nonzero over $\mathbb Q_\ell$, so the cycle
classes are linearly independent over $\mathbb Q_\ell$ in
$H^{2m}(M,\mathbb Q_\ell(m))$.  This bounds every such finite set.
Hence every $V_s$ is finite dimensional and its quotient by the
radical has a nondegenerate symmetric form.
The matrix $G_s$ is therefore invertible.

Two cycles supported on distinct fibres have intersection zero.
For a smooth fibre $i:F\hookrightarrow M$, every vertical cycle is
$i_*u$, and the self-intersection formula gives
\[
 (i_*u)(i_*v)=i_*\bigl(uv\,c_1(N_{F/M})\bigr)=0,
\]
since $N_{F/M}$ is trivial.  Thus all smooth-fibre classes lie in the
radical of the full vertical space $V=\sum_sV_s$.  Moreover
$\bigoplus_s V_s/R_s\to V/\operatorname{rad}V$ is an isomorphism:
surjectivity is immediate, and a relation among the summands pairs
separately with each $V_s$, forcing each term to be radical.  The same
argument identifies the radical of $V$ with the sum of the $R_s$ and
the smooth-fibre spaces.  Only finitely many singular fibres occur.

Set $\alpha=\pi^*Z$, $\beta=\pi^*Z'$.  By the cited existence theorem,
choose an admissible extension $\alpha_0$ of $z$.  Localization for Chow
groups (Fulton, Proposition 1.8) shows that
$\alpha-\alpha_0$ is represented by a finite sum of fibre-supported
cycles.  Consequently $\alpha r=0$ for every $r\in\operatorname{rad}V$:
the $\alpha_0$ term vanishes by admissibility and the difference term
vanishes by the definition of the radical.  This proves in particular
that $a_s(Z)$ does not depend on the choices modulo $R_s$.

For the column $c_s=G_s^{-1}a_s(Z)$ define actual Chow classes
\[
 w_Z=\sum_{s,i}(c_s)_iD_{s,i},\qquad
 \alpha^h=\alpha-w_Z.
\]
They satisfy $\alpha^h D_{s,j}=0$ for every $s,j$, and $\alpha^h r=0$
for every vertical radical class $r$.  Thus $\alpha^h$ is admissible
and still restricts to $z$.  Apply the identical construction to
$\beta$, obtaining $w_{Z'}$ and $\beta^h$.

Orthogonality to $V$ gives
\[
 \deg(\alpha\beta)=\deg(\alpha^h\beta^h)+\deg(w_Zw_{Z'}),\qquad
 \deg(w_Zw_{Z'})=\sum_s a_s(Z)^{\mathsf t}G_s^{-1}a_s(Z').
\]
The first term is $h_{\mathrm{geom}}(z,z')$ by the cited construction;
the projection formula gives $\deg(\alpha\beta)=\deg_X(ZZ')$.
Equation \eqref{eq:44-signed-height} proves
\eqref{eq:44-corrected}; over $\mathbb C$ also apply
\eqref{eq:44-nt-comparison}.  A change of basis conjugates $G_s$
and changes the two columns contragrediently, leaving the scalar
unchanged; replacing a representative by a radical class has already
been handled.  This also proves the asserted choice independence.
\end{proof}

\paragraph{Other codimensions.}
For the hyperplane pencil in the original theorem, in its notation
$i+j=\dim X$, if $i\ne j$, one of $2i,2j$
is strictly less than $\dim X$.  Weak Lefschetz makes restriction in
that degree injective.  Algebraic triviality on $Y$ therefore makes
one original cycle homologically trivial on $X$.  Its pullback to
$M$ is admissible and has zero intersection with every complementary
cycle.  The geometric height and the raw intersection are both zero.
This proves the vanishing recorded in the margin of scan 004, with
the precise comparison above.  The correction issue is concentrated
in the middle case.

\subsection{SN-44-02: the vertical intersection form of an ordinary node}
\label{sn:44-02}

\paragraph{Statement.}
Let $m\ge1$, and let $f:M\to B$ have smooth projective total space of dimension
$2m$ over an algebraically closed field of characteristic different
from $2$.  Suppose $F_s$ is smooth except for one ordinary double
point $p_s$.  Blow up $p_s$ in $M$.  The exceptional divisor is
$E_s\simeq\mathbb P^{2m-1}$, and its intersection with the strict
transform $\overline F_s$ is the smooth quadric
$Q_s\simeq Q^{2m-2}$.  Choose its two middle ruling classes $a,b$.
For $m=1$, $Q_s$ consists of two reduced points and $a,b$ are their
classes; the smooth strict transform $\overline F_s$ may be disconnected.

For an $m$-dimensional cycle $D$ supported on $F_s$, take its strict
transform $\overline D=\overline i_*d$ and write
\[
 j^*d=u a+v b\quad\hbox{in }\operatorname{CH}^{m-1}(Q_s)_{\mathbb Q},
 \qquad \nu_s(D)=u-v.                                      \tag{44.5}
\]
Here $j:Q_s\hookrightarrow\overline F_s$ and
$\overline i:\overline F_s\hookrightarrow\operatorname{Bl}_{p_s}M$.
Then for any two such cycles
\[
 \boxed{\quad\deg_M(DD')=(-1)^m\nu_s(D)\nu_s(D').\quad}       \tag{44.6}
 \label{eq:44-node-form}
\]
The functional $\nu_s$ is well defined on $V_s$ up to simultaneous
change of sign.  Thus $V_s/R_s$ is either zero or one dimensional.
If it is nonzero, a rational cycle $D_s$ can be normalized by
$\nu_s(D_s)=1$, and then $D_s^2=(-1)^m$.

\begin{proof}
Write $p:\widetilde M\to M$ for this blowup, $E$ for its exceptional
divisor, and $h=c_1(\mathcal O_E(1))$.  The Chow group of the quadric
in the indicated degree is freely generated by its two ruling
classes.  Their elementary intersection table is
\[
 \begin{gathered}
 a+b=h^{m-1}|_Q,\qquad
 \deg(a h^{m-1})=\deg(b h^{m-1})=1,\\
 a^2=b^2=\frac{1+(-1)^{m-1}}2,\qquad
 ab=\frac{1-(-1)^{m-1}}2.
 \end{gathered}                                             \tag{44.7}
\]
For $m\ge2$, this is the quadric calculation already recorded and
proved in HODGE\_SIGN, SG-02, applied with its $r=m-1$; it is used
here as a cited calculation.  For $m=1$ the two point classes are
orthogonal idempotents, which gives the same table directly.
In particular each ruling plane has degree one.

Put $\delta_D=u+v$.  The restriction of $\overline D$ to $E$ is
$\delta_Dh^m$, because each ruling plane has degree one in $E$.
The blowup formula for Chow groups says that the kernel of
$p_*:\operatorname{CH}^{m}(\widetilde M)\to\operatorname{CH}^{m}(M)$
is generated by $E^m$.  Therefore $p^*D-\overline D=cE^m$.
Its restriction to $E$ is zero: the restriction of $p^*D$ factors
through the positive-codimension Chow group of a point.  Since
$E|_E=-h$, we obtain
\[
 p^*D=\overline D+(-1)^{m+1}\delta_D E^m.                   \tag{44.8}
\]
These uses of the blowup and self-intersection formulas are precisely
Fulton, Proposition 6.7 and Corollary 6.3; their displayed
specialization here also verifies the sign directly.

The nodal fibre has multiplicity two at $p_s$, so
$p^*F_s=\overline F_s+2E$.  Because a fibre divisor restricts
trivially to a fibre,
\[
 N_{\overline F_s/\widetilde M}=\mathcal O_{\overline F_s}(-2Q).
\]
Consequently
\[
 \deg(\overline D\,\overline D')=-2\deg_Q((ua+vb)(u'a+v'b)).
 \tag{44.9}
\]
Furthermore $E^{2m}=-1$ and
$\deg(\overline D E^m)=(-1)^{m-1}\delta_D$.  On expanding
$\deg(p^*D\,p^*D')$ with formulas (44.8)--(44.9), the two cross terms and
the exceptional square contribute $\delta_D\delta_{D'}$ in total.
Thus, more explicitly,
\[
 \deg(DD')=(u+v)(u'+v')-2\deg_Q((ua+vb)(u'a+v'b))
          =(-1)^m(u-v)(u'-v').
\]
This proves the claimed formula without a local factor $1/2$.

Strict transforms can change under rational equivalence by cycles
supported on $Q$.  For $m\ge2$, their restrictions to $Q$ add
multiples of $h^{m-1}=a+b$, since
$N_{Q/\overline F_s}=\mathcal O_Q(-1)$ and
$\operatorname{CH}^{m-2}(Q)$ is generated by $h^{m-2}$.
For $m=1$ the possible correction group is
$\operatorname{CH}^{-1}(Q)=0$; equivalently, there are no
one-dimensional cycles supported on the two-point scheme $Q$.
Hence $u-v$ is unchanged.  Independently, the formula just proved implies
that a numerically zero vertical class has $\nu_s(D)^2=0$ and hence
$\nu_s(D)=0$; it descends to $V_s$.  The radical is exactly its
kernel whenever its image is nonzero, and otherwise the entire
form is zero.  Scaling any cycle with nonzero $\nu_s$ proves the
last assertion.  Interchanging the two rulings changes all $\nu_s$
by a common sign and leaves every formula unchanged.
\end{proof}

\subsection{SN-44-03: the two vanishing-cycle cases}
\label{sn:44-03}

Assume now that the pencil is a Lefschetz pencil over $\mathbb C$,
with fibres of dimension $2m-1\ge1$.  Let $E_{\mathrm{van}}$ denote
the span of its vanishing cycles in $H^{2m-1}(Y_t,\mathbb Q)$.
This notation avoids confusion with the exceptional divisors in the
previous proof.  Let $S_*$ be the singular fibres for which the
\emph{actual algebraic} vertical space has nonzero $\nu_s$.
Choose the normalized $D_s$ of SN-44-02 and set
\[
 c_s(Z)=\deg(\pi^*Z\,D_s).
\]
Then the corrected Theorem 5 takes the particularly simple form
\[
 \boxed{\quad
 \langle z,z'\rangle_G
   =(-1)^m\deg_X(ZZ')-\sum_{s\in S_*}c_s(Z)c_s(Z').
 \quad}                                                    \tag{44.10}
 \label{eq:44-node-correction}
\]
Its admissible extension is
$\pi^*Z-(-1)^m\sum_{s\in S_*}c_s(Z)D_s$.

If $E_{\mathrm{van}}\ne0$, every local vanishing cycle is nonzero,
$S_*=\varnothing$, and the uncorrected formula is valid with the
height convention (44.3).  If $E_{\mathrm{van}}=0$, the height in
(44.10) is zero, and
\[
 (-1)^m\deg_X(ZZ')=\sum_{s\in S_*}c_s(Z)c_s(Z').             \tag{44.11}
\]
No claim that a prescribed integral generator of a local cohomology
group is algebraic is hidden in the notation: $D_s$ is chosen only
from the actual image of cycles.  A nonzero integral value of
$\nu_s$ may require rational scaling.  In SN-44-04 its value is
exactly one at all three fibres, so no denominator occurs.

\begin{proof}
Equation (44.10) follows from (44.4) and $G_s=((-1)^m)$.
The local specialization sequence at a node is Deligne, \emph{Weil I},
(4.1), or in etale notation (4.3.3).  Its homology dual identifies
the image of the local defect map from middle cycles of $F_s$ with
the kernel of
$\mathbb Q\longrightarrow H_{2m-1}(Y_t,\mathbb Q)$,
$1\mapsto\delta_s$.  Here is the link with the cycle calculation
(44.5).  For $m\ge2$, the circle-bundle sequence for
$\mathcal O_Q(-1)$ identifies the relevant local quotient with
\[
 H^{2m-2}(Q,\mathbb Q)\big/
       hH^{2m-4}(Q,\mathbb Q)
   = (\mathbb Q a\oplus\mathbb Q b)/\mathbb Q(a+b).
\]
The same description follows from the etale Gysin sequence, with
the Tate twists understood.  For $m=1$ use the reduced quotient
$H^0(Q,\mathbb Q)/\mathbb Q\cdot1$, which is again the displayed
one-dimensional quotient.  Restriction to $Q$ and this quotient
are the local defect map, so $ua+vb$ maps to $u-v$, up to a
nonzero choice of orientation.  This proves the required
identification rather than assuming a local algebraic generator.
If
$\delta_s\ne0$, the indicated map is injective and therefore
$\nu_s$ is zero on all vertical cycles.  The conjugacy of the local
vanishing cycles (SGA 7 II, Expose XVIII, 6.6; Deligne, \emph{Weil I},
(5.4) and (5.8)) shows that this applies to every singular fibre
when $E_{\mathrm{van}}\ne0$.

If $E_{\mathrm{van}}=0$, the fixed-part theorem for the pencil gives
an isomorphism of rational Hodge structures and of variations
\[
 H^{2m-1}(X,\mathbb Q)|_U\xrightarrow{\sim}
 R^{2m-1}f_*\mathbb Q|_U,
\]
where $U$ is the smooth locus of the base.  Here injectivity is weak
Lefschetz, and surjectivity is the invariant-cycle statement:
the vanishing part is the orthogonal complement of the fixed part.
Equivalently one can use Deligne, \emph{Weil I}, (5.8)(b)(1), together
with the global invariant-cycle theorem.  This makes the variation
constant as a variation of Hodge structure, not just as a local
system.

The $H^1$ of the higher Picard variety $P=\operatorname{Pic}^m_Y$,
with its Tate twist, is a direct summand of this odd cohomology.
Here are the maps.  Let $g=\dim P$; the case $g=0$ is immediate.
A Poincare cycle $\Gamma\in\operatorname{CH}^{m}(P\times Y)_{\mathbb Q}$
and its transpose induce
\[
 u:H^{2g-1}(P)\longrightarrow H^{2m-1}(Y),\qquad
 v:H^{2m-1}(Y)\longrightarrow H^1(P),
\]
with the required Tate twists understood.  The divisor
correspondence giving $vu$ induces a nonzero multiple of
$\lambda_Y^m:P\to P^\vee$, since $\Gamma$ induces a nonzero
multiple of the identity on the higher Picard variety.  Hence
$vu$ is invertible.  Thus $u$ is split injective, and the
Lefschetz isomorphism on the abelian variety $P$ identifies its
source with a twist of $H^1(P)$.  The cycle may be chosen over
$\mathbb C(t)$ with rational coefficients: take the trace of a
Poincare cycle defined over a finite Galois extension, which only
changes its nonzero multiplicity.  The construction therefore
gives the required maps of variations over a nonempty open $U$.
The existence and isogeny statements here are Saito,
\emph{Abelian varieties attached to cycles of intermediate
dimension}, Section 4, or Chen, Lemmas 3.16 and 3.19.
Thus the polarized weight-one variation of $P$ is constant, and
$P/\mathbb C(t)$ is isogenous to a constant abelian variety.
Indeed a subvariation of a constant polarizable Hodge structure
is specified by a fixed rational subspace, whose Hodge filtration
is its intersection with the fixed filtration; polarized weight-one
Hodge structures recover abelian varieties up to isogeny.

Every $\mathbb C(t)$-point of a constant abelian variety is constant:
the associated rational map from $\mathbb P^1$ extends by properness,
and its differential is zero because invariant differentials on
an abelian variety span its cotangent space whereas
$H^0(\mathbb P^1,\Omega^1)=0$.  The same assertion modulo a finite
group holds after isogeny.  Consequently the Poincare Neron--Tate
pairing of any two $\mathbb C(t)$-points of $P$ is zero.  Formula
(44.2) gives $\langle z,z'\rangle_G=0$, proving (44.11).
\end{proof}

\paragraph{What ``Cartier'' does and does not mean here.}
For $m=2$, vertical middle cycles are Weil divisors on threefold
fibres.  The zero-defect classes may be replaced, numerically for
this calculation, by the radical; this suffices.  The stronger
claim that \emph{every} such Weil divisor is Cartier is unnecessary
and is not asserted.  For $m>2$, middle vertical cycles are not
divisors at all.  The precise invariant needed by Theorem 5 is
the vertical intersection form, not a factoriality assertion.

\subsection{SN-44-04: three nodal fibres, three individual units}
\label{sn:44-04}

Use the pencil and notation of HODGE\_SIGN AR-01:
\[
 X=\mathbb P^2_x\times\mathbb P^2_y,\qquad
 Z=a^2-ab+b^2,\qquad
 M=\left\{\sum_{j=0}^2 d_j(\lambda,\mu)x_jy_j=0\right\}
       \subset X\times\mathbb P^1,
\]
with $(d_0,d_1,d_2)=(\lambda,\lambda+\mu,\lambda+2\mu)$.
The incidence model $M$ is the blowup of the smooth base locus.
Write $t=f^*[\mathrm{point}]$.  Its three singular parameters are
$s_0=[0:1]$, $s_1=[-1:1]$, and $s_2=[-2:1]$.  For each parameter
define the \emph{actual} surface
\[
 D_i=\{x=e_i,\ s=s_i\}\times\mathbb P^2_y\ \subset M.
 \tag{44.12}
\]
Then
\begin{gather}
 D_iD_j=0\ (i\ne j),\qquad \deg D_i^2=1,\qquad
 \deg(\pi^*Z\,D_i)=1,                                    \tag{44.13}\\
 \sum_{i=0}^2[D_i]=\pi^*Z+t(2a-b)
          \quad\hbox{in }\operatorname{CH}^2(M),           \tag{44.14}\\
 (\pi^*Z)^h=\pi^*Z-\sum_iD_i=t(b-2a).                    \tag{44.15}
\end{gather}
The last class is vertical and orthogonal to every vertical class.
Each nodal fibre therefore contributes precisely $1$ to the
subtraction in (44.10), and the corrected formula is
\[
 \boxed{\qquad 0=3-(1+1+1).\qquad}                       \tag{44.16}
\]

\begin{proof}
Let $Q_y$ be the rank-two universal quotient bundle on
$\mathbb P^2_y$, so that
\[
 0\longrightarrow\mathcal O(-b)\longrightarrow
 \mathcal O^{\oplus3}\longrightarrow Q_y\longrightarrow0.
\]
On $M$, the linear functional
$(d_0x_0,d_1x_1,d_2x_2)$ vanishes on the tautological line represented
by $y$, by the defining equation of $M$.  It therefore defines a
section of the rank-two bundle
\[
 \mathcal F=Q_y^\vee\otimes\mathcal O_M(a+t).
\]
Its zero scheme is exactly $D_0\sqcup D_1\sqcup D_2$.
To check the scheme structure, around $(e_i,s_i)$ on
$\mathbb P^2_x\times\mathbb P^1$, the three functions $d_jx_j$
are regular parameters: the two $d_j(s_i)$ with $j\ne i$ and
the derivative of $d_i$ are nonzero.  In $M$, their relation with
$y$ eliminates one of them on each standard $y$-chart.  The
remaining two are regular parameters normal to $D_i$.
Thus the section is regular and its zero scheme is reduced.

The top Chern class formula for a regular section gives
\[
 \begin{split}
 \sum_i[D_i]=c_2(\mathcal F)
 &=c_2(Q_y^\vee)+(a+t)c_1(Q_y^\vee)+(a+t)^2\\
 &=b^2-(a+t)b+(a+t)^2\\
 &=a^2-ab+b^2+t(2a-b),
 \end{split}
\]
using $t^2=0$.  This proves the integral Chow identity (44.14).
Since $a$ and $t$ restrict trivially to $D_i$, its normal bundle is
\[
 N_{D_i/M}=\mathcal F|_{D_i}=Q_y^\vee
             \simeq\Omega^1_{\mathbb P^2}(1).
\]
The Euler sequence gives $c(Q_y^\vee)=1-b+b^2$, and hence
$\deg D_i^2=\int_{\mathbb P^2}b^2=1$.  Distinct $D_i$ lie on
distinct fibres and are disjoint.  Finally
$\pi^*Z|_{D_i}=b^2$, proving all of (44.13).

For any vertical cycle $v$, projection to a point of the base gives
$t|_{\operatorname{supp}v}=0$ in the operational Chow ring.  Thus
$t(b-2a)$ has intersection zero with every vertical class; it is an
admissible extension of the zero generic class.  Its square is zero
because $t^2=0$.  Consequently the vertical projection of $\pi^*Z$
has components exactly $D_i$, and their three squares are the three
individual correction terms.  The already proved AR-01 facts
$Z|_Y=0$ and $\deg Z^2=3$ now give (44.16).
\end{proof}

This proof also identifies the local normalization in SN-44-02:
the strict transform of $D_i$ meets the exceptional quadric in one
ruling line, so $\nu_{s_i}(D_i)=1$ after choosing that ruling first.
The normal bundle computation independently verifies its norm.

\subsection{SN-44-05: the correction as a Tate-motive projector}
\label{sn:44-05}

\paragraph{Statement and proof.}
For a normalized nodal class $D_s$ in SN-44-02, put
\[
 p_s=(-1)^mD_s\times D_s\in\operatorname{CH}^{2m}(M\times M)_{\mathbb Q}.
\]
Composition of product correspondences multiplies their outside
factors by the middle intersection number.  Since
$D_s^2=(-1)^m$ and different fibres have zero intersection,
$p_s^2=p_s$ and $p_sp_r=0$ for $s\ne r$.
Each projector factors through, and has image isomorphic to,
the Tate motive $\mathbb L^m$: the two morphisms are supplied by
$D_s$ and $(-1)^mD_s$, whose composite on $\mathbb L^m$ is $1$.
Its action on a middle cycle class is
$p_s\alpha=(-1)^m\deg(\alpha D_s)D_s$.  Thus the corrected
extension in (44.10) is precisely the image of $\alpha$ under
$1-\sum_sp_s$, modulo the vertical radical.  This proves the
motivic assertion by explicit algebraic correspondences; it does
not assume a general standard conjecture. \qed

\paragraph{Schemes over $\mathbb Z$, $\lambda$-rings and $K$-theory.}
The example, all three $D_i$, the regular section of $\mathcal F$,
and its integral Chow identity are defined over
$\operatorname{Spec}\mathbb Z[1/2]$.  Over this base the three
parameters remain distinct, the coefficient differences needed
above are units, and the same local parameter check proves
relative smoothness of the total space and regularity of the
three zero components.  In $K_0(M)$ the Koszul resolution reads
\[
 [\mathcal O_{\coprod D_i}]
     =\lambda_{-1}(\mathcal F^\vee)
     =1-[\mathcal F^\vee]+[\det\mathcal F^\vee].
\]
Its degree-two Chern character is
$\operatorname{ch}_2[\mathcal O_{\coprod D_i}]
 =c_2(\mathcal F)=\pi^*Z+t(2a-b)$.
These identities are the $K$-theoretic and $\lambda$-ring versions
of (44.14).  Regular immersions and the normal-bundle sequence
belong to the framework of EGA IV$_4$, 19.1.1;
the intersection operations used here are the refined Gysin and
self-intersection operations of Fulton, Sections 6.2--6.3.

\paragraph{Etale sheaves, weights, motives and $L$-functions.}
For a fibre of the example over $\mathbb F_q$ of odd characteristic,
the classes $D_i$ give three algebraic Tate summands
$\mathbb Q_\ell(-2)\subset H^4(M_{\overline{\mathbb F}_q},\mathbb Q_\ell)$.
Their Frobenius eigenvalues are $q^2$, so their weights are $4$.
At sheaf level these are the three skyscraper defects in
$R^4f_*\mathbb Q_\ell$ allowed by Deligne, \emph{Weil I},
(5.8)(b)(2), in the case $E_{\mathrm{van}}=0$; the local
specialization sequence is (4.3.3), and the SGA locus is SGA 7 II,
Expose XIV, 3.2, and Expose XVIII, 6.6.  Their degree-four zeta
factor is $(1-q^2T)^{-3}$, and the corresponding motive is
$\mathbb L^2\oplus\mathbb L^2\oplus\mathbb L^2$.
Purity assigns these classes the expected weight but cannot delete
their intersection contribution.  This is the precise reason a
sign or weight calculation cannot repair the omitted term in
cote 44 by itself.  These are geometric height assertions over a
curve and fibrewise statements over $\mathbb Z[1/2]$; no assertion
about an Arakelov height over $\operatorname{Spec}\mathbb Z$ is made.

\subsection{SN-44-06: extension of the height comparison and the
\texorpdfstring{$E=0$}{E=0} conclusion to positive characteristic}
\label{sn:44-06}

This subsection supplies a separate algebraic proof of the extension;
it does not extend Chen's complex theorem by a change of notation.
The local calculation SN-44-02 and the explicit example already work
in characteristic different from $2$.  We now show that the
Neron--Tate interpretation and the zero-height assertion of SN-44-03
also hold in that range.  The height comparison itself works in every
characteristic.

\paragraph{Finite-extension normalization and descent.}
Let $k$ be algebraically closed, $B/k$ a smooth projective connected
curve, and $K=k(B)$.  An incidence Picard variety and its Poincare
cycle, incidence duality, and the finitely many points being paired
may first be defined over a finite extension $L/K$.  Compute their
usual geometric Poincare Neron--Tate pairing over the normalization
$B_L$ of $B$, and divide its degree by $[L:K]$.  This gives the
pairing relative to the original base $B$.  All finite extensions,
including inseparable ones, are normalized by their full degree.
The Poincare pairing and Kahn's geometric height have the same
base-change rule.  For the latter this is Proposition 3.9(iv),
extended to arbitrary regular generic fibres by Theorem 4.14(b)
of \cite{KahnHeight2024}; thus no resolution of the base-changed
high-dimensional model is needed.  The normalized number is
independent of the finite field of definition: compare over a
common finite extension.  It is an ordinary Neron--Tate pairing
on an abelian variety with this usual extension-of-fields
normalization.

For the $E=0$ assertion we need to preserve the rational base curve,
so descent only to an arbitrary finite separable extension would
not suffice.  Here is the more precise descent that does suffice.
Over any perfect field $F$, the incidence Picard variety of
$Y_{\overline F}$ descends to $F$.  Indeed the universal property
for incidence homomorphisms supplies, for each
$\sigma\in\operatorname{Gal}(\overline F/F)$, a unique compatible
isomorphism from its $\sigma$-twist, and uniqueness supplies the
cocycle identity.  This datum is continuous: the variety and a
correspondence realizing its incidence-universal map are defined
over a finite extension, and its stabilizer acts by their specified
base change.  Galois descent for abelian varieties is effective.
The same argument makes the incidence map equivariant, so a cycle
defined over $F$ has an $F$-rational image.  This is the
perfect-field descent recorded in Chen, Remark 3.21; Saito's
incidence construction and duality, Theorem 3.6 and Proposition 4.4,
apply over algebraically closed fields of arbitrary characteristic
\cite{SaitoPicard1979}.

In characteristic $p>0$, apply this to
$F=K^{\mathrm{perf}}$.  Trace a Poincare cycle from a finite Galois
extension to obtain a rational Poincare cycle over $F$; its induced
endomorphism is a positive integral multiple of the identity.
The abelian variety, this cycle, the duality, and the finitely many
points and identities all descend, by finite presentation, to some
finite purely inseparable subextension $K_r/K$ of $F/K$.
When $K=k(t)$ we can enlarge it to
$K_r=k(t^{1/p^r})$; its smooth projective base is again
$\mathbb P^1_k$.  This argument avoids asserting an unproved
descent statement over the imperfect field $k(t)$ itself.

\paragraph{The arbitrary-characteristic height comparison.}
With the finite-extension normalization just specified, (44.2)
holds over $K=k(B)$ in every characteristic for algebraically
trivial cycles.  Here is the reduction to the classical curve
pairing in full.  After a finite extension, algebraic equivalence
expresses $z$ as a finite sum of terms $\Gamma_*\gamma$, with
$C$ a smooth projective geometrically connected curve,
$\gamma=[c_0]-[c_1]\in\operatorname{CH}^1(C)$,
$c_0,c_1\in C(K)$, and
$\Gamma\in\operatorname{CH}^m(C\times Y)_{\mathbb Q}$.
It is enough to prove the identity for one such term, by
bilinearity.  This is the curve reduction in Kahn's proof of
Theorem 5.6(a), also used in Chen's proof of Theorem 7.9.

Admissibility is preserved by these correspondences, by Kahn,
Proposition 3.6 (and Proposition 4.12(c) after alterations).
The projection formula gives the adjunction
\[
 h_Y(\Gamma_*\gamma,w)=h_C(\gamma,\Gamma^*w).               \tag{44.17}
\]
To make the degree shift explicit, take the closure of $\Gamma$
in the fibre product of regular models and admissible lifts of
$\gamma,w$.  Both sides of (44.17) are the pushforward to $B$ of
the product of this closure with the two pulled-back lifts,
formed by refined Gysin maps.  Commutativity of these operations
and the projection formula identify them.  This is exactly the
calculation in Kahn's proof of Lemma 3.3, with the complementary
codimensions shifted by the degree of $\Gamma$; it is not a use
of the degree-zero-only statement of Proposition 3.8 at a wrong
degree.  If a regular model is unavailable, pull back to admissible
alterations, apply this calculation, and divide by their degrees,
as in Theorem 4.14(b).

For $C$, the regular surface model exists in all characteristics.
At each place, a rational extension of a degree-zero divisor that
has degree zero on every fibre component computes Neron's local
pairing by intersection.  This is P\'epin,
\emph{Neron's pairing and relative algebraic equivalence},
Definition 2.1 and Theorem 2.3 \cite{PepinNeron2011}.
The condition on component degrees is precisely the admissibility
condition for this surface.  Summing the local identities is the
global Neron pairing, with normalized valuations and degree on $B$.
Equivalently, it is the degree of the canonical rational extension
of the Poincare biextension.  Under the Abel--Jacobi map its
duality is $\lambda_C^1$, the negative of the theta polarization:
the transpose of the universal divisor pulls back to the given
degree-zero divisor on $C$.  Functoriality of Neron's pairing
therefore gives, in precisely the normalization of (44.2),
\[
 h_C(\gamma,\eta)
 =\langle\operatorname{AJ}\gamma,
                 \lambda_C^1\operatorname{AJ}\eta\rangle_{\mathcal P}.
                                                               \tag{44.18}
\]
This uses the established pairing on curves, not a
higher-codimension height conjecture or a conjecture on component
group duality.

Let $a=\operatorname{Pic}(\Gamma):J(C)\to P$ and
$b=\operatorname{Pic}(\Gamma^{\mathsf t}):P\to J(C)$.
Saito's functorial incidence-duality identity, Proposition 4.4,
specializes, since $k_C^1=1$, to
\[
 k_Y^m\lambda_C^1 b=a^\vee\lambda_Y^m.                    \tag{44.19}
\]
The incidence maps satisfy
$\theta(\Gamma_*\gamma)=a\operatorname{AJ}\gamma$ and
$\operatorname{AJ}(\Gamma^*w)=b\theta w$.
Combining (44.17)--(44.19) and functoriality of the Poincare
biextension proves
\[
\begin{split}
h_Y(z,w)
 &=\langle\operatorname{AJ}\gamma,
                 \lambda_C^1 b\theta w\rangle_{\mathcal P}\\
 &=\frac1{k_Y^m}\langle\operatorname{AJ}\gamma,
                 a^\vee\lambda_Y^m\theta w\rangle_{\mathcal P}\\
 &=\frac1{k_Y^m}\langle\theta z,
                 \lambda_Y^m\theta w\rangle_{\mathcal P}.
\end{split}
\]
Finally divide by the finite extension degree.  This proves
(44.2) in the stated generality.  If the traced Poincare cycle
has multiplicity $N k_Y^m$ instead of $k_Y^m$, its duality is
$N\lambda_Y^m$, so the normalized expression is unchanged.

The same argument covers complementary, possibly unequal,
codimensions $i+j=\dim Y+1$.  Choose a Poincare $i$-cycle on
$\operatorname{Pic}^i_Y\times Y$ of multiplicity $k_Y^i$; its
transpose gives
$\lambda_Y^j:\operatorname{Pic}^j_Y\to
(\operatorname{Pic}^i_Y)^\vee$.  Saito's same Proposition 4.4
then gives
\[
 h_Y(z,w)=\frac1{k_Y^i}
       \langle\theta_i z,\lambda_Y^j\theta_j w\rangle_{\mathcal P}.
\]
Indeed the curve correspondence has codimension $i$, its transpose
sends codimension $j$ to divisors, and (44.19) uses $k_Y^i$.
This supplies the Neron--Tate comparison for the nonmiddle
vanishing assertion in SN-44-01 as well.

\paragraph{The fixed part over $\mathbb P^1$ in arbitrary characteristic.}
Let the hypotheses of SN-44-03 hold over an algebraically closed
field of characteristic different from $2$, with $\ell$-adic
cohomology in place of Betti cohomology, and suppose
$E_{\mathrm{van}}=0$.  Deligne, \emph{Weil I}, (5.8)(b)(1),
says that $R^{2m-1}f_*\mathbb Q_\ell$ is constant.  Thus
$H^{2m-1}(Y_{\overline K},\mathbb Q_\ell)$ is a trivial
Galois representation.  Use the descent above to work over
$K_r=k(t^{1/p^r})$ in characteristic $p$ (and $K_r=K$ in
characteristic zero), so that the incidence Picard variety $P$
and a rational Poincare cycle are defined there.

The maps $u,v$ in the proof of SN-44-03 have the same construction
in $\ell$-adic cohomology.  The isogeny $vu$ makes $v$
surjective.  In particular $H^1(P_{\overline K},\mathbb Q_\ell)$,
up to its Tate twist, is a quotient of the trivial middle
representation, so is itself trivial.  All the Tate twists are
trivial on this geometric Galois group because $k$ is algebraically
closed.  The integral Tate module is also unramified: an
endomorphism of its free $\mathbb Z_\ell$ lattice that is zero
after tensoring with $\mathbb Q_\ell$ is zero.

The Neron--Ogg--Shafarevich criterion (SGA 7 I, Expose IX,
Corollary 2.2.9) now gives good reduction of $P$ at every point
of $\mathbb P^1_k$.  Its Neron model
$g:\mathcal A\to\mathbb P^1_k$ is therefore an abelian scheme.
It is projective: the global Neron model over a noetherian regular
curve is quasi-projective (SGA 7 I, IX, Proposition 3.2 and
3.2.1, printed pp.~347--348, citing Raynaud, VIII.2(3))
\cite{SGA7IHeight}, and good reduction makes it
proper.  Its total space is consequently smooth projective over
$k$.  We spell out the fixed-part argument to avoid confusing
constant up to isogeny with constant as a family.

The low-degree Leray exact sequence for $g$ is
\[
0\longrightarrow H^1(\mathbb P^1,\mathbb Q_\ell)
 \longrightarrow H^1(\mathcal A,\mathbb Q_\ell)
 \longrightarrow H^0(\mathbb P^1,R^1g_*\mathbb Q_\ell)
 \xrightarrow{d_2}H^2(\mathbb P^1,\mathbb Q_\ell)
 \longrightarrow H^2(\mathcal A,\mathbb Q_\ell).
\]
The first term is zero.  The last arrow is injective because the
zero section is a left inverse to $g^*$, so $d_2=0$.
Moreover $R^1g_*\mathbb Q_\ell$ is constant, as is already true
at the generic point.  Hence restriction gives
\[
 H^1(\mathcal A,\mathbb Q_\ell)
   \xrightarrow{\sim}H^1(\mathcal A_s,\mathbb Q_\ell)
 \quad\text{for every }s\in\mathbb P^1(k).                 \tag{44.20}
\]

Put $A_0=\operatorname{Pic}^0_{\mathcal A/k,\mathrm{red}}$,
an abelian variety over $k$.  A Poincare bundle exists after
rigidification at a $k$-point of $\mathcal A$, and restriction
to the fibres defines a homomorphism of abelian schemes
\[
 A_0\times_k\mathbb P^1\longrightarrow\mathcal A^\vee.
\]
By the Kummer identification
$V_\ell(\operatorname{Pic}^0(T)_{\mathrm{red}})
 =H^1(T,\mathbb Q_\ell(1))$ for smooth projective $T/k$,
(44.20) says that its maps on rational Tate modules are
isomorphisms.  In particular its generic fibre is an isogeny.
Dualizing, $P$ is isogenous over $K_r$ to a constant abelian
variety.  This conclusion permits nonconstant abelian schemes
that are constant only up to isogeny.

Every $K_r$-point of a constant abelian variety is constant:
the rational map from $\mathbb P^1_k$ extends by properness,
and every morphism $\mathbb P^1_k\to A_0$ is constant, by the
Albanese universal property and
$\operatorname{Alb}(\mathbb P^1)=0$.  This argument works also
for inseparable morphisms.  A quasi-inverse to the isogeny shows
that $P(K_r)$ modulo its $K_r/k$-trace is torsion.  The canonical
Poincare Neron--Tate pairing kills the trace and torsion, so it
vanishes on the points coming from $z,z'$.  The comparison proved
above and division by $[K_r:K]$ now give
$h_{\mathrm{geom}}(z,z')=0$ over the original base.

The $E\ne0$ part of SN-44-03 also carries over: use the etale
local specialization sequence (Deligne, (4.3.3)) and the
conjugacy in (5.8); the difference-of-rulings defect is the
algebraic realization of the same one-dimensional local term.
Thus (44.10) and (44.11), including their Neron--Tate
interpretation with the specified finite-extension normalization,
hold for Lefschetz pencils of odd-dimensional members at least
one over algebraically closed fields of characteristic
different from $2$.  The general Gram formula (44.4) remains
valid in characteristic $2$ whenever its stated smooth model
exists; the ordinary-node calculation used here is not asserted
in that characteristic. \qed


\part{Unestablished assertions and exact boundaries}
\section{What is, and is not, concluded}

This part records assertions not presented as proved findings. It does
not weaken the proved counterexamples or silently turn a restricted
intertwiner into a stronger one.

\subsection{The general index route}
The literal lower-critical ordinary-Gysin version is refuted by
SN-C02. The weighted replacement is completely proved in SN-C03--04.
For the dimension-consistent ordinary upper-Gysin variant, the exact
extra condition in the diagram is $A_0(W)$ on the hyperplane
$W\subset X\times X$, not on $Y\times Y$. No geometric
counterexample to the standard conjecture is claimed from the
formal module in SN-C05.

There is a useful additional cited sufficient condition:
Kleiman 1968, Proposition 3.6, gives
$D(W)+\text{hard Lefschetz}\Rightarrow A(W,L_W)$.
Thus numerical equals homological equivalence on $W$ supplies
the missing $A_0(W)$; no index inequality is needed for this
direction. Under $I(W,L_W)$, Corollary 3.9 gives the converse
$A(W,L_W)\Rightarrow D(W)$. These are classical implications,
not new proofs or consequences of trace purity.

In the abelian cases the missing operators and positivity have
both been supplied by the cited chain of SN-10--11. Outside
those cases this paper does not assert the general standard
conjectures $B$ or $I$. SN-14 refutes the proposed formal
trace-only implication to the positive-adjoint argument. It
does not refute geometric semisimplicity or the standard conjectures.
Numerical motivic semisimplicity and positivity on the specified
homological realization remain different assertions.

\subsection{The remaining stronger theta comparison}
IF-01 and IF-04 give genuine integral intertwining with Lefschetz
operators for their specified Fock and row-root actions. IF-02--03 and SN-15
answer the order-index question on the actual exterior lattice, with
an explicit bounded matrix formula in every rank and evaluated indices
through rank six.
IF-08 constructs the finite coefficient-kernel transport, and
IF-09 fixes the actual two-step partitions, determinant insertions,
and an explicitly evaluated nonzero kernel test. IF-10 gives an
explicit CM-character-normalized geometric coefficient realization.
IF-11 applies it to an inhabited genuine two-step theta-jet family.
SF-Q6-01--03 identify three exact arithmetic admissibility obstructions
for the existing numerical and small-weight specializations.
IF-12 supplies a genuine CM theta kernel with a nonzero fixed-root
commutator, and proves existence of such a torsion translate for any
nonzero source theta vector. This refutes automatic root invariance
of that actual theta-translation family; it does not assert that a
nonzero term survives a prescribed global characteristic-kernel sum.
SF-Q6-04 gives the actual stabilizer prime bound and a positive
complete numerical expression under its explicitly stated ties.
SF-Q6-07 proves that the convenient constant-twist row model cannot
satisfy those ties. SF-Q6-05 removes the common-pin condition for
a single family and transfers a given nonzero bottom period to a
complex terminal tail. SF-Q6-06 gives the exact diamond Fourier
test and translated-period formula. Its refined-level nonzero
period theorem has a definite ambient group; the actual two-step
Hida ambient is indefinite of signature $(k,1)$, and its prescribed
level has no arbitrary refinement slot. These remaining hypotheses
are not removed by either complex surjectivity or the finite
coefficient construction.
SF-Q6-08 removes unused global rank and split-bound assumptions
from the fixed inert-prime valuation proof. Its actual theta
factorization, kernel-lattice, and terminal-pairing hypotheses
remain explicit; it is not an assertion that an unevaluated
candidate supplies those inputs.

The still stronger assertion would identify each actual global
automorphic block matrix, with its moments, characters and
stabilizers, with a single prescribed Lefschetz operator or with
one universal Lefschetz order. The explicit formula IF-08i
transports those matrices; it does not discard their entries.
An exact Smith index of a fixed global matrix is the determinantal
divisor of that transported matrix, so its finite jets and class
data still matter. The parameter $m=2$ alone does not select them.

The approaches actually tried are recorded in the work notes and
ledger: multiaffine Fock; full polynomial Weyl and dual lattices;
all-rank saturation; exact small-rank orders; actual auxiliary and
column weights; noncompact oscillator obstruction; actual jet and
global block formulas; the explicit exponential-kernel test; genuine
inhabited algebraic theta jets; CM-character realization; the three
arithmetic size obstructions; and
the full numerical-bound route with its gain, error and stabilizer
terms. A positive raw contribution is not substituted for a
positive full bound. An exponential test kernel is not substituted
for the globally automorphic theta series.

The formula-only and lift-count-only claims are contradicted by
the specified lattice and normalization tests. A universal
nonexistence theorem for every possible normalized geometric
realization of the global chain is not asserted.

\section{Framework map}

\begin{longtable}{@{}p{0.16\textwidth}p{0.77\textwidth}@{}}
\toprule
Finding IDs & Algebraic-geometric and Grothendieck-framework interpretation\\
\midrule
\endfirsthead
\toprule Finding IDs & Interpretation\\\midrule
\endhead
SN-C01--05 &
Cycle classes, restriction and refined Gysin, at cote 16
scans 034, 036, 040, 049--050.
The extra $L$ changes codimension in the associated graded of
the $\gamma$ filtration. The maps have the usual Tate twists
in SGA 5 VII 7.1--7.2; the $\mathbf P^1$ example exists over
$\mathbf Z$. The corrected condition supplies inverse morphisms
in homological motives, not positivity or a new $L$-factor identity.\\
SN-10--12 &
Algebraic Lefschetz correspondences, primitive cycle pairings,
and abelian homological motives. Under the stipulated CM Hodge
input the homological and numerical categories agree on all
abelian products. Spreading uses EGA IV$_3$ 8.8.2 and
SGA 4 XVI 2.2. The explicit CM fourfold separates algebraicity
from divisor generation.\\
SN-13--14 &
The positive rational trace form used on cote 16 scan 049,
versus what a Frobenius-module $K_0$ class can see.
SGA 6 V 2.2.1 supplies the exact-sequence exterior filtration.
Pure nonsplit point sheaves have the same local $L$-factor
as their semisimplifications. Algebraic splittings imply
factorization of motivic local polynomials; the converse fails.\\
IF-01--04, SN-15 &
Integral exterior powers and divided powers over
$\operatorname{Spec}\mathbf Z$; the coefficient coalgebra
of $\mathrm{SL}_{2,\mathbf Z}$ and its dual order.
These refine cote 16 scans 014--018 and the predecessor
SL-04, SL-08--09. Integral lattices are retained before
passing to representations or motives.\\
IF-05--12, SF-Q6-01--08 &
Actual polynomial coefficient functors and finite theta jets,
with determinant lines and local lattice tests. Their exterior
and tensor operations give $\lambda$-operations in $K_0$;
base change to $\mathbf Z_\ell$ sends Smith factors to local
lengths. Rational realization loses those lengths.
No algebraicity of an arbitrary theta moment matrix is
inferred from the formal tensor construction.\\
SN-44-01--06 &
Admissible extensions over a curve, refined Gysin, blowup
intersection theory, N\'eron models (SGA 7 I IX 2.2.9, 3.2),
and local vanishing cycles (SGA 7 II XIV 3.2, XVIII 6.6).
The three explicit corrections exist over $\mathbf Z[1/2]$,
are given by a Koszul $\lambda_{-1}$ class, and split off
three Tate motives. Their finite-field factor is
$(1-q^2T)^{-3}$; purity does not remove their intersections.\\
\bottomrule
\end{longtable}

The attached registers give individual question states, every actual
check event, and the implications for the predecessor's open rows.
No general integral Hodge conjecture, higher Arakelov index theorem,
or general number-field motivic GRH is a consequence claimed here.



\begin{thebibliography}{99}
\bibitem{GrothCoates1966}
A. Grothendieck, letter to J. Coates dated 6.1.1966 as written,
cote 16, scans 048--050; especially scans 049--050.
CSG transcription L11d, printed p.~3, Proposition.
\url{https://csg.igrothendieck.org/wp-content/uploads/2024/04/L11d.pdf}.

\bibitem{GrothCoates1967}
A. Grothendieck, letter to J. Coates, 4.1.1967,
cote 16 scans 044--047. CSG transcription L10d,
printed p.~2, Proposition (i), (ii bis), and flag corollary.
\url{https://csg.igrothendieck.org/wp-content/uploads/2024/04/L10d.pdf}.

\bibitem{Archive}
A. Grothendieck, Montpellier archive, cotes 16 and 44,
supplied transcriptions numbered by scan.
Cote 16: scans 014--018, 034--040, 049--050.
Cote 44: scan 004, Theorem 5, manuscript p.~2;
scans 008--009 for incidence duality and primitive signs.
\url{https://grothendieck.umontpellier.fr/16.pdf}.

\bibitem{Kleiman1968}
S. L. Kleiman, \emph{Algebraic cycles and the Weil conjectures},
in \emph{Dix expos\'es sur la cohomologie des sch\'emas},
North-Holland/Masson, 1968, pp.~359--386.
Exact locators: 1.4.7; 2.3--2.5; 2.7; 2.12;
Theorem 2.13(i),(ii$'$),(iv), pp.~372--373;
Theorem 2A11(2), p.~377;
Proposition 3.6 and Corollary 3.9, pp.~379--380;
Theorem 3.11 and Corollaries 3.12, 3.14, pp.~380--382;
Lemma 4.6 and Theorem 4.7, pp.~385--386.
\url{https://download.uni-mainz.de/mathematik/Algebraische%20Geometrie/Lehre/WS23.Padische.1966.kleiman.weil.pdf}.

\bibitem{Milne1999a}
J. S. Milne, \emph{Lefschetz classes on abelian varieties},
Duke Mathematical Journal \textbf{96} (1999), 639--675,
Theorem 5.9 (author PDF p.~26), read with the author's correction
to its proof.
\url{https://www.jmilne.org/math/articles/1999aP.pdf};
\url{https://www.jmilne.org/math/articles/1999a.html}.

\bibitem{Milne2002}
J. S. Milne, \emph{Polarizations and Grothendieck's standard
conjectures}, Annals of Mathematics (2) \textbf{155} (2002), 599--610.
Later author version (September 19, 2001):
Theorem 3.3(a),(b), pp.~8--9; Corollary 3.4;
Remarks 3.7--3.8, p.~9.
\url{https://www.jmilne.org/math/articles/2002aP.pdf}.
Original author version (March 26/August 14, 2001):
Theorem 4.6(a),(b), pp.~13--14.
\url{https://www.jmilne.org/math/articles/2002aS.pdf}.
Addendum to Theorem 2.1:
\url{https://www.jmilne.org/math/articles/2002a.html}.

\bibitem{CMHodge032}
OpenAI, \emph{The rational Hodge conjecture for CM abelian
varieties}, October 6, 2026, supplied file 032.
Theorem 1.1, p.~1 and concluding proof p.~47;
Corollary 8.4, p.~50. Used as stipulated cited input.

\bibitem{VanHoftenSempliner}
P. van Hoften and J. Sempliner,
\emph{The Hodge conjecture for CM abelian varieties},
supplied Hexagon 2610.00259v2, Theorem I and \S1.1.2, p.~1.

\bibitem{MilneCM}
J. S. Milne, \emph{Complex Multiplication}, author course notes,
Propositions 3.12--3.13, p.~30.
\url{https://www.jmilne.org/math/CourseNotes/CM.pdf}.

\bibitem{MilneAV}
J. S. Milne, \emph{Abelian Varieties}, course notes,
version 1.1, July 27, 1998, Proposition 13.1, printed p.~53;
for exterior cohomology, Theorem 11.1(b).
The exact version consulted is
\url{https://www.jmilne.org/math/CourseNotes/AV110.pdf}.

\bibitem{Jannsen1992}
U. Jannsen, \emph{Motives, numerical equivalence, and
semi-simplicity}, Inventiones Mathematicae \textbf{107}
(1992), 447--452, Theorem 1.
\url{https://epub.uni-regensburg.de/26642/1/jannsen10.pdf}.

\bibitem{Alpoge}
L. Alp\"oge, \emph{Remark}, supplied four-page note and Lean archive,
20 November 2026. Steps 2--5, 7--10, 13, 17 as numbered in the
supplied note; the definitions and theorem declarations actually
used are listed in the integral sections and manifests.
\url{https://alpo.ge/remark.pdf}.
The supplied \path{ALPOGE_LEAN_MAP_20261120.md} and
\path{ALPOGE_LEAN_SCAN.tsv} are navigation/audit inputs,
not new certificates produced by this run.

\bibitem{HodgeSign}
\emph{HODGE\_SIGN}, supplied
\path{RETURN_HODGE_SIGN_final_20261121T0057Z/HODGE_SIGN.tex}.
Used locators: SG-02; SG-04 with its scope remark;
SG-08--SG-11; SL-04, SL-06--SL-09;
AR-01; U-01--U-05 and Q6g.
GCT\_NSRH F03 and the supplied independent check are cited
as directed by the current prompt.

\bibitem{WeilI}
P. Deligne, \emph{La conjecture de Weil. I},
Publications Math\'ematiques de l'IH\'ES \textbf{43}
(1974), 273--307.
Theorem 1.6; (4.1), (4.3.3); (5.4), (5.8)(a),(b).
\url{https://www.numdam.org/item/PMIHES_1974__43__273_0.pdf}.

\bibitem{WeilII}
P. Deligne, \emph{La conjecture de Weil. II},
Publications Math\'ematiques de l'IH\'ES \textbf{52}
(1980), 137--252.
Theorem 3.3.1 over a point; Theorem 4.1.1.
\url{https://www.numdam.org/item/PMIHES_1980__52__137_0.pdf}.

\bibitem{SGA4}
M. Artin, A. Grothendieck and J.-L. Verdier (eds.),
\emph{Th\'eorie des topos et cohomologie \'etale des sch\'emas},
SGA 4, Lecture Notes in Mathematics 269, 270, 305.
Expos\'e XVI, Corollary 2.2, with locally constant coefficients;
Expos\'e XVIII for duality.

\bibitem{SGA5}
A. Grothendieck, edited by L. Illusie,
\emph{Cohomologie $\ell$-adique et fonctions $L$},
SGA 5, Lecture Notes in Mathematics 589, 1977.
Expos\'e III, \S3, 5.1--5.3;
Expos\'e VII, Theorem 7.1 and Corollary 7.2.

\bibitem{SGA6}
P. Berthelot, A. Grothendieck and L. Illusie,
\emph{Th\'eorie des intersections et th\'eor\`eme de
Riemann--Roch}, SGA 6, Lecture Notes in Mathematics 225, 1971.
Expos\'e V, Definition 2.1, Example 2.2(b), Lemma 2.2.1.

\bibitem{SGA7}
A. Grothendieck, P. Deligne and N. Katz (eds.),
\emph{Groupes de monodromie en g\'eom\'etrie alg\'ebrique},
SGA 7 I--II, Lecture Notes in Mathematics 288 and 340.
I, Expos\'e IX, Corollary 2.2.9 and Proposition 3.2,
3.2.1; II, Expos\'e XIV, 3.2,
and Expos\'e XVIII, 6.6.

\bibitem{EGA}
A. Grothendieck and J. Dieudonn\'e,
\emph{\'El\'ements de g\'eom\'etrie alg\'ebrique}.
EGA IV$_3$, 8.8.2; EGA IV$_4$, 19.1.1.

\bibitem{Fulton}
W. Fulton, \emph{Intersection Theory}, second edition,
Springer, 1998.
Proposition 1.8; Sections 6.2--6.3 and Corollary 6.3;
Proposition 6.7; regular-section top Chern classes in \S14.1.

% The singular-fibre bibliography is appended by build_paper.py.


% Append inside the root paper's thebibliography environment, adjusting keys
% if the root already provides these references.
\bibitem{KahnHeight2024}
B. Kahn, with an appendix by Q. Liu,
\emph{Refined height pairing}, Algebra \& Number Theory \textbf{18}
(2024), no. 6, 1039--1079,
doi:10.2140/ant.2024.18.1039.
Locators: Definition 2.2; Proposition 2.8; Definition 2.10 and (2-9);
Propositions 3.6, 3.8, 3.9 and 4.12; Theorem 4.14(b);
Theorem 5.6(a); for the classical sign,
Theorem 5.10(b).
\url{https://msp.org/ant/2024/18-6/ant-v18-n6-p01-p.pdf}.

\bibitem{ChenHeight2025}
Z. Chen, \emph{Degeneration of the archimedean height pairing of
algebraically trivial cycles}, arXiv:2512.22788v1,
submitted 28 December 2025; downloaded PDF dated 30 December 2025.
Locators: Lemma 3.16; Definition 3.17; Lemma 3.19; Example 3.20;
Remark 3.21 (perfect-field descent);
Theorems 5.10, 5.13, 7.5, 7.9, and 7.19(a),(b).
The Neron--Tate comparison in 7.9 is distinct from the further
Griffiths-conditional Lear comparison in 7.19(c).
\url{https://arxiv.org/pdf/2512.22788v1}.

\bibitem{KunnemannPicard1996}
K. Kunnemann, \emph{Higher Picard varieties and the height pairing},
American Journal of Mathematics \textbf{118} (1996), no. 4,
781--797, doi:10.1353/ajm.1996.0033.  Locators: Lemma 8.1 and
Theorem 8.2.  Cited as the number-field antecedent of Chen's
geometric Theorem 7.9; no extension of the number-field theorem
to function fields is silently assumed.

\bibitem{SaitoPicard1979}
H. Saito, \emph{Abelian varieties attached to cycles of intermediate
dimension}, Nagoya Mathematical Journal \textbf{75} (1979),
95--119, doi:10.1017/S0027763000024855.
Locators: Theorem 3.6; (4.1); Proposition 4.4; Corollary 4.5.
The ground field is explicitly algebraically closed of arbitrary
characteristic (p.~96).  The Poincare-cycle and incidence-duality
formulas are restated in Chen, Lemmas 3.16 and 3.19.

\bibitem{PepinNeron2011}
C. P\'epin, \emph{Neron's pairing and relative algebraic equivalence},
arXiv:1103.0570v1 (2011), Definition 2.1 and Theorem 2.3,
proof pp.~6--8; see also the curve intersection formula on p.~2.
The base is a complete discrete valuation ring with algebraically
closed residue field, without a characteristic restriction.
\url{https://arxiv.org/pdf/1103.0570}.

\bibitem{SGA7IHeight}
A. Grothendieck, M. Raynaud and D.~S. Rim,
\emph{Groupes de monodromie en geometrie algebrique, I},
SGA 7 I, Lecture Notes in Mathematics \textbf{288}, Springer, 1972.
Expose IX: Corollary 2.2.9, pp.~335--336 (good reduction);
Proposition 3.2 and 3.2.1, pp.~347--348 (global Neron model on a
regular curve and quasi-projectivity).  In the retrieved scan these
are PDF pages 340--341 and 352--353, respectively; the expose's
internal page numbers are 17--18 and 29--30.
Original scan retrieved from
\url{https://library.slmath.org/nonmsri/sga/sga/pdf/sga7-1.pdf}.

% Existing shared bibliography keys to reuse if present:
% Fulton: W. Fulton, Intersection Theory, 2nd ed., 1998,
% Proposition 1.8; Corollary 6.3; Proposition 6.7(e);
% Example 6.3.4(a) / Proposition 14.1(c).
% Deligne Weil I: Publ. Math. IHES 43 (1974), 273--307,
% (4.1), (4.3.3), (5.4), (5.8)(a),(b).
% SGA 7 II, LNM 340 (1973), Expose XIV 3.2; Expose XVIII 6.6.

\end{thebibliography}
\end{document}
